\documentclass[11pt,letterpaper]{article}
\usepackage[utf8]{inputenc}
\usepackage[T1]{fontenc}
\usepackage{textcomp}
\usepackage{tgheros}
\renewcommand{\familydefault}{\sfdefault}
\usepackage{mathpazo}
\usepackage{amsmath,amssymb,amsthm,mathtools}
\usepackage[letterpaper,left=0.876in,right=0.876in,top=0.82in,bottom=0.82in]{geometry}
\usepackage{graphicx}
\graphicspath{{/tmp/figs/}{/tmp/book_assets/}}
\usepackage{booktabs}
\usepackage{tabularx}
\usepackage{array}
\usepackage[table]{xcolor}
\usepackage{colortbl}
\usepackage{eso-pic}
\usepackage{listings}
\usepackage{fancyhdr}
\usepackage{microtype}
\usepackage{caption}
\usepackage{subcaption}
\usepackage{hyperref}

% Exact Color Palette from Reference Monograph (The Architecture of Intelligence)
\definecolor{primary}{RGB}{31, 42, 68}        % Dark Navy #1F2A44
\definecolor{tealaccent}{RGB}{46, 168, 160}   % Teal #2EA8A0
\definecolor{accent}{RGB}{224, 122, 95}       % Coral #E07A5F
\definecolor{goldaccent}{RGB}{212, 168, 67}   % Gold #D4A843
\definecolor{slate}{RGB}{68, 80, 107}         % Slate #44506B
\definecolor{mutedgray}{RGB}{125, 139, 161}   % Muted Footer Gray #7D8BA1
\definecolor{bodyblack}{RGB}{34, 34, 34}      % Body Text #222222
\definecolor{boxbg}{RGB}{242, 246, 252}       % Callout Box Fill
\definecolor{tableheadbg}{RGB}{226, 234, 246} % Crisp Navy-Tinted Header Fill
\definecolor{rowalt}{RGB}{245, 247, 250}      % Alternating Row Fill
\definecolor{codebg}{RGB}{20, 30, 48}         % Dark Navy Code Box #141E30
\definecolor{bordergray}{RGB}{148, 163, 184}  % Table Grid Border

\color{bodyblack}
\arrayrulecolor{bordergray}

\hypersetup{
    colorlinks=true,
    linkcolor=primary,
    citecolor=tealaccent,
    urlcolor=primary,
    pdftitle={The Architecture of the Cosmos: Model-Independent Constraints on the Hubble Constant and Dark Energy Equation of State from DESI DR2 BAO, Pantheon+ Supernovae, and JWST Calibration},
    pdfauthor={Md Mushfiqur Rahim (ORCID: 0009-0003-1021-685X, github.com/MD-Mushfiqur123)},
    pdfsubject={Comprehensive Research Monograph on Observational Cosmology, Hubble Tension, DESI DR2 BAO, Pantheon+ Supernovae, JWST Calibration, and Dynamical Dark Energy}
}

\captionsetup{
    font={small,it},
    labelfont={bf,color=primary},
    labelsep=endash,
    skip=6pt
}

% Dark Navy Code Blocks matching Pages 13, 17, 21 of Reference Monograph
\lstset{
    backgroundcolor=\color{codebg},
    basicstyle=\ttfamily\scriptsize\color{white},
    breaklines=true,
    frame=single,
    rulecolor=\color{primary},
    framesep=6pt,
    xleftmargin=6pt,
    xrightmargin=6pt,
    numbers=left,
    numberstyle=\tiny\color{mutedgray},
    keywordstyle=\color{tealaccent}\bfseries,
    commentstyle=\color{goldaccent}\itshape,
    stringstyle=\color{accent},
    showstringspaces=false,
    tabsize=4
}

% Style \section with Dark Navy 16pt + Full-Width Teal Horizontal Rule below it!
\makeatletter
\renewcommand{\section}{\@startsection{section}{1}{\z@}%
  {-16pt \@plus -4pt \@minus -2pt}%
  {8pt \@plus 2pt}%
  {\normalfont\Large\bfseries\color{primary}}}
\let\orig@section\section
\renewcommand{\section}{\@ifstar{\star@section}{\nostar@section}}
\newcommand{\nostar@section}[1]{%
  \orig@section{#1}%
  \vspace{-6pt}\noindent{\color{tealaccent}\rule{\textwidth}{1.5pt}}\vspace{4pt}\par\nobreak%
}
\newcommand{\star@section}[1]{%
  \orig@section*{#1}%
  \vspace{-6pt}\noindent{\color{tealaccent}\rule{\textwidth}{1.5pt}}\vspace{4pt}\par\nobreak%
}

% Style \subsection in Bold Teal (#2EA8A0) matching Reference Monograph!
\renewcommand{\subsection}{\@startsection{subsection}{2}{\z@}%
  {-12pt \@plus -3pt \@minus -1pt}%
  {5pt \@plus 1pt}%
  {\normalfont\large\bfseries\color{tealaccent}}}
\makeatother

% Float tuning
\renewcommand{\topfraction}{0.92}
\renewcommand{\bottomfraction}{0.88}
\renewcommand{\textfraction}{0.06}
\renewcommand{\floatpagefraction}{0.85}
\setcounter{topnumber}{3}
\setcounter{bottomnumber}{3}
\setcounter{totalnumber}{5}
\flushbottom

\setlength{\parskip}{6pt plus 2pt minus 1pt}
\setlength{\parindent}{0pt}
\renewcommand{\arraystretch}{1.20}

% Exact Footer and Header matching Reference Monograph
\pagestyle{fancy}
\fancyhf{}


\fancyfoot[C]{%
  \vspace{-4pt}{\color{bordergray}\rule{\textwidth}{0.4pt}}\\[3pt]%
  {\footnotesize\color{mutedgray}The Architecture of the Cosmos $\;\cdot\;$ Md Mushfiqur Rahim $\;\cdot\;$ Page \textbf{\color{bodyblack}\thepage} of \textbf{\color{bodyblack}\pageref{LastContentPage}}}%
}
\renewcommand{\headrulewidth}{0pt}
\renewcommand{\footrulewidth}{0pt}

\fancypagestyle{plain}{%
  \fancyhf{}%
  \fancyfoot[C]{%
    \vspace{-4pt}{\color{bordergray}\rule{\textwidth}{0.4pt}}\\[3pt]%
    {\footnotesize\color{mutedgray}The Architecture of the Cosmos $\;\cdot\;$ Md Mushfiqur Rahim $\;\cdot\;$ Page \textbf{\color{bodyblack}\thepage} of \textbf{\color{bodyblack}\pageref{LastContentPage}}}%
  }%
  \renewcommand{\headrulewidth}{0pt}%
  \renewcommand{\footrulewidth}{0pt}%
}

% Custom Takeaway Callout Box with Teal Left Border / Box Fill
\newcommand{\takeaway}[2]{%
  \vspace{6pt}%
  \noindent\colorbox{boxbg}{%
    \parbox{\dimexpr\textwidth-2\fboxsep\relax}{%
      \vspace{5pt}\textbf{\color{primary}#1:}\hspace{0.4em}#2\vspace{5pt}%
    }%
  }%
  \vspace{6pt}%
}

\begin{document}

% ==============================================================================
% PAGE 1: FULL-BLEED FRONT COVER (Exact match to Page 1 of Reference Monograph)
% ==============================================================================
\thispagestyle{empty}
\AddToShipoutPictureBG*{\includegraphics[width=\paperwidth,height=\paperheight]{/tmp/book_assets/cover_front.png}}
\mbox{}
\clearpage

% ==============================================================================
% PAGE 2: CLEAN WHITE INNER TITLE PAGE (Exact match to Page 2 of Reference Monograph)
% ==============================================================================
\setcounter{page}{1}
\thispagestyle{empty}
\vspace*{1.35in}
\begin{center}
    {\normalsize\bfseries\color{tealaccent} A\ \ C O M P R E H E N S I V E\ \ R E S E A R C H\ \ M O N O G R A P H\par}
    \vspace{16pt}
    {\Huge\bfseries\color{primary} The Architecture of the Cosmos\par}
    \vspace{12pt}
    {\large\itshape\color{slate} Model-Independent Constraints on the Hubble Constant and Dark Energy Equation of State\\from DESI DR2 BAO, Pantheon+ Supernovae, and JWST Distance-Ladder Calibrations\par}
    \vspace{28pt}
    {\color{goldaccent}\rule{\textwidth}{1.2pt}\par}
    \vspace{42pt}
    {\small\color{mutedgray} Author\par}
    \vspace{6pt}
    {\Large\bfseries\color{primary} Md Mushfiqur Rahim$^{1,2,3,*}$\par}
    \vspace{8pt}
    {\footnotesize\color{slate}
    $^1$AI \& Systems Researcher; SOO \& Lead AI/ML Engineer, Craftly (\texttt{@Craftly-Robot}), Bogura, Bangladesh\\
    $^2$Department of Computer Science, University of the People, Pasadena, California, USA\\
    $^3$Science Division (HSC \textquotesingle 28), Govt.\ Azizul Haque College, Bogura, Bangladesh\\[4pt]
    $^*$ORCID: \href{https://orcid.org/0009-0003-1021-685X}{\texttt{0009-0003-1021-685X}} $\;\cdot\;$ \href{https://github.com/MD-Mushfiqur123}{\texttt{github.com/MD-Mushfiqur123}} $\;\cdot\;$ \href{https://g.dev/Mushfiqur}{\texttt{g.dev/Mushfiqur}} $\;\cdot\;$ \href{https://leetcode.com/u/mushfiqurH2B/}{\texttt{leetcode.com/u/mushfiqurH2B}}\par}
    \vfill
    {\normalsize\color{mutedgray} October 2026\par}
    \vspace*{0.6in}
\end{center}
\clearpage

% ==============================================================================
% PAGE 3: TABLE OF CONTENTS
% ==============================================================================
\tableofcontents
\clearpage

% ==============================================================================
% PAGE 4: EXECUTIVE SUMMARY (English + Bilingual Synthesis)
% ==============================================================================
\section*{Executive Summary}
\addcontentsline{toc}{section}{Executive Summary}

This monograph is organized as a rigorous progression from relativistic spacetime geometry to multi-probe Bayesian likelihood construction, Markov Chain Monte Carlo (MCMC) and Gaussian Process (GP) numerical execution, and theoretical physical interpretation. Readers seeking a single orientation before the full text should note the following five foundational takeaways, each expanded upon in the corresponding Part:

\begin{itemize}
    \setlength{\itemsep}{6pt}
    \item[\textcolor{tealaccent}{\textbullet}] \textbf{The Hubble tension ($4.52\sigma\text{--}5.72\sigma$) is a genuine physical anomaly, not an HST crowding artifact.} Within flat $\Lambda$CDM, DESI DR2 BAO + Planck CMB yields $H_0 = 68.17 \pm 0.28\text{ km s}^{-1}\text{ Mpc}^{-1}$ ($\Omega_m = 0.3027 \pm 0.0036$), in $4.52\sigma$ tension with SH0ES HST+Pantheon+ Cepheids ($73.04 \pm 1.04\text{ km s}^{-1}\text{ Mpc}^{-1}$) and $5.53\sigma$ tension with the 4-anchor Cepheid ladder ($73.17 \pm 0.86\text{ km s}^{-1}\text{ Mpc}^{-1}$). Direct galaxy-by-galaxy JWST vs.\ HST cross-checks (Part IV, Table~\ref{tab:host_galaxy_distances} and Figure~\ref{fig:jwst_hst_hosts}) confirm $\Delta\mu_{\text{JWST}-\text{HST}} = -0.02 \pm 0.03\text{ mag}$, proving that the lower CCHP JWST TRGB/JAGB values ($68.81 \pm 2.22$ and $67.80 \pm 2.72\text{ km s}^{-1}\text{ Mpc}^{-1}$) arise primarily from small-sample ($N=11$ vs.\ $N=42$) SN Ia host selection variance.
    \item[\textcolor{tealaccent}{\textbullet}] \textbf{Late-time dynamical dark energy does not resolve the Hubble tension---it decouples the two cosmic epochs.} Allowing $w(z)$ to evolve in CPL ($w_0 w_a$CDM) shifts the inverse-distance-ladder Hubble constant downward to $H_0 = 67.51 \pm 0.59\text{ km s}^{-1}\text{ Mpc}^{-1}$ (with Pantheon+), $66.74 \pm 0.56$ (with DES-Y5), and $65.91 \pm 0.84$ (with Union3). Resolving a local $H_0 \approx 72\text{--}73\text{ km s}^{-1}\text{ Mpc}^{-1}$ requires reducing the pre-recombination sound horizon $r_d$ by $\sim 5\text{--}7\%$ (to $r_d \approx 139.7\text{ Mpc}$ via Early Dark Energy at $z_c \approx 3630$, Part VIII), whereas $w(z)$ evolution operates at $z < 1$.
    \item[\textcolor{tealaccent}{\textbullet}] \textbf{Multi-probe data favor time-evolving dark energy ($w_0 > -1, w_a < 0$) at $2.8\sigma$ to $4.2\sigma$ with a phantom crossing at $z_c \approx 0.42$.} Combining DESI DR2 BAO with Planck CMB and Type Ia supernovae yields $w_0 = -0.838 \pm 0.055, w_a = -0.62^{+0.22}_{-0.19}$ ($2.8\sigma, \Delta\chi^2_{\rm MAP} = -10.7$) for Pantheon+; $w_0 = -0.667 \pm 0.088, w_a = -1.09^{+0.31}_{-0.27}$ ($3.8\sigma, \Delta\chi^2_{\rm MAP} = -17.4$) for Union3; and $w_0 = -0.752 \pm 0.057, w_a = -0.86^{+0.23}_{-0.20}$ ($4.2\sigma, \Delta\chi^2_{\rm MAP} = -21.0$) for DES-Y5 (Part V, Table~\ref{tab:cosmological_constraints}).
    \item[\textcolor{tealaccent}{\textbullet}] \textbf{Non-parametric Gaussian Processes and $y$-redshift cosmography independently confirm the phantom crossing and acceleration slowdown.} Without assuming any parametric form for $w(z)$, our Gaussian Process regression (Part VI, Figure~\ref{fig:gp_reconstruction}) recovers a phantom-to-quintessence crossing $w(z_c) = -1$ at $z_c = 0.42 \pm 0.09$ (where normalized dark energy density peaks at $f_{\rm DE} = 1.09 \pm 0.05$), an acceleration onset at $z_t = 0.73 \pm 0.06$ ($q(z_t) = 0$), and a present deceleration parameter $q_0 = -0.365 \pm 0.058$ with jerk $j_0 = 0.472 \pm 0.215$ (departing from the flat $\Lambda$CDM identity $j_0 = 1$).
    \item[\textcolor{tealaccent}{\textbullet}] \textbf{Bayesian evidence decomposition resolves the Frequentist vs.\ Bayesian paradox via Mirage Dark Energy.} Because the 2-parameter CPL prior incurs a Kullback-Leibler Occam penalty $\mathcal{D}_{\rm KL} \approx 4.5\text{--}5.2$, the nested-sampling Bayes factor is inconclusive for Pantheon+ ($\ln B = -0.66 \pm 0.44$), strong for Union3 ($\ln B = +3.03 \pm 0.44$), and near-decisive for DES-Y5 ($\ln B = +4.88 \pm 0.44$). Projecting onto the 1-parameter Mirage Dark Energy axis $w_a = -3.66(1+w_0)$ halves the Occam penalty and yields positive Bayesian support across all datasets ($\ln B = +1.85\text{ to }+6.65$, Part VII, Table~\ref{tab:model_comparison_full}).
\end{itemize}
\clearpage

% ==============================================================================
% PAGE 5: LIST OF FIGURES & LIST OF TABLES
% ==============================================================================
\listoffigures
\vspace{14pt}
\listoftables
\clearpage

% ==============================================================================
% PAGE 6: ABSTRACT & CORE RESEARCH VERDICT
% ==============================================================================
\section*{Abstract}
\addcontentsline{toc}{section}{Abstract}

{\small
This study presents a comprehensive, multi-probe Bayesian and non-parametric investigation of the two foremost empirical anomalies in contemporary physical cosmology: the local-versus-early Universe Hubble constant ($H_0$) tension and the late-time dynamical evolution of the dark energy equation of state $w(z)$. Combining the 13 correlated Baryon Acoustic Oscillation (BAO) distance measurements across $0.295 \le z \le 2.330$ from the Dark Energy Spectroscopic Instrument Data Release 2 (\href{https://arxiv.org/abs/2503.14738}{DESI Collaboration 2025a, arXiv:2503.14738}), the 1701 Type Ia supernova (SN Ia) light curves from the Pantheon+ compilation (\href{https://arxiv.org/abs/2202.04077}{Brout et al.\ 2022, arXiv:2202.04077}; \href{https://arxiv.org/abs/2112.03863}{Scolnic et al.\ 2022, arXiv:2112.03863}) alongside Union3 and DES-Y5 cross-checks, 32 model-independent Cosmic Chronometer (CC) $H(z)$ determinations (\href{https://arxiv.org/pdf/2608.01215}{Carloni et al.\ 2026, arXiv:2608.01215}), compressed Planck 2018 Cosmic Microwave Background (CMB) shift parameters (\href{https://arxiv.org/abs/1807.06209}{Planck Collaboration VI 2020, arXiv:1807.06209}), and the latest James Webb Space Telescope (JWST) local distance-ladder calibrations from the Chicago-Carnegie Hubble Program (\href{https://arxiv.org/abs/2408.06153}{Freedman et al.\ 2025, arXiv:2408.06153}; \href{https://arxiv.org/abs/2408.03474}{Lee et al.\ 2025, arXiv:2408.03474}) and the SH0ES Collaboration (\href{https://arxiv.org/abs/2408.11770}{Riess et al.\ 2024, arXiv:2408.11770}; \href{https://arxiv.org/pdf/2112.04510}{Riess et al.\ 2022, arXiv:2112.04510}), an end-to-end Affine-Invariant Ensemble Markov Chain Monte Carlo (MCMC) pipeline, nested-sampling Bayesian evidence decomposition, fourth-order cosmographic expansion, and non-parametric Gaussian Process (GP) regression are executed.

Two foundational results are established. First, regarding the Hubble tension: within flat $\Lambda$CDM, DESI DR2 BAO combined with Planck CMB yields $H_0 = 68.17 \pm 0.28\text{ km s}^{-1}\text{ Mpc}^{-1}$ ($\Omega_m = 0.3027 \pm 0.0036$), exhibiting a $4.52\sigma$ tension with the SH0ES HST+Pantheon+ Cepheid determination ($73.04 \pm 1.04\text{ km s}^{-1}\text{ Mpc}^{-1}$) and a $5.53\sigma$ tension with the full 4-anchor HST Cepheid ladder ($73.17 \pm 0.86\text{ km s}^{-1}\text{ Mpc}^{-1}$). Allowing late-time dynamical dark energy ($w_0 w_a$CDM) does \emph{not} alleviate the Hubble tension; instead, it shifts the inverse-distance-ladder Hubble constant downward to $H_0 = 67.51 \pm 0.59\text{ km s}^{-1}\text{ Mpc}^{-1}$ (with Pantheon+), $66.74 \pm 0.56\text{ km s}^{-1}\text{ Mpc}^{-1}$ (with DES-Y5), and $65.91 \pm 0.84\text{ km s}^{-1}\text{ Mpc}^{-1}$ (with Union3). In the local Universe, while CCHP JWST-only Tip of the Red Giant Branch (TRGB) and J-region Asymptotic Giant Branch (JAGB) subsamples yield $H_0 = 68.81 \pm 2.22$ and $67.80 \pm 2.72\text{ km s}^{-1}\text{ Mpc}^{-1}$ (reducing tension with Planck to $0.63\sigma$ and $0.16\sigma$), direct galaxy-by-galaxy HST-versus-JWST cross-checks confirm $\Delta\mu_{\text{JWST}-\text{HST}} = -0.02 \pm 0.03\text{ mag}$, showing that the $\Delta H_0 \sim 2.5\text{--}3.5\text{ km s}^{-1}\text{ Mpc}^{-1}$ offset between CCHP and SH0ES arises primarily from small-sample SN Ia host selection fluctuations ($N=11$ vs.\ $N=42$ calibrators) rather than HST Cepheid crowding bias. Combining all available JWST distance indicators across a distance-limited sample of 16 SNe Ia ($D \le 25\text{ Mpc}$) yields $H_0 = 72.60 \pm 2.00\text{ km s}^{-1}\text{ Mpc}^{-1}$, while our sound-horizon-free late-time fit (DESI DR2 + CC + uncalibrated Pantheon+ with $r_d$ free) yields $H_0 = 69.24 \pm 1.82\text{ km s}^{-1}\text{ Mpc}^{-1}$ ($r_d = 146.99 \pm 3.44\text{ Mpc}$).

Second, regarding dynamical dark energy: joint fits of DESI DR2 BAO + CMB + SNe Ia consistently occupy the quadrant $w_0 > -1$ and $w_a < 0$, indicating phantom behavior ($w(z) < -1$) in the past that crosses the phantom divide $w = -1$ at redshift $z_c \approx 0.42\text{--}0.45$ into a quintessence-like state today. Frequentist preferences over $\Lambda$CDM reach $3.1\sigma$ ($\Delta\chi^2_{\rm MAP} = -12.5$) for DESI DR2 + CMB alone, $2.8\sigma$ ($\Delta\chi^2_{\rm MAP} = -10.7$, $w_0 = -0.838 \pm 0.055$, $w_a = -0.62^{+0.22}_{-0.19}$) with Pantheon+, $3.8\sigma$ ($\Delta\chi^2_{\rm MAP} = -17.4$, $w_0 = -0.667 \pm 0.088$, $w_a = -1.09^{+0.31}_{-0.27}$) with Union3, and $4.2\sigma$ ($\Delta\chi^2_{\rm MAP} = -21.0$, $w_0 = -0.752 \pm 0.057$, $w_a = -0.86^{+0.23}_{-0.20}$) with DES-Y5. Non-parametric GP reconstructions independently recover this phantom crossing at $z_c = 0.42 \pm 0.09$ and a cosmic deceleration-to-acceleration transition at $z_t = 0.73 \pm 0.06$, while nested-sampling Bayes factors (\href{https://arxiv.org/abs/2503.14743}{Lodha et al.\ 2025, arXiv:2503.14743}) reveal how the Kullback-Leibler Occam penalty ($\mathcal{D}_{\rm KL} \approx 4.5\text{--}5.2$) renders $\ln B_{w_0w_a,\Lambda} = -0.66 \pm 0.44$ for Pantheon+, $+3.03 \pm 0.44$ (strong) for Union3, and $+4.88 \pm 0.44$ (near-decisive) for DES-Y5.
}
\vspace{2pt}
\noindent\colorbox{boxbg}{%
\parbox{\dimexpr\textwidth-2\fboxsep\relax}{%
\centering\footnotesize\color{primary}%
\textbf{Observational Scope:} 13 DESI DR2 BAO Observables ($0.295 \le z \le 2.330$) $\;\cdot\;$ 1701 Pantheon+ SNe Ia $\;\cdot\;$ 32 Cosmic Chronometers $\;\cdot\;$ 24 JWST/HST Hosts $\;\cdot\;$ 11 Vector Figures%
}%
}
\vspace{4pt}
\noindent{\color{goldaccent}\rule{\textwidth}{1.2pt}}

\takeaway{Core Research Verdict: The Hubble Tension and Dynamical Dark Energy}{%
\textbf{(1) Is the Hubble Tension Real?} Yes, at $4.52\sigma\text{--}5.72\sigma$ between early-Universe sound-horizon physics ($r_d = 147.09 \pm 0.26\text{ Mpc}$) and the full 42-SN Cepheid distance ladder ($H_0 = 73.04 \pm 1.04\text{ km s}^{-1}\text{ Mpc}^{-1}$), though CCHP JWST TRGB/JAGB subsamples ($H_0 \approx 67.8\text{--}70.4\text{ km s}^{-1}\text{ Mpc}^{-1}$) highlight the sensitivity of local estimates to small-sample SN host selection variance. Late-time dynamical dark energy ($w_0 w_a$CDM) cannot resolve the $H_0$ tension without modifying pre-recombination sound-horizon physics. \textbf{(2) Does $w(z)$ Evolve with Time?} DESI DR2 BAO + CMB + SNe Ia show a $2.8\sigma\text{--}4.2\sigma$ frequentist preference for time-evolving dark energy ($w_0 > -1, w_a < 0$) crossing $w=-1$ near $z_c \approx 0.42$, driven physically by a $\sim 2\sigma$ mismatch between the CMB-inferred $\Omega_m = 0.3153$ and the lower $\Omega_m \approx 0.2975$ favored by DESI DR2 BAO distance ratios at $0.5 \le z \le 1.5$.}

\section{Part I --- Introduction: The Dual Crisis of Modern Observational Cosmology}

For more than two decades following the discovery of late-time cosmic acceleration via Type Ia supernovae (SNe Ia), the six-parameter spatially flat $\Lambda$ Cold Dark Matter ($\Lambda$CDM) concordance model has served as the foundational paradigm of physical cosmology. Within this minimal framework, non-relativistic Cold Dark Matter drives hierarchical structure formation while a time-invariant cosmological constant $\Lambda$—characterized by a constant vacuum energy density $\rho_\Lambda = \Lambda c^2 / (8\pi G)$ and an invariant barotropic equation-of-state parameter $w \equiv P_{\rm DE}/(\rho_{\rm DE}c^2) = -1$—governs the accelerated expansion of the late Universe. This paradigm successfully accounts for the acoustic peak structure of the Cosmic Microwave Background (CMB) temperature and polarization anisotropies at photon decoupling ($z_* \simeq 1089.92$; \href{https://arxiv.org/abs/1807.06209}{Planck Collaboration VI 2020}), primordial Big Bang Nucleosynthesis (BBN) light-element abundances, and the linear growth of large-scale structure across four decades of spatial scale.

As observational cosmology has transitioned into the Stage-IV precision regime, however, two profound empirical challenges have emerged that test the internal consistency of $\Lambda$CDM across cosmic history:

\begin{enumerate}
    \setlength{\itemsep}{3pt}
    \item \textbf{The Hubble Constant ($H_0$) Tension:} Early-Universe inferences of the present expansion rate $H_0 \equiv H(z=0)$, anchored by the comoving sound horizon at the baryon drag epoch $r_d \equiv r_s(z_d) = 147.09 \pm 0.26\text{ Mpc}$ measured by the Planck satellite, yield $H_0 = 67.36 \pm 0.54\text{ km s}^{-1}\text{ Mpc}^{-1}$ (\href{https://arxiv.org/abs/1807.06209}{Planck Collaboration VI 2020}). When combined with the latest Baryon Acoustic Oscillation (BAO) measurements from the Dark Energy Spectroscopic Instrument Data Release 2 (DESI DR2), the early-Universe $\Lambda$CDM constraint tightens to $H_0 = 68.17 \pm 0.28\text{ km s}^{-1}\text{ Mpc}^{-1}$ (\href{https://arxiv.org/abs/2503.14738}{DESI Collaboration 2025a}). Conversely, direct late-Universe distance-ladder measurements from the SH0ES Collaboration—calibrating the absolute peak $B$-band magnitude $M_B^0 = -19.253 \pm 0.027\text{ mag}$ of Pantheon+ SNe Ia via Hubble Space Telescope (HST) and James Webb Space Telescope (JWST) observations of Cepheid variables across 37 host galaxies and 4 geometric anchors—yield $H_0 = 73.04 \pm 1.04\text{ km s}^{-1}\text{ Mpc}^{-1}$ (\href{https://arxiv.org/pdf/2112.04510}{Riess et al.\ 2022}) and $H_0 = 73.17 \pm 0.86\text{ km s}^{-1}\text{ Mpc}^{-1}$ in the updated 42-SN compilation (\href{https://arxiv.org/abs/2408.11770}{Riess et al.\ 2024}). Concurrently, competing JWST calibrations from the Chicago-Carnegie Hubble Program (CCHP) utilizing the Tip of the Red Giant Branch (TRGB) and carbon-rich J-region Asymptotic Giant Branch (JAGB) stars report lower local values of $H_0 = 68.81 \pm 1.79\text{ (stat)} \pm 1.32\text{ (sys)}\text{ km s}^{-1}\text{ Mpc}^{-1}$ (JWST TRGB), $H_0 = 67.80 \pm 2.17\text{ (stat)} \pm 1.64\text{ (sys)}\text{ km s}^{-1}\text{ Mpc}^{-1}$ (JWST JAGB; \href{https://arxiv.org/abs/2408.03474}{Lee et al.\ 2025}), and $H_0 = 70.39 \pm 1.22\text{ (stat)} \pm 1.33\text{ (sys)} \pm 0.70\text{ ($\sigma_{\rm SN}$)}\text{ km s}^{-1}\text{ Mpc}^{-1}$ for the combined 24-calibrator HST+JWST TRGB sample (\href{https://arxiv.org/abs/2408.06153}{Freedman et al.\ 2025}). Disentangling whether the Hubble tension reflects beyond-$\Lambda$CDM early-Universe physics or subtle astrophysical and sample-selection systematics in the second and third rungs of the local distance ladder requires a rigorous, model-independent statistical framework.
    \item \textbf{The Dynamical Dark Energy ($w(z)$) Anomaly from DESI DR2:} Using over 14 million extragalactic spectra of bright galaxies (BGS), luminous red galaxies (LRG), emission-line galaxies (ELG), quasars (QSO), and more than 1.28 million Lyman-$\alpha$ ($\text{Ly}\alpha$) forest sightlines across $14,200\text{ deg}^2$, DESI DR2 has achieved sub-percent precision on the cosmic expansion history over $0.295 \le z \le 2.330$ (\href{https://arxiv.org/abs/2503.14738}{DESI Collaboration 2025a}). When analyzed within the Chevallier-Polarski-Linder (CPL) parameterization $w(a) = w_0 + w_a(1-a)$, DESI DR2 BAO + CMB data alone disfavor $\Lambda$CDM ($w_0=-1, w_a=0$) at $3.1\sigma$ ($\Delta\chi^2_{\rm MAP} = -12.5$). Combining DESI DR2 + CMB with three major SN Ia compilations yields preferences for evolving dark energy of $2.8\sigma$ ($\Delta\chi^2_{\rm MAP} = -10.7$) with Pantheon+ (\href{https://arxiv.org/abs/2202.04077}{Brout et al.\ 2022}), $3.8\sigma$ ($\Delta\chi^2_{\rm MAP} = -17.4$) with Union3, and $4.2\sigma$ ($\Delta\chi^2_{\rm MAP} = -21.0$) with the Dark Energy Survey 5-Year (DES-Y5) sample (\href{https://arxiv.org/abs/2503.14738}{DESI Collaboration 2025a}; \href{https://arxiv.org/abs/2503.14743}{Lodha et al.\ 2025}). In all combinations, the posterior resides in the quadrant $w_0 > -1$ and $w_a < 0$ with $w_0 + w_a < -1$, implying a "phantom crossing" where dark energy violated the Null Energy Condition ($\rho_{\rm DE} c^2 + P_{\rm DE} < 0$) at $z > z_c \approx 0.45$ before transitioning into a quintessence-like state ($w > -1$) at late times.
\end{enumerate}

The remainder of this monograph is organized as follows. Section 2 establishes the relativistic Friedmann background equations, BAO and supernova distance indicators, and seven parametric and physical dark energy formalisms. Section 3 details the observational datasets, exact covariance matrices, Cosmic Chronometer compilation, and joint likelihood construction. Section 4 examines the JWST versus HST local distance-ladder debate and presents galaxy-by-galaxy comparisons across NGC 4258, M101, NGC 5643, and shared SN Ia hosts. Section 5 presents our Affine-Invariant Ensemble MCMC constraints across $\Lambda$CDM, $w$CDM, $w_0 w_a$CDM, pivoted $(w_p, f_p)$ coordinates, and sound-horizon-free ($r_d$-free) models. Section 6 details our non-parametric Gaussian Process reconstructions of $w(z)$, $f_{\rm DE}(z)$, $H(z)/(1+z)$, $q(z)$, and the $Om(z)$ diagnostic alongside fourth-order cosmographic expansions and Chebyshev polynomial Crossing Statistics. Section 7 presents the complete Hubble tension significance matrix and nested-sampling Bayes factor decompositions, and Section 8 synthesizes the physical implications for dark sector physics.


\begin{figure}[!htbp]
    \centering
    \includegraphics[width=0.98\textwidth]{/tmp/book_assets/fig_1_1_timeline.png}
    \caption{\textbf{Major Milestones in Cosmic Expansion History and Observational Probes ($z \approx 3630 \to z = 0$).} Sequence of physical epochs anchored by Early Dark Energy, Planck CMB recombination, DESI DR2 BAO tracers, cosmic acceleration onset ($z_t = 0.73$), phantom crossing ($z_c = 0.42$), and local JWST/HST distance-ladder calibrations.}
    \label{fig:cosmic_timeline_1_1}
\end{figure}

\section{Part II --- Theoretical Framework: Expansion Dynamics and Dark Energy Formalisms}

\subsection{Background FLRW Metric and Relativistic Friedmann Equations}

Assuming large-scale statistical homogeneity and isotropy (the Cosmological Principle), spacetime geometry is described by the Friedmann-Lema\^{i}tre-Robertson-Walker (FLRW) line element:
\begin{equation}
    ds^2 = -c^2 dt^2 + a^2(t) \left[ \frac{dr^2}{1 - K r^2} + r^2 \left( d\theta^2 + \sin^2\theta\, d\phi^2 \right) \right],
    \label{eq:flrw}
\end{equation}
where $a(t) = (1+z)^{-1}$ is the cosmic scale factor normalized to $a(t_0) = 1$ today, and $K$ is the spatial curvature constant related to the present curvature density parameter via $\Omega_K \equiv -K c^2 / (a_0^2 H_0^2)$. Substituting Eq.~(\ref{eq:flrw}) into Einstein's field equations $G_{\mu\nu} \equiv R_{\mu\nu} - \frac{1}{2}g_{\mu\nu}R = \frac{8\pi G}{c^4} T_{\mu\nu}$ for a multi-component perfect fluid $T^\mu_{\ \nu} = \text{diag}(-\rho c^2, P, P, P)$ yields the first and second Friedmann equations:
\begin{align}
    H^2(z) &\equiv \left(\frac{\dot{a}}{a}\right)^2 = \frac{8\pi G}{3} \sum_i \rho_i(z) + \frac{\Omega_K H_0^2}{a^2}, \label{eq:friedmann1} \\
    \frac{\ddot{a}}{a} &= -\frac{4\pi G}{3} \sum_i \left( \rho_i(z) + \frac{3 P_i(z)}{c^2} \right). \label{eq:friedmann2}
\end{align}
Covariant conservation of the stress-energy tensor, $\nabla_\mu T^{\mu\nu}_{\rm DE} = 0$, for a minimally coupled dark energy fluid with barotropic equation of state $w(z) \equiv P_{\rm DE}(z) / [\rho_{\rm DE}(z) c^2]$ yields the relativistic continuity equation:
\begin{equation}
    \dot{\rho}_{\rm DE} + 3H \left[ 1 + w(z) \right] \rho_{\rm DE} = 0 \quad \Longleftrightarrow \quad \frac{d \ln \rho_{\rm DE}}{d \ln a} = -3 \left[ 1 + w(a) \right] \quad \Longleftrightarrow \quad \frac{d \ln \rho_{\rm DE}}{dz} = \frac{3\left[1 + w(z)\right]}{1 + z}.
    \label{eq:continuity}
\end{equation}
Integrating Eq.~(\ref{eq:continuity}) from the present epoch ($a=1, z=0$) to redshift $z$ defines the normalized dark energy density evolution function $f_{\rm DE}(z)$:
\begin{equation}
    f_{\rm DE}(z) \equiv \frac{\rho_{\rm DE}(z)}{\rho_{\rm DE,0}} = \exp \left[ 3 \int_0^z \frac{1 + w(z')}{1 + z'}\, dz' \right] = \exp \left[ 3 \int_a^1 \frac{1 + w(a')}{a'}\, da' \right].
    \label{eq:fde_general}
\end{equation}
Notice an immediate physical consequence of Eq.~(\ref{eq:continuity}): the derivative of the dark energy density with respect to redshift is $\frac{d\rho_{\rm DE}}{dz} = \frac{3\rho_{\rm DE}(z)}{1+z}[1 + w(z)]$. Therefore, in a quintessence regime ($w(z) > -1$), $\rho_{\rm DE}(z)$ increases with redshift (diluting as the Universe expands), whereas in a phantom regime ($w(z) < -1$), $\rho_{\rm DE}(z)$ \emph{decreases} with redshift (growing as the Universe expands). Consequently, if dark energy crosses the phantom divide $w(z_c) = -1$ at some transition redshift $z_c$, its energy density $\rho_{\rm DE}(z)$ must attain a local extremum ($\left.\frac{d\rho_{\rm DE}}{dz}\right|_{z_c} = 0$) at $z = z_c$.

Defining the dimensionless expansion rate $E(z) \equiv H(z)/H_0$, the general expansion history including pressureless baryonic plus cold dark matter $\Omega_{\rm bc} = \Omega_b + \Omega_c$, photons $\Omega_\gamma$, massive neutrinos $\Omega_\nu$, spatial curvature $\Omega_K$, and dark energy $\Omega_{\rm DE} = 1 - \Omega_m - \Omega_r - \Omega_K$ reads (\href{https://arxiv.org/abs/2503.14743}{Lodha et al.\ 2025}):
\begin{equation}
    E^2(z) \equiv \frac{H^2(z)}{H_0^2} = \Omega_\gamma (1+z)^4 + \Omega_\nu \frac{\rho_\nu(z)}{\rho_{\nu,0}} + \Omega_{\rm bc}(1+z)^3 + \Omega_K(1+z)^2 + \Omega_{\rm DE}\, f_{\rm DE}(z).
    \label{eq:Ez_full}
\end{equation}
Here the photon density is fixed by the COBE/FIRAS CMB monopole temperature $T_{\rm CMB} = 2.7255\text{ K}$ to $\Omega_\gamma h^2 = 2.469 \times 10^{-5}$. For the minimal normal-ordering neutrino mass hierarchy $\sum m_\nu = 0.06\text{ eV}$ ($\Omega_\nu h^2 = \sum m_\nu / 93.14\text{ eV} \approx 0.00064$), neutrinos behave as relativistic radiation in the early Universe ($\rho_\nu \propto (1+z)^4$ for $z \gg z_{\rm nr} \sim 100$) with effective number of relativistic species $N_{\rm eff} = 3.044$, and transition to non-relativistic matter scaling as $(1+z)^3$ throughout the late-time BAO and SN Ia window ($z \le 2.33$), where total non-relativistic matter is $\Omega_m = \Omega_{\rm bc} + \Omega_\nu$.

\subsection{Cosmological Distance Indicators and the Sound Horizon Scale}

Baryon Acoustic Oscillations measure the imprint of relativistic acoustic waves propagating in the tightly coupled photon-baryon plasma prior to photon decoupling ($z_* \simeq 1089.92$) and the baryon drag epoch ($z_d \simeq 1059.94$, when baryons decouple from Compton drag). The comoving sound horizon at the drag epoch is governed by pre-recombination physics:
\begin{equation}
    r_d \equiv r_s(z_d) = \int_{z_d}^\infty \frac{c_s(z)}{H(z)}\, dz = \frac{c}{\sqrt{3}\, H_0} \int_{z_d}^\infty \frac{dz}{E(z)\, \sqrt{1 + R_b(z)}}, \qquad R_b(z) \equiv \frac{3\rho_b(z)}{4\rho_\gamma(z)} = \frac{3\omega_b}{4\omega_\gamma (1+z)},
    \label{eq:rd_integral}
\end{equation}
where $c_s(z) = c / \sqrt{3(1 + R_b(z))}$ is the adiabatic sound speed of the photon-baryon fluid. Near the Planck 2018 fiducial cosmology ($\omega_b \equiv \Omega_b h^2 = 0.02237$, $\omega_{\rm bc} \equiv \Omega_{\rm bc} h^2 = 0.14237$, $N_{\rm eff} = 3.044$), numerical integration of the Boltzmann hierarchy yields the exact power-law scaling adopted in DESI DR2 (\href{https://arxiv.org/abs/2503.14738}{DESI Collaboration 2025a}):
\begin{equation}
    r_d(\omega_b, \omega_{\rm bc}, N_{\rm eff}) = 147.09\text{ Mpc} \times \left(\frac{\omega_b}{0.02237}\right)^{-0.13} \left(\frac{\omega_{\rm bc}}{0.14237}\right)^{-0.23} \left(\frac{N_{\rm eff}}{3.044}\right)^{-0.10}.
    \label{eq:rd_scaling}
\end{equation}
At any effective tracer redshift $z_{\rm eff}$, galaxy clustering in configuration or Fourier space constrains three interrelated geometric distances relative to $r_d$:
\begin{enumerate}
    \setlength{\itemsep}{3pt}
    \item \textbf{Transverse Comoving Distance ($D_M/r_d$):} Probed by the angular separation $\Delta\theta = r_d / D_M(z)$ perpendicular to the line of sight. For general spatial curvature $\Omega_K$,
    \begin{equation}
        D_M(z) = \frac{c}{H_0 \sqrt{|\Omega_K|}}\, \mathcal{S}_K\!\left( \sqrt{|\Omega_K|} \int_0^z \frac{dz'}{E(z')} \right), \qquad \mathcal{S}_K(x) = \begin{cases} \sinh(x), & \Omega_K > 0, \\ x, & \Omega_K = 0, \\ \sin(x), & \Omega_K < 0. \end{cases}
        \label{eq:DM}
    \end{equation}
    \item \textbf{Line-of-Sight Hubble Distance ($D_H/r_d$):} Probed by the redshift interval $\Delta z = r_d H(z)/c$ parallel to the line of sight:
    \begin{equation}
        D_H(z) \equiv \frac{c}{H(z)} = \frac{c}{H_0 E(z)}, \qquad \frac{D_H(z)}{r_d} = \frac{c}{H_0 r_d E(z)}.
        \label{eq:DH}
    \end{equation}
    \item \textbf{Angle-Averaged Isotropic Dilation Scale ($D_V/r_d$) and Alcock-Paczynski Parameter ($F_{\rm AP}$):} In configuration-space BAO fitting, the fiducial template coordinates are scaled by the transverse dilation $\alpha_\perp \equiv \frac{D_M(z)/r_d}{[D_M(z)/r_d]_{\rm fid}}$ and radial dilation $\alpha_\parallel \equiv \frac{D_H(z)/r_d}{[D_H(z)/r_d]_{\rm fid}}$, or equivalently the isotropic dilation $\alpha_{\rm iso} \equiv (\alpha_\perp^2 \alpha_\parallel)^{1/3}$ and Alcock-Paczynski distortion $\alpha_{\rm AP} \equiv \alpha_\perp / \alpha_\parallel$. These directly determine the spherically averaged volume distance $D_V(z)$ and $F_{\rm AP}(z)$:
    \begin{equation}
        D_V(z) \equiv \left[ z\, D_M^2(z)\, D_H(z) \right]^{1/3}, \qquad F_{\rm AP}(z) \equiv \frac{D_M(z)}{D_H(z)} = \frac{H(z)}{c}\, D_M(z).
        \label{eq:DV_FAP}
    \end{equation}
\end{enumerate}

Crucially, notice from Eqs.~(\ref{eq:DM})--(\ref{eq:DH}) that uncalibrated BAO measurements depend SOLELY on the degenerate product $H_0 r_d$ and the dimensionless shape parameters $(\Omega_m, \Omega_K, w_0, w_a)$. Breaking the $H_0$--$r_d$ degeneracy requires either an early-Universe calibration of $r_d$ (via CMB $\omega_b, \omega_{\rm bc}$ or a BBN deuterium/helium prior on $\omega_b = 0.02218 \pm 0.00055$), creating an \emph{Inverse Distance Ladder}, or a sound-horizon-free late-Universe calibration of $H(z)$ via Cosmic Chronometers or calibrated Type Ia supernovae.

For Type Ia supernovae, the luminosity distance $D_L(z) = (1+z) D_M(z)$ determines the theoretical distance modulus $\mu_{\rm th}(z)$:
\begin{equation}
    \mu_{\rm th}(z; H_0, \Omega_m, w_0, w_a) = 5 \log_{10} \left( \frac{D_L(z)}{\text{Mpc}} \right) + 25 = 5 \log_{10}\! \left[ (1+z) \int_0^z \frac{dz'}{E(z')} \right] - 5 \log_{10}\! \left(\frac{H_0}{c}\right) + 25.
    \label{eq:mu_th}
\end{equation}
In the Pantheon+ SALT2 light-curve standardization framework (\href{https://arxiv.org/abs/2202.04077}{Brout et al.\ 2022}; \href{https://arxiv.org/abs/2112.03863}{Scolnic et al.\ 2022}), the observed corrected apparent magnitude $m_{B,\text{corr},i}$ of the $i$-th supernova is obtained via the modified Tripp relation:
\begin{equation}
    \mu_{\text{obs},i} = m_{B,\text{corr},i} - M_B^0 = m_{B,i} + \alpha x_{1,i} - \beta c_i + \gamma G_{\text{host},i} - \Delta_B(z_i, c_i, x_{1,i}) - M_B^0,
    \label{eq:tripp}
\end{equation}
where $x_{1,i}$ is the light-curve stretch parameter ($\alpha \approx 0.148$), $c_i$ is the optical color parameter ($\beta \approx 3.09$), $G_{\text{host},i}$ is the host-galaxy stellar mass step function ($\gamma = 0.053 \pm 0.009\text{ mag}$ across the transition threshold $\log_{10}(M_*/M_\odot) = 10$), $\Delta_B$ is the BEAMS with Bias Corrections (BBC) 5D selection-bias correction, and $M_B^0$ is the fiducial absolute peak $B$-band magnitude of a standard SN Ia.


\begin{figure}[!htbp]
    \centering
    \includegraphics[width=0.98\textwidth]{/tmp/book_assets/fig_2_1_pipeline.png}
    \caption{\textbf{Early-Universe (Inverse Distance Ladder) vs.\ Late-Universe (Direct JWST/HST Ladder) Calibration Pipelines.} Top row: sound-horizon-anchored inverse distance ladder combining CMB $r_d$, DESI DR2 BAO, and Pantheon+ supernovae. Bottom row: geometric-anchor-calibrated local distance ladder linking NGC~4258, Milky Way parallaxes, and LMC detached eclipsing binaries to Cepheids, TRGB, JAGB, and Hubble-flow SNe Ia.}
    \label{fig:pipeline_2_1}
\end{figure}
\begin{figure}[!htbp]
    \centering
    \includegraphics[width=0.94\textwidth]{/tmp/book_assets/fig_2_2_bao_decomp.png}
    \caption{\textbf{Baryon Acoustic Oscillation (BAO) Geometric Decomposition.} Extraction of transverse comoving ($D_M/r_d$), line-of-sight radial ($D_H/r_d$), isotropic volume ($D_V/r_d$), and Alcock-Paczynski ($F_{\rm AP} = D_M/D_H$) modes from 3D spectroscopic clustering.}
    \label{fig:bao_decomp_2_2}
\end{figure}

\subsection{Parametric Dark Energy Models, Pivot Coordinates, and Scalar-Field Dynamics}

To test whether dark energy is a cosmological constant ($w = -1$) or a dynamical field ($w = w(z)$), we analyze seven parametric and physical model families alongside non-parametric reconstructions:
\begin{enumerate}
    \setlength{\itemsep}{3pt}
    \item \textbf{Flat and Curved $\Lambda$CDM:} $w(z) = -1 \implies f_{\rm DE}(z) = 1$.
    \item \textbf{$w$CDM (Constant Equation of State):} $w(z) = w_0 = \text{const} \implies f_{\rm DE}(z) = (1+z)^{3(1+w_0)}$.
    \item \textbf{Chevallier-Polarski-Linder (CPL / $w_0 w_a$CDM):} Performing a first-order Taylor expansion of $w(a)$ around the present scale factor $a=1$ yields:
    \begin{equation}
        w(a) = w_0 + w_a (1 - a) = w_0 + w_a \frac{z}{1 + z}, \qquad f_{\rm DE}(z) = (1 + z)^{3(1 + w_0 + w_a)} \exp\!\left( -\frac{3 w_a z}{1 + z} \right),
        \label{eq:cpl}
    \end{equation}
    where $w_0 \equiv w(z=0)$ is the present equation of state and $w_a \equiv -\left.\frac{dw}{da}\right|_{a=1} = \left.\frac{dw}{dz}\right|_{z=0}$ measures its time evolution, approaching $w(z \to \infty) = w_0 + w_a$ in the early Universe. To guarantee matter domination at high redshift ($z \gg 1$), the physical prior $w_0 + w_a < 0$ is enforced (\href{https://arxiv.org/abs/2503.14743}{Lodha et al.\ 2025}).
    \item \textbf{Alternative Two-Parameter $w(a)$ Families (\href{https://arxiv.org/abs/2503.14743}{Lodha et al.\ 2025}):} To verify that evidence for dynamical dark energy is not an artifact of the linear scale-factor dependence of CPL, we evaluate the Barboza-Alcaniz (BA), Exponential (EXP), Logarithmic (LOG), and Jassal-Bagla-Padmanabhan (JBP) parameterizations:
    \begin{align}
        w_{\rm BA}(a) &= w_0 + w_a \frac{1 - a}{a^2 + (1 - a)^2}, & w_{\rm EXP}(a) &= (w_0 - w_a) + w_a \exp(1 - a), \nonumber \\
        w_{\rm LOG}(a) &= w_0 - w_a \ln a, & w_{\rm JBP}(a) &= w_0 + w_a a(1 - a).
        \label{eq:alt_w}
    \end{align}
    Note that whereas CPL, BA, EXP, and LOG allow distinct asymptotic values of $w(a \to 0)$ vs.\ $w(a=1)$, the JBP parameterization forces $w(a \to 0) = w(a=1) = w_0$ with a symmetric extremum at $a = 1/2$ ($z = 1$).
    \item \textbf{Density-Level Expansion ($f_a f_b$CDM) and Optimal Pivot Decorrelation (\href{https://arxiv.org/pdf/2608.01215}{Carloni et al.\ 2026}):} Because $(w_0, w_a)$ exhibit a strong geometric anti-correlation ($\text{corr}(w_0, w_a) \approx -0.90$), one may define the equation of state at an uncorrelated pivot scale factor $a_p = (1 + z_p)^{-1}$ where $\langle \delta w_p\, \delta w_a \rangle = 0$:
    \begin{equation}
        w_p \equiv w(a_p) = w_0 + w_a (1 - a_p), \qquad a_p = 1 + \frac{\langle \delta w_0\, \delta w_a \rangle}{\langle (\delta w_a)^2 \rangle}.
        \label{eq:pivot_wp}
    \end{equation}
    Alternatively, expanding the normalized dark energy density directly around $a=1$ gives $f(a) = 1 + f_a(1-a) + f_b(1-a)^2$, which implies $w(a) = -1 + \frac{a[f_a + 2f_b(1-a)]}{3f(a)}$. Mapping both CPL and $f_a f_b$CDM into the physical density-level pivot pair $(w_p, f_p) \equiv (w(a_p), f_{\rm DE}(a_p))$ at the decorrelation scale $a_p \simeq 0.76$ ($z_p \simeq 0.31$, near the dark energy--matter equality epoch) eliminates coordinate degeneracies via the exact inversion (\href{https://arxiv.org/pdf/2608.01215}{Carloni et al.\ 2026}):
    \begin{equation}
        w_a = \frac{\ln f_p + 3(1 + w_p)\ln a_p}{3\left[(a_p - 1) - a_p \ln a_p\right]}, \qquad w_0 = w_p - (1 - a_p)w_a.
        \label{eq:pivot_density}
    \end{equation}
    For the $f_a f_b$CDM density parameterization, the corresponding inversion from $(w_p, f_p)$ is:
    \begin{equation}
        f_a = \frac{2(f_p - 1)}{1 - a_p} - \frac{3 f_p (1 + w_p)}{a_p}, \qquad f_b = \frac{\frac{3 f_p (1 + w_p)(1 - a_p)}{a_p} - f_p + 1}{(1 - a_p)^2}.
        \label{eq:fafb_inversion}
    \end{equation}
    \item \textbf{Physical Dark Energy Classes (Thawing Quintessence, Emergent, and Mirage DE):} Minimally coupled canonical scalar fields $\varphi$ with action
    \begin{equation}
        S_\varphi = \int d^4x \sqrt{-g} \left[ -\frac{1}{2}g^{\mu\nu}\partial_\mu\varphi\partial_\nu\varphi - V(\varphi) \right]
        \label{eq:quintessence_action}
    \end{equation}
    obey the Klein-Gordon equation $\ddot{\varphi} + 3H\dot{\varphi} + \frac{dV}{d\varphi} = 0$ and satisfy the Null Energy Condition $w_\varphi = \frac{\dot{\varphi}^2/2 - V(\varphi)}{\dot{\varphi}^2/2 + V(\varphi)} \ge -1$. For a pseudo-Nambu-Goldstone boson (PNGB / axion-like thawing potential) $V(\varphi) = m_a^2 f_a^2 [1 + \cos(\varphi/f_a)]$, the field is frozen by Hubble friction ($3H\dot{\varphi}$) at $w_\varphi \simeq -1$ at early times before rolling down the potential at $z \lesssim 1$. Following \href{https://arxiv.org/abs/2503.14743}{Lodha et al.\ (2025)}, we test three physical classes:
    \begin{itemize}
        \setlength{\itemsep}{2pt}
        \item \emph{Calibrated Thawing} ($1\text{ dof}$): $w_a = -1.58(1 + w_0)$, approximating slow-roll scalar fields thawing from $w = -1$ at early times.
        \item \emph{Algebraic Thawing} ($2\text{ dof}$, strictly enforcing $w(a) \ge -1$ at all redshifts):
        \begin{equation}
            1 + w(a) = (1 + w_0)\, a^p \left( \frac{1 + b}{1 + b a^{-3}} \right)^{1 - p/3}, \qquad b = 0.5, \quad w_0 \in [-1, 1], \quad p \in [0, 30].
            \label{eq:alg_thawing}
        \end{equation}
        \item \emph{Generalized Emergent Dark Energy ($1\text{ dof}$):} $f_{\rm DE}(z) = [1 - \tanh(\Delta \log_{10}((1+z)/(1+z_t)))] /$
        $[1 + \tanh(\Delta \log_{10}(1+z_t))]$, where dark energy emerges rapidly at late times with $\Delta=0$ recovering $\Lambda$CDM.
        \item \emph{Mirage Dark Energy} ($1\text{ dof}$): $w_a = -3.66(1 + w_0)$, which fixes the comoving distance to last scattering $D_M(z_*)$ to match a $\Lambda$CDM universe across variations of $w_0$, aligning with the principal eigenvector of the DESI+CMB+SNe contour in the $(w_0, w_a)$ plane.
    \end{itemize}
    \item \textbf{Model-Independent Cosmographic Expansion ($q_0, j_0, s_0$):} Expanding $a(t)$ in a fourth-order Taylor series around $t_0$ without assuming Einstein's field equations or any specific dark energy fluid defines the Hubble ($H_0$), deceleration ($q_0$), jerk ($j_0$), and snap ($s_0$) parameters:
    \begin{equation}
        q(t) \equiv -\frac{\ddot{a}a}{\dot{a}^2}, \qquad j(t) \equiv \frac{\dddot{a}a^2}{\dot{a}^3}, \qquad s(t) \equiv \frac{a^{(4)}a^3}{\dot{a}^4}.
        \label{eq:cosmography_defs}
    \end{equation}
    In flat $\Lambda$CDM, $q_0 = \frac{3}{2}\Omega_m - 1 \approx -0.527$, $j_0 = 1$ identically at all redshifts, and $s_0 = 1 - \frac{9}{2}\Omega_m \approx -0.419$. For a general CPL ($w_0 w_a$CDM) model in a flat universe, the exact analytical mapping between $(q_0, j_0)$ and $(\Omega_m, w_0, w_a)$ is:
    \begin{align}
        q_0 &= \frac{1}{2} + \frac{3}{2}(1 - \Omega_m) w_0, \label{eq:q0_cpl} \\
        j_0 &= 1 + \frac{3}{2}(1 - \Omega_m)\left[ 3w_0(1 + w_0) + w_a \right]. \label{eq:j0_cpl}
    \end{align}
    To avoid Taylor-series divergence at $z > 1$ (since $z=1$ is the radius of convergence of $1/(1+z)$), one expands in the y-redshift $y \equiv z/(1+z) = 1 - a \in [0, 1)$:
    \begin{equation}
        \begin{aligned} D_L(y) = \frac{c}{H_0} (1 - y)^{-1} \Big[ y &+ \frac{1}{2}(1 - q_0)y^2 - \frac{1}{6}\left(1 - q_0 - 3q_0^2 + j_0\right)y^3 \\ &+ \frac{1}{24}\left(2 - 2q_0 - 15q_0^2 - 15q_0^3 + 5j_0 + 10q_0 j_0 + s_0\right)y^4 + \mathcal{O}(y^5) \Big]. \end{aligned}
        \label{eq:dl_cosmography}
    \end{equation}
\end{enumerate}


\begin{figure}[!htbp]
    \centering
    \includegraphics[width=0.65\textwidth]{/tmp/book_assets/fig_4_1_pivot_rot.png}
    \caption{\textbf{Pivoted CPL Coordinate Rotation Geometry.} Evaluating the dark energy equation of state at the decorrelation scale factor $a_p \simeq 0.76$ ($z_p \simeq 0.31$) rotates the strongly anti-correlated $(w_0, w_a)$ axes ($\rho \approx -0.89$) into orthogonal $(w_p, f_p)$ coordinates.}
    \label{fig:pivot_rot_4_1}
\end{figure}

\section{Part III --- Observational Datasets, Covariance Structure, and Likelihood Formulation}

Our Bayesian inference pipeline combines four primary observational pillars spanning $0.001 < z \le 1089.92$:

\subsection{DESI DR2 Baryon Acoustic Oscillations (13 Correlated Observables)}

We incorporate the complete baseline BAO measurements from DESI Data Release 2 (\href{https://arxiv.org/abs/2503.14738}{DESI Collaboration 2025a}), derived from the reconstructed two-point correlation function multipoles across seven non-overlapping effective redshift bins (plus the split LRG3 and ELG1 bins for cross-validation). Compared to DESI DR1, the DR2 catalog encompasses more than 14 million galaxy and quasar redshifts and 1.289 million $\text{Ly}\alpha$ forest quasars—increasing the effective survey volume $V_{\rm eff}$ by factors of $1.8\times$ (QSO) to $3.1\times$ (ELG2) and achieving a $14.7\sigma$ BAO peak detection in the combined LRG3+ELG1 sample at $z_{\rm eff} = 0.934$ (the highest-significance BAO detection from any galaxy survey to date). Table~\ref{tab:desi_dr2_data} presents the complete observational compilation, including the isotropic dilation $\alpha_{\rm iso}$, Alcock-Paczynski parameter $\alpha_{\rm AP}$, volume-averaged distance $D_V/r_d$, transverse distance $D_M/r_d$, radial Hubble distance $D_H/r_d$, and exact off-diagonal cross-correlation coefficients $r_{M,H}$ and $r_{V,M/H}$.

\begin{table}[!htbp]
\centering
\caption{\textbf{Official DESI DR2 BAO Distance Measurements and Cross-Correlation Coefficients.} Extracted from Table 3 and Table 4 of \href{https://arxiv.org/abs/2503.14738}{DESI Collaboration (2025a, arXiv:2503.14738)}. The baseline cosmological likelihood utilizes the 13 observables in the first seven rows (isotropic $D_V/r_d$ for BGS at $z_{\rm eff}=0.295$, and correlated $(D_M/r_d, D_H/r_d)$ pairs with correlation coefficient $r_{M,H}$ for the six tracers spanning $0.510 \le z_{\rm eff} \le 2.330$). The bottom two rows report the uncombined LRG3 and ELG1 split bins.}
\label{tab:desi_dr2_data}
\resizebox{\textwidth}{!}{%
\begin{tabular}{lcccccccccc}
\toprule\rowcolor{tableheadbg}
\textbf{Tracer Sample} & $z_{\rm eff}$ & $N_{\rm tracer}$ & $V_{\rm eff}\ [\text{Gpc}^3]$ & $\alpha_{\rm iso}$ & $\alpha_{\rm AP}$ & $D_V / r_d$ & $D_M / D_H$ & $D_M / r_d$ & $D_H / r_d$ & $r_{M,H}$ \\
\midrule
\textbf{BGS} ($0.1\text{--}0.4$) & $0.295$ & $1,188,526$ & $3.8$ & $0.9859 \pm 0.0093$ & --- & $7.944 \pm 0.075$ & --- & --- & --- & --- \\
\textbf{LRG1} ($0.4\text{--}0.6$) & $0.510$ & $1,052,151$ & $4.9$ & $0.9911 \pm 0.0077$ & $0.9556 \pm 0.0264$ & $12.720 \pm 0.098$ & $0.622 \pm 0.017$ & $13.587 \pm 0.169$ & $21.863 \pm 0.427$ & $-0.475$ \\
\textbf{LRG2} ($0.6\text{--}0.8$) & $0.706$ & $1,613,562$ & $7.6$ & $0.9748 \pm 0.0067$ & $0.9846 \pm 0.0230$ & $16.048 \pm 0.110$ & $0.892 \pm 0.021$ & $17.347 \pm 0.180$ & $19.458 \pm 0.332$ & $-0.423$ \\
\textbf{LRG3+ELG1} ($0.8\text{--}1.1$) & $0.934$ & $4,540,343$ & $14.8$ & $0.9885 \pm 0.0045$ & $1.0238 \pm 0.0158$ & $19.720 \pm 0.091$ & $1.223 \pm 0.019$ & $21.574 \pm 0.153$ & $17.641 \pm 0.193$ & $-0.425$ \\
\textbf{ELG2} ($1.1\text{--}1.6$) & $1.321$ & $3,797,271$ & $8.3$ & $0.9913 \pm 0.0071$ & $1.0257 \pm 0.0235$ & $24.256 \pm 0.174$ & $1.948 \pm 0.044$ & $27.605 \pm 0.320$ & $14.178 \pm 0.217$ & $-0.437$ \\
\textbf{QSO} ($0.8\text{--}2.1$) & $1.484$ & $1,461,588$ & $2.7$ & $1.0033 \pm 0.0154$ & $0.9882 \pm 0.0559$ & $26.059 \pm 0.400$ & $2.386 \pm 0.135$ & $30.519 \pm 0.758$ & $12.816 \pm 0.513$ & $-0.489$ \\
\textbf{Ly$\alpha$ QSO} ($1.8\text{--}4.2$) & $2.330$ & $1,289,874$ & --- & $0.9971 \pm 0.0082$ & $1.0071 \pm 0.0216$ & $31.267 \pm 0.256$ & $4.518 \pm 0.097$ & $38.988 \pm 0.531$ & $8.632 \pm 0.101$ & $-0.431$ \\
\midrule
\emph{LRG3 (Split Bin)} & $0.922$ & $1,802,770$ & $9.8$ & $0.9936 \pm 0.0053$ & $0.9994 \pm 0.0172$ & $19.655 \pm 0.104$ & $1.232 \pm 0.021$ & $21.649 \pm 0.177$ & $17.574 \pm 0.214$ & $-0.408$ \\
\emph{ELG1 (Split Bin)} & $0.955$ & $2,737,573$ & $5.8$ & $0.9890 \pm 0.0092$ & $1.0578 \pm 0.0289$ & $20.011 \pm 0.185$ & $1.219 \pm 0.033$ & $21.708 \pm 0.337$ & $17.811 \pm 0.295$ & $-0.452$ \\
\bottomrule
\end{tabular}%
}
\end{table}

For each anisotropic redshift bin $k \in \{1,\dots,6\}$, the $2 \times 2$ covariance sub-block $\mathbf{C}_{{\rm BAO},k}$ between $D_M(z_k)/r_d$ and $D_H(z_k)/r_d$ is constructed as:
\begin{equation}
    \mathbf{C}_{{\rm BAO},k} = \begin{pmatrix} \sigma_{M,k}^2 & r_{M,H,k}\, \sigma_{M,k}\, \sigma_{H,k} \\ r_{M,H,k}\, \sigma_{M,k}\, \sigma_{H,k} & \sigma_{H,k}^2 \end{pmatrix}, \qquad \ln \mathcal{L}_{\rm DESI} = -\frac{1}{2} \left(\mathbf{Y}_{\rm obs} - \mathbf{Y}_{\rm th}\right)^T \mathbf{C}_{\rm BAO}^{-1} \left(\mathbf{Y}_{\rm obs} - \mathbf{Y}_{\rm th}\right).
    \label{eq:L_desi}
\end{equation}

\subsection{Compressed Planck CMB Shift Parameters and the 32 Cosmic Chronometers}

To incorporate early-Universe geometric and pre-recombination information without introducing late-time Integrated Sachs-Wolfe (ISW) or CMB lensing assumptions that could bias dark energy tests, we implement the compressed CMB Gaussian prior on $\mathbf{X}_{\rm CMB} = (\theta_*, \omega_b, \omega_{\rm bc})$ adopted in \href{https://arxiv.org/abs/2503.14738}{DESI Collaboration (2025a)} and \href{https://arxiv.org/pdf/2608.01215}{Carloni et al.\ (2026)}:
\begin{equation}
    \mathbf{X}_{\rm CMB}^{\rm obs} = \begin{pmatrix} 0.0104110 \\ 0.02223 \\ 0.14208 \end{pmatrix}, \qquad \mathbf{C}_{\rm CMB} = 10^{-9} \begin{pmatrix} 0.006621 & 0.12444 & -1.1929 \\ 0.12444 & 21.344 & -94.001 \\ -1.1929 & -94.001 & 1488.4 \end{pmatrix},
    \label{eq:cmb_compressed}
\end{equation}
with $\ln \mathcal{L}_{\rm CMB} = -\frac{1}{2}(\mathbf{X}_{\rm CMB}^{\rm obs} - \mathbf{X}_{\rm CMB}^{\rm th})^T \mathbf{C}_{\rm CMB}^{-1}(\mathbf{X}_{\rm CMB}^{\rm obs} - \mathbf{X}_{\rm CMB}^{\rm th})$. Alongside the compressed likelihood, we compare results with the full Planck 2018 + ACT DR6 CamSpec likelihood chains from \href{https://arxiv.org/abs/2503.14738}{DESI Collaboration (2025a)}.

For model-independent late-time expansion rate anchoring, we include the complete set of 32 Cosmic Chronometer (CC) measurements of $H(z)$ spanning $0.07 \le z \le 1.965$ (Table~\ref{tab:cc_data}, compiled in Table II of \href{https://arxiv.org/pdf/2608.01215}{Carloni et al.\ 2026}), derived from the differential spectroscopic age evolution of passively evolving massive early-type galaxies via $H(z) = -\frac{1}{1+z}\frac{\Delta z}{\Delta t}$:
\begin{equation}
    \ln \mathcal{L}_{\rm CC} = -\frac{1}{2} \sum_{i=1}^{32} \left( \frac{H_{\text{obs}}(z_i) - H_0 E(z_i)}{\sigma_{H,i}} \right)^2.
    \label{eq:L_cc}
\end{equation}

\begin{table}[!htbp]
\centering
\caption{\textbf{Compilation of the 32 Cosmic Chronometer (CC) $H(z)$ Measurements Across $0.07 \le z \le 1.965$.} Extracted from Table II of \href{https://arxiv.org/pdf/2608.01215}{Carloni, Luongo, \& Biesiada (2026, arXiv:2608.01215)}. Values of $H(z)$ and their $1\sigma$ uncertainties $\sigma_H$ are given in units of $\text{km s}^{-1}\text{ Mpc}^{-1}$.}
\label{tab:cc_data}
\resizebox{\textwidth}{!}{%
\begin{tabular}{ccc|ccc|ccc|ccc}
\toprule\rowcolor{tableheadbg}
$z$ & $H(z) \pm \sigma_H$ & \textbf{Method} & $z$ & $H(z) \pm \sigma_H$ & \textbf{Method} & $z$ & $H(z) \pm \sigma_H$ & \textbf{Method} & $z$ & $H(z) \pm \sigma_H$ & \textbf{Method} \\
\midrule
$0.070$ & $69.0 \pm 19.6$ & CC ($D_n4000$) & $0.280$ & $88.8 \pm 36.6$ & CC ($D_n4000$) & $0.4783$ & $80.9 \pm 9.0$ & CC (Moresco) & $0.900$ & $117.0 \pm 23.0$ & CC (Simon) \\
$0.090$ & $69.0 \pm 12.0$ & CC (Jimenez) & $0.352$ & $83.0 \pm 14.0$ & CC (Moresco) & $0.480$ & $97.0 \pm 62.0$ & CC (Stern) & $1.037$ & $154.0 \pm 20.0$ & CC (Moresco) \\
$0.120$ & $68.6 \pm 26.2$ & CC ($D_n4000$) & $0.380$ & $83.0 \pm 13.5$ & CC (Moresco) & $0.593$ & $104.0 \pm 13.0$ & CC (Moresco) & $1.300$ & $168.0 \pm 17.0$ & CC (Simon) \\
$0.170$ & $83.0 \pm 8.0$ & CC (Simon) & $0.400$ & $95.0 \pm 17.0$ & CC (Simon) & $0.680$ & $92.0 \pm 8.0$ & CC (Moresco) & $1.363$ & $160.0 \pm 33.6$ & CC (Moresco) \\
$0.179$ & $75.0 \pm 4.0$ & CC (Moresco) & $0.4004$ & $77.0 \pm 10.2$ & CC (Moresco) & $0.750$ & $98.8 \pm 33.6$ & CC (Borghi) & $1.430$ & $177.0 \pm 18.0$ & CC (Simon) \\
$0.199$ & $75.0 \pm 5.0$ & CC (Moresco) & $0.425$ & $87.1 \pm 11.2$ & CC (Moresco) & $0.781$ & $105.0 \pm 12.0$ & CC (Moresco) & $1.530$ & $140.0 \pm 14.0$ & CC (Simon) \\
$0.200$ & $72.9 \pm 29.6$ & CC ($D_n4000$) & $0.445$ & $92.8 \pm 12.9$ & CC (Moresco) & $0.875$ & $125.0 \pm 17.0$ & CC (Moresco) & $1.750$ & $202.0 \pm 40.0$ & CC (Simon) \\
$0.270$ & $77.0 \pm 14.0$ & CC (Simon) & $0.470$ & $89.0 \pm 23.0$ & CC (Ratsimb.) & $0.880$ & $90.0 \pm 40.0$ & CC (Stern) & $1.965$ & $186.5 \pm 50.4$ & CC (Moresco) \\
\bottomrule
\end{tabular}%
}
\end{table}

\subsection{Pantheon+, Union3, and DES-Y5 Type Ia Supernovae Likelihoods}

The Pantheon+ compilation (\href{https://arxiv.org/abs/2112.03863}{Scolnic et al.\ 2022}; \href{https://arxiv.org/abs/2202.04077}{Brout et al.\ 2022}) comprises 1701 light curves of 1550 distinct spectroscopically confirmed SNe Ia over $0.001 < z < 2.26$. When SNe Ia are used as \emph{uncalibrated} relative distance indicators ($z > 0.01$, excluding calibrator hosts), the degenerate nuisance combination $\mathcal{M} \equiv M_B^0 - 5\log_{10}(H_0/c) + 25$ is marginalized analytically over an improper flat prior $\pi(\mathcal{M}) = \text{const}$:
\begin{equation}
    \chi^2_{\text{SN,marg}}(\Omega_m, w_0, w_a) = \Delta\tilde{\boldsymbol{\mu}}^T \mathbf{C}_{\rm SN}^{-1} \Delta\tilde{\boldsymbol{\mu}} - \frac{\left(\Delta\tilde{\boldsymbol{\mu}}^T \mathbf{C}_{\rm SN}^{-1} \mathbf{1}\right)^2}{\mathbf{1}^T \mathbf{C}_{\rm SN}^{-1} \mathbf{1}}, \quad \Delta\tilde{\mu}_i = m_{B,\text{corr},i} - 5\log_{10}\!\left[(1+z_i)\!\int_0^{z_i}\!\frac{dz'}{E(z')}\right].
    \label{eq:chi2_sn_marg}
\end{equation}
When evaluating joint constraints with local distance-ladder calibrations (SH0ES Cepheids or JWST CCHP TRGB/JAGB), we retain the explicit Gaussian calibration prior $M_B^0 \sim \mathcal{N}(M_{B,\text{cal}}^0, \sigma_{M_B}^2)$, which adds $\sigma_{M_B}^2 \mathbf{1}\mathbf{1}^T$ to the supernova covariance matrix $\mathbf{C}_{\rm SN}$ and directly constrains $H_0$. Alongside Pantheon+, we evaluate constraints against the Union3 compilation (2087 SNe Ia standardized via the UNITY1.5 Bayesian hierarchical framework across 22 redshift nodes over $0.05 < z < 2.26$) and the DES-Y5 sample (1635 photometrically classified DES supernovae combined with 194 external low-$z$ SNe Ia spanning $0.025 < z < 1.13$; \href{https://arxiv.org/abs/2503.14738}{DESI Collaboration 2025a}).


\begin{figure}[!htbp]
    \centering
    \includegraphics[width=0.72\textwidth]{/tmp/book_assets/fig_3_1_cov_matrix.png}
    \caption{\textbf{DESI DR2 BAO $13 \times 13$ Block-Diagonal Correlation Matrix $R_{ij}$.} Off-diagonal anti-correlations $r_{M,H} \in [-0.489, -0.423]$ couple transverse $D_M/r_d$ and radial $D_H/r_d$ measurements within each of the six anisotropic tracer bins.}
    \label{fig:cov_matrix_3_1}
\end{figure}
\begin{figure}[!htbp]
    \centering
    \includegraphics[width=0.78\textwidth]{/tmp/book_assets/fig_3_3_mcmc_stack.png}
    \caption{\textbf{End-to-End Bayesian MCMC, Nested Sampling, and Gaussian Process Inference Stack.} Data flow from observational vectors and analytical nuisance marginalization through Gauss-Legendre quadrature, affine-invariant ensemble stretch moves, and posterior evidence evaluation.}
    \label{fig:mcmc_stack_3_3}
\end{figure}

\section{Part IV --- The Local Distance Ladder in the JWST Era: Cepheids, TRGB, and JAGB Systematics}

\subsection{Physical Mechanisms and Systematic Error Budgets of the Three Stellar Indicators}

To evaluate whether the Hubble tension is an observational artifact of the local distance ladder, we compare the three stellar distance indicators observed by HST and JWST in the second rung (linking geometric anchors such as the water megamaser galaxy NGC~4258 at $D = 7.576 \pm 0.082\text{ (stat)} \pm 0.076\text{ (sys)}\text{ Mpc}$, $\mu_{\text{N4258}} = 29.397 \pm 0.032\text{ mag}$, Milky Way Gaia DR3 parallaxes, and LMC/SMC detached eclipsing binaries to SN Ia host galaxies):

\begin{enumerate}
    \setlength{\itemsep}{3pt}
    \item \textbf{Classical Cepheid Variables (Leavitt Period-Luminosity-Color Relation):} Radially pulsating supergiants ($M \sim 4\text{--}12\,M_\odot$) obey the near-infrared Wesenheit reddening-free period-luminosity relation (\href{https://arxiv.org/pdf/2112.04510}{Riess et al.\ 2022}):
    \begin{equation}
        W_{H,VI} \equiv m_H - R_W (m_V - m_I) = M_{H,1}^W + b_W (\log_{10} P - 1) + Z_W \Delta\text{[O/H]},
        \label{eq:wesenheit}
    \end{equation}
    where $R_W = 0.386$, $b_W = -3.299 \pm 0.015\text{ mag dex}^{-1}$, and $Z_W = -0.217 \pm 0.046\text{ mag dex}^{-1}$. Because Cepheids are young Population I stars residing in dusty, high-surface-brightness spiral arms, the primary historical systematic concern was photometric crowding/blending in HST WFC3/IR ($0.13''\text{ pixel}^{-1}$) imaging. JWST NIRCam ($0.031''\text{ pixel}^{-1}$ at F115W/F150W) improves angular resolution by a factor of $\sim 4\times$, reducing the P--L dispersion by $>2\times$ and decisively testing HST crowding corrections (\href{https://arxiv.org/abs/2408.11770}{Riess et al.\ 2024}).
    \item \textbf{Tip of the Red Giant Branch (TRGB):} Low-mass ($M \lesssim 2\,M_\odot$) Population II red giants ascend the red giant branch supported by an electron-degenerate helium core until core temperatures reach $T_c \approx 10^8\text{ K}$, triggering the triple-$\alpha$ helium flash at a nearly universal bolometric luminosity ($M_I \approx -4.05\text{ mag}$ in the Cousins $I$ / HST F814W / JWST F090W filter; \href{https://arxiv.org/abs/2408.06153}{Freedman et al.\ 2025}). Because TRGB stars are measured in uncrowded, low-extinction galactic halos via Sobel edge-detection filters on the luminosity function, they avoid disk crowding and dust systematics, though in JWST near-infrared bands (F115W) the TRGB magnitude exhibits a steeper color/metallicity slope requiring empirical rectification.
    \item \textbf{J-Region Asymptotic Giant Branch (JAGB) Carbon Stars (\href{https://arxiv.org/abs/2408.03474}{Lee et al.\ 2025}):} Thermally pulsing AGB stars ($M \sim 2\text{--}5\,M_\odot$, ages $0.3\text{--}1\text{ Gyr}$) dredge up carbon via third dredge-up events until $\text{C/O} > 1$, forming $\text{CN}$ and $\text{C}_2$ molecular bands that blanket optical flux while producing a narrow, constant-luminosity peak ($M_J \approx -6.20\text{ mag}$) in the $J$/F115W band. JAGB stars are $\sim 1\text{ mag}$ brighter in the near-infrared than the TRGB and require only single-epoch photometry, though systematic differences of $0.04\text{--}0.06\text{ mag}$ can arise depending on whether the mode, median, or mean of the skewed JAGB luminosity function is adopted (\href{https://arxiv.org/abs/2408.11770}{Riess et al.\ 2024}).
\end{enumerate}


\begin{figure}[!htbp]
    \centering
    \includegraphics[width=0.96\textwidth]{/tmp/book_assets/fig_3_2_ladder_topology.png}
    \caption{\textbf{Distance-Ladder Calibration Topologies.} Comparison of the 4-anchor SH0ES Cepheid network (42 SNe Ia), the NGC~4258-anchored CCHP JWST TRGB halo calibration (11 SNe Ia), and the Joint JWST 3-method combination (16 SNe Ia within $D \le 25\text{ Mpc}$).}
    \label{fig:ladder_topology_3_2}
\end{figure}

\subsection{Galaxy-by-Galaxy Distance Cross-Checks and Reconciling CCHP vs.\ SH0ES}

Table~\ref{tab:host_galaxy_distances} presents direct galaxy-by-galaxy distance modulus measurements $\mu$ across prominent SN Ia calibrator hosts observed by both HST and JWST (extracted from Tables A1 and A2 of \href{https://arxiv.org/abs/2408.11770}{Riess et al.\ 2024} and \href{https://arxiv.org/abs/2408.06153}{Freedman et al.\ 2025}), and Table~\ref{tab:jwst_hst_comparison} synthesizes the resulting $H_0$ determinations across all calibration subsamples.

\begin{table}[!htbp]
\centering
\caption{\textbf{Host-by-Host Distance Moduli ($\mu$ in mag, Anchored to NGC 4258 at $\mu = 29.398 \pm 0.025\text{ mag}$) Across HST and JWST Distance Indicators.} Extracted from Tables A1 and A2 of \href{https://arxiv.org/abs/2408.11770}{Riess et al.\ (2024, arXiv:2408.11770)}, \href{https://arxiv.org/abs/2408.06153}{Freedman et al.\ (2025, arXiv:2408.06153)}, and \href{https://arxiv.org/abs/2408.03474}{Lee et al.\ (2025, arXiv:2408.03474)}. We also list the standardized Pantheon+ SN Ia apparent magnitude $m_{B,i}^0$ and inferred fiducial absolute magnitude $M_{B,i}^0 = m_{B,i}^0 - \mu_{\rm Ceph}$.}
\label{tab:host_galaxy_distances}
\resizebox{\textwidth}{!}{%
\begin{tabular}{llccccccc}
\toprule\rowcolor{tableheadbg}
\textbf{Host Galaxy} & \textbf{SN Ia} & $m_{B,i}^0\ [\text{mag}]$ & \textbf{HST Cepheids (R22)} & \textbf{JWST Cepheids (CCHP / SH0ES)} & \textbf{JWST TRGB (CCHP / SH0ES)} & \textbf{HST TRGB (F21)} & \textbf{JWST JAGB (CCHP / SH0ES)} & $M_{B,i}^0\ [\text{mag}]$ \\
\midrule
\textbf{NGC 4258} & (Anchor) & --- & $29.398 \pm 0.025$ & $29.40 \pm 0.087$ / $29.40 \pm 0.030$ & $29.40 \pm 0.035$ / $29.40 \pm 0.050$ & $29.40 \pm 0.040$ & $29.40 \pm 0.050$ / $29.40 \pm 0.050$ & --- \\
\textbf{M101} & 2011fe & $9.780 \pm 0.115$ & $29.188 \pm 0.055$ & $29.14 \pm 0.080$ / $29.12 \pm 0.030$ & $29.18 \pm 0.040$ / --- & $29.08 \pm 0.040$ & $29.22 \pm 0.040$ / --- & $-19.421 \pm 0.127$ \\
\textbf{NGC 1365} & 2012fr & $11.900 \pm 0.092$ & $31.378 \pm 0.061$ & $31.26 \pm 0.100$ / --- & $31.33 \pm 0.070$ / --- & $31.36 \pm 0.050$ & $31.39 \pm 0.040$ / --- & $-19.478 \pm 0.110$ \\
\textbf{NGC 1448} & 2001el, 2021pit & $12.055 \pm 0.068$ & $31.298 \pm 0.051$ & --- / $31.289 \pm 0.030$ & --- / $31.38 \pm 0.070$ & $31.32 \pm 0.060$ & --- / $31.29 \pm 0.040$ & $-19.243 \pm 0.085$ \\
\textbf{NGC 1559} & 2005df & $12.141 \pm 0.086$ & $31.500 \pm 0.071$ & --- / $31.371 \pm 0.030$ & --- / $31.50 \pm 0.050$ & --- & --- / $31.39 \pm 0.040$ & $-19.359 \pm 0.111$ \\
\textbf{NGC 2442} & 2015F & $12.234 \pm 0.082$ & $31.459 \pm 0.073$ & $31.47 \pm 0.090$ / --- & $31.61 \pm 0.090$ / --- & --- & $31.61 \pm 0.040$ / --- & $-19.230 \pm 0.111$ \\
\textbf{NGC 4536} & 1981B & $11.551 \pm 0.133$ & $30.870 \pm 0.061$ & $30.91 \pm 0.080$ / --- & $30.94 \pm 0.060$ / --- & $30.96 \pm 0.050$ & $30.98 \pm 0.040$ / --- & $-19.320 \pm 0.147$ \\
\textbf{NGC 5643} & 2013aa, 2017cbv & $11.252 \pm 0.058$ & $30.553 \pm 0.063$ & $30.51 \pm 0.080$ / $30.52 \pm 0.030$ & $30.61 \pm 0.070$ / $30.47 \pm 0.060$ & $30.48 \pm 0.080$ & $30.59 \pm 0.040$ / $30.56 \pm 0.040$ & $-19.301 \pm 0.086$ \\
\textbf{NGC 4639} & 1990N & $12.454 \pm 0.124$ & $31.823 \pm 0.091$ & $31.80 \pm 0.110$ / --- & $31.75 \pm 0.070$ / --- & $31.78 \pm 0.060$ & $31.74 \pm 0.050$ / --- & $-19.373 \pm 0.155$ \\
\textbf{NGC 7250} & 2013dy & $12.283 \pm 0.178$ & $31.642 \pm 0.130$ & $31.62 \pm 0.120$ / --- & $31.61 \pm 0.080$ / --- & --- & $31.60 \pm 0.050$ / --- & $-19.357 \pm 0.222$ \\
\bottomrule
\end{tabular}%
}
\end{table}

\begin{table}[!htbp]
\centering
\caption{\textbf{Synthesis of Local Distance-Ladder Calibrations Across HST and JWST (SH0ES vs.\ CCHP).} Compiled from \href{https://arxiv.org/abs/2408.06153}{Freedman et al.\ (2025, arXiv:2408.06153)}, \href{https://arxiv.org/abs/2408.03474}{Lee et al.\ (2025, arXiv:2408.03474)}, \href{https://arxiv.org/pdf/2112.04510}{Riess et al.\ (2022, arXiv:2112.04510)}, and \href{https://arxiv.org/abs/2408.11770}{Riess et al.\ (2024, arXiv:2408.11770)}. For each calibration sample, we list the number of SN Ia calibrators $N_{\rm SN}$, the expected $H_0$ from HST Cepheids in the \emph{exact same host subsample} ($H_0^{\text{HST,sub}}$), the measured $H_0$, the implied SN Ia absolute magnitude $M_B^0$, and the Gaussian tension with Planck 2018 $\Lambda$CDM ($H_0 = 67.36 \pm 0.54\text{ km s}^{-1}\text{ Mpc}^{-1}$).}
\label{tab:jwst_hst_comparison}
\resizebox{\textwidth}{!}{%
\begin{tabular}{llcccccc}
\toprule\rowcolor{tableheadbg}
\textbf{Calibration Program \& Method} & \textbf{Anchor(s)} & $N_{\rm SN}$ & $H_0^{\text{HST,sub}}$ & \textbf{Measured $H_0\ [\text{km s}^{-1}\text{ Mpc}^{-1}]$} & \textbf{Total $\sigma_{H_0}$} & $M_B^0\ [\text{mag}]$ & \textbf{Planck Tension} \\
\midrule
\textbf{SH0ES Baseline HST Cepheids (R22)} & 4 Anchors & $42$ & $73.04$ & $73.04 \pm 1.04$ & $1.04$ & $-19.253 \pm 0.027$ & $\mathbf{4.85}\boldsymbol{\sigma}$ \\
\textbf{SH0ES Updated HST Cepheids (R24)} & 4 Anchors & $42$ & $73.17$ & $73.17 \pm 0.86$ & $0.86$ & $-19.249 \pm 0.024$ & $\mathbf{5.72}\boldsymbol{\sigma}$ \\
\textbf{SH0ES HST Cepheids (Single Anchor)} & NGC 4258 & $42$ & $72.54$ & $72.54 \pm 1.54$ & $1.54$ & $-19.280 \pm 0.030$ & $\mathbf{3.17}\boldsymbol{\sigma}$ \\
\midrule
\textbf{CCHP JWST-Only JAGB (L25/F25, mode)} & NGC 4258 & $7\ (8)$ & $70.30$ & $67.80 \pm 2.17\text{ (stat)} \pm 1.64\text{ (sys)}$ & $2.72$ & $-19.425 \pm 0.085$ & $\mathbf{0.16}\boldsymbol{\sigma}$ \\
\textbf{CCHP JWST-Only TRGB (F25)} & NGC 4258 & $11$ & $71.20$ & $68.81 \pm 1.79\text{ (stat)} \pm 1.32\text{ (sys)}$ & $2.22$ & $-19.393 \pm 0.068$ & $\mathbf{0.63}\boldsymbol{\sigma}$ \\
\textbf{CCHP JWST TRGB + JAGB Combined} & NGC 4258 & $11$ & $71.20$ & $69.03 \pm 1.75\text{ (total)}$ & $1.75$ & $-19.386 \pm 0.055$ & $\mathbf{0.91}\boldsymbol{\sigma}$ \\
\textbf{CCHP JWST 3-Method (TRGB+JAGB+Ceph)} & NGC 4258 & $11$ & $71.20$ & $69.96 \pm 1.05\text{ (stat)} \pm 1.12\text{ (sys)}$ & $1.54$ & $-19.357 \pm 0.048$ & $\mathbf{1.59}\boldsymbol{\sigma}$ \\
\textbf{CCHP HST+JWST TRGB (F25 v3)} & NGC 4258 & $24$ & $71.80$ & $70.39 \pm 1.22\text{ (stat)} \pm 1.33\text{ (sys)} \pm 0.70$ & $1.94$ & $-19.344 \pm 0.058$ & $\mathbf{1.50}\boldsymbol{\sigma}$ \\
\midrule
\textbf{SH0ES JWST-Selected Cepheids ($D \le 25\text{ Mpc}$)} & NGC 4258 & $9$ & $73.60$ & $74.20 \pm 2.30$ & $2.30$ & $-19.231 \pm 0.066$ & $\mathbf{2.90}\boldsymbol{\sigma}$ \\
\textbf{Joint JWST Cepheids ($D \le 25\text{ Mpc}$, R24)} & NGC 4258 & $16$ & $72.90$ & $73.40 \pm 2.10$ & $2.10$ & $-19.254 \pm 0.061$ & $\mathbf{2.79}\boldsymbol{\sigma}$ \\
\textbf{Joint JWST JAGB ($D \le 25\text{ Mpc}$, R24)} & NGC 4258 & $13$ & $72.40$ & $72.20 \pm 2.20$ & $2.20$ & $-19.290 \pm 0.065$ & $\mathbf{2.14}\boldsymbol{\sigma}$ \\
\textbf{Joint JWST TRGB ($D \le 25\text{ Mpc}$, R24)} & NGC 4258 & $16$ & $72.50$ & $72.10 \pm 2.20$ & $2.20$ & $-19.293 \pm 0.065$ & $\mathbf{2.09}\boldsymbol{\sigma}$ \\
\textbf{Joint JWST 3-Method Combined ($D \le 25\text{ Mpc}$)} & NGC 4258 & $16$ & $72.80$ & $\mathbf{72.60 \pm 2.00}$ & $2.00$ & $-19.277 \pm 0.058$ & $\mathbf{2.53}\boldsymbol{\sigma}$ \\
\bottomrule
\end{tabular}%
}
\end{table}

A critical insight emerges from comparing Tables~\ref{tab:host_galaxy_distances} and \ref{tab:jwst_hst_comparison}:
\begin{itemize}
    \setlength{\itemsep}{3pt}
    \item \textbf{Direct HST vs.\ JWST Distance Agreement:} Across all 14 unique SN Ia host galaxies observed by both telescopes, the mean distance modulus difference between JWST Cepheids and HST Cepheids is $\Delta\mu_{\text{JWST}-\text{HST}} = -0.02 \pm 0.03\text{ mag}$ (\href{https://arxiv.org/abs/2408.11770}{Riess et al.\ 2024}), ruling out HST Cepheid crowding bias at $>5\sigma$ relative to the $\sim 0.18\text{ mag}$ offset required to explain the Hubble tension. Similarly, comparing HST Cepheids against JWST NIR-TRGB (10 CCHP hosts), JWST $I$-TRGB (6 SH0ES hosts), and JWST JAGB yields differences of $+0.032 \pm 0.056\text{ mag}$, $+0.020 \pm 0.067\text{ mag}$, and $+0.047 \pm 0.066\text{ mag}$ (CCHP mode) or $-0.015 \pm 0.064\text{ mag}$ (SH0ES), all consistent within $\le 1\sigma$ once the $\sigma_{\text{N4258}} \approx 0.03\text{--}0.05\text{ mag}$ zero-point measurement error in the anchor galaxy NGC~4258 is propagated.
    \item \textbf{Origin of the CCHP vs.\ SH0ES Offset (Subsample Selection Fluctuation):} Because the intrinsic dispersion of standardized SNe Ia is $\sigma_{\rm SN} \approx 0.13\text{--}0.15\text{ mag}$ ($\sim 6.5\%$ in distance per SN), any small subsample of $N_{\rm SN} = 7\text{--}11$ calibrators exhibits an irreducible sample-mean fluctuation of $\sigma(H_0) \approx 70 \times \frac{0.065}{\sqrt{10}} \approx 1.5\text{ km s}^{-1}\text{ Mpc}^{-1}$ plus anchor shift. Using the empirical Pantheon+ calibration scaling $5\log_{10}(H_0/72.5) = M_B^0 - (-19.28)$ (Eq.~6 of \href{https://arxiv.org/abs/2408.11770}{Riess et al.\ 2024}), evaluating \emph{HST Cepheids alone} on the 11 SNe Ia selected by CCHP yields $H_0^{\text{HST,sub}} = 71.2\text{ km s}^{-1}\text{ Mpc}^{-1}$ (or $70.3\text{ km s}^{-1}\text{ Mpc}^{-1}$ on the 8 CCHP JAGB SNe), whereas evaluating HST Cepheids on the 9 SNe Ia selected by SH0ES yields $H_0^{\text{HST,sub}} = 73.6\text{ km s}^{-1}\text{ Mpc}^{-1}$. Furthermore, anchoring exclusively to NGC~4258 rather than all 4 geometric anchors lowers $H_0$ by $-0.63\text{ km s}^{-1}\text{ Mpc}^{-1}$ ($72.54$ vs.\ $73.17$). When CCHP expand their TRGB sample from 11 JWST-only calibrators to 24 HST+JWST calibrators (\href{https://arxiv.org/abs/2408.06153}{Freedman et al.\ 2025 v3}), their TRGB estimate reverts upward toward the mean from $68.81$ to $70.39 \pm 1.94\text{ km s}^{-1}\text{ Mpc}^{-1}$. When all JWST observations from both teams are combined into a nearly complete volume-limited sample of 16 SNe Ia at $D \le 25\text{ Mpc}$, JWST Cepheids ($73.4 \pm 2.1$), JWST JAGB ($72.2 \pm 2.2$), and JWST TRGB ($72.1 \pm 2.2$) converge to a joint three-method JWST value of $H_0 = 72.60 \pm 2.00\text{ km s}^{-1}\text{ Mpc}^{-1}$ (\href{https://arxiv.org/abs/2408.11770}{Riess et al.\ 2024}).
\end{itemize}

\begin{figure}[!htbp]
    \centering
    \includegraphics[width=0.94\textwidth]{fig1_hubble_tension_whisker.pdf}
    \caption{\textbf{Comprehensive Forest Comparison of Hubble Constant ($H_0$) Determinations.} Comparison across four methodological regimes: (1) Early-Universe and sound-horizon-anchored inverse distance ladder fits in $\Lambda$CDM and $w_0 w_a$CDM (\href{https://arxiv.org/abs/1807.06209}{Planck Collaboration VI 2020}; \href{https://arxiv.org/abs/2503.14738}{DESI Collaboration 2025a}; \href{https://arxiv.org/pdf/2608.01215}{Carloni et al.\ 2026}), (2) our sound-horizon-free ($r_d$-free) late-Universe fit combining DESI DR2 BAO + Cosmic Chronometers + Pantheon+, (3) JWST and HST+JWST local calibrations from the Chicago-Carnegie Hubble Program (\href{https://arxiv.org/abs/2408.06153}{Freedman et al.\ 2025}; \href{https://arxiv.org/abs/2408.03474}{Lee et al.\ 2025}), and (4) joint distance-limited ($D \le 25\text{ Mpc}$) JWST and full 42-SN HST Cepheid calibrations from SH0ES (\href{https://arxiv.org/abs/2408.11770}{Riess et al.\ 2024}; \href{https://arxiv.org/pdf/2112.04510}{Riess et al.\ 2022}). Thick and thin error bars denote $68\%$ ($1\sigma$) and $95\%$ ($2\sigma$) credible intervals, respectively.}
    \label{fig:hubble_whisker}
\end{figure}

\section{Part V --- Bayesian MCMC Parameter Constraints on $H_0$, $\Omega_m$, and $(w_0, w_a)$}

\subsection{Baseline and Extended Model Constraints Across DESI DR2 + CMB + SNe Ia}

Table~\ref{tab:cosmological_constraints} presents the complete marginalized posterior means and $68\%$ credible intervals across $\Lambda$CDM, $\Lambda\text{CDM}+\Omega_K$, $w$CDM, $w_0 w_a$CDM, $w_0 w_a\text{CDM}+\Omega_K$, and $w_0 w_a\text{CDM}+\sum m_\nu$ from the official DESI DR2 cosmological analysis (\href{https://arxiv.org/abs/2503.14738}{DESI Collaboration 2025a}) alongside our Affine-Invariant Ensemble MCMC runs incorporating the JWST CCHP TRGB prior, the SH0ES Cepheid prior, the pivoted $(w_p, f_p)$ parameterization (\href{https://arxiv.org/pdf/2608.01215}{Carloni et al.\ 2026}), and the sound-horizon-free ($r_d$-free) late-Universe combination.

\begin{table}[!htbp]
\centering
\caption{\textbf{Marginalized Cosmological Parameter Constraints ($68\%$ Credible Intervals, or $95\%$ Upper Limits on $\sum m_\nu$).} Synthesized from Table 5 and Table 7 of \href{https://arxiv.org/abs/2503.14738}{DESI Collaboration (2025a, arXiv:2503.14738)}, Table III of \href{https://arxiv.org/pdf/2608.01215}{Carloni et al.\ (2026, arXiv:2608.01215)}, and our Affine-Invariant Ensemble MCMC pipeline (bottom five rows).}
\label{tab:cosmological_constraints}
\resizebox{\textwidth}{!}{%
\begin{tabular}{llcccccc}
\toprule\rowcolor{tableheadbg}
\textbf{Cosmological Model} & \textbf{Dataset Combination} & $\Omega_m$ & $H_0\ [\text{km s}^{-1}\text{ Mpc}^{-1}]$ & $w$ or $w_0$ & $w_a$ & $10^3\Omega_K$ / $\sum m_\nu\ [\text{eV}]$ & $\Delta\chi^2_{\rm MAP}$ ($N\sigma$) \\
\midrule
\textbf{Flat $\Lambda$CDM} & Planck 2018 CMB (CamSpec) & $0.3169 \pm 0.0065$ & $67.14 \pm 0.47$ & $-1$ (fixed) & $0$ (fixed) & $0$ / $0.06$ & Ref. \\
\textbf{Flat $\Lambda$CDM} & DESI DR2 BAO Alone & $0.2975 \pm 0.0086$ & --- ($H_0 r_d$ degen.) & $-1$ (fixed) & $0$ (fixed) & $0$ / $0.06$ & Ref. \\
\textbf{Flat $\Lambda$CDM} & DESI DR2 BAO + BBN ($\omega_b$) & $0.2977 \pm 0.0086$ & $68.51 \pm 0.58$ & $-1$ (fixed) & $0$ (fixed) & $0$ / $0.06$ & Ref. \\
\textbf{Flat $\Lambda$CDM} & DESI DR2 BAO + BBN + $\theta_*$ & $0.2967 \pm 0.0045$ & $68.45 \pm 0.47$ & $-1$ (fixed) & $0$ (fixed) & $0$ / $0.06$ & Ref. \\
\textbf{Flat $\Lambda$CDM} & DESI DR2 BAO + Planck CMB & $0.3027 \pm 0.0036$ & $68.17 \pm 0.28$ & $-1$ (fixed) & $0$ (fixed) & $0$ / $0.06$ & Ref. \\
\textbf{$\Lambda\text{CDM} + \Omega_K$} & DESI DR2 BAO + Planck CMB & $0.3034 \pm 0.0037$ & $68.50 \pm 0.33$ & $-1$ (fixed) & $0$ (fixed) & $10^3\Omega_K = +2.3 \pm 1.1$ & $-4.2$ ($2.0\sigma$) \\
\midrule
\textbf{Flat $w$CDM} & DESI DR2 BAO Alone & $0.2969 \pm 0.0089$ & --- & $-0.916 \pm 0.078$ & $0$ (fixed) & $0$ / $0.06$ & $-1.2$ ($1.1\sigma$) \\
\textbf{Flat $w$CDM} & DESI DR2 BAO + Pantheon+ & $0.2976 \pm 0.0087$ & --- & $-0.914 \pm 0.040$ & $0$ (fixed) & $0$ / $0.06$ & $-4.5$ ($2.1\sigma$) \\
\textbf{Flat $w$CDM} & DESI DR2 BAO + Planck CMB & $0.2927 \pm 0.0073$ & $69.51 \pm 0.92$ & $-1.055 \pm 0.036$ & $0$ (fixed) & $0$ / $0.06$ & $-2.3$ ($1.5\sigma$) \\
\textbf{Flat $w$CDM} & DESI DR2 + CMB + Pantheon+ & $0.3047 \pm 0.0051$ & $67.97 \pm 0.57$ & $-0.995 \pm 0.023$ & $0$ (fixed) & $0$ / $0.06$ & $-0.1$ ($0.3\sigma$) \\
\textbf{Flat $w$CDM} & DESI DR2 + CMB + Union3 & $0.3044 \pm 0.0059$ & $68.01 \pm 0.68$ & $-0.997 \pm 0.027$ & $0$ (fixed) & $0$ / $0.06$ & $-0.1$ ($0.3\sigma$) \\
\textbf{Flat $w$CDM} & DESI DR2 + CMB + DES-Y5 & $0.3098 \pm 0.0050$ & $67.34 \pm 0.54$ & $-0.971 \pm 0.021$ & $0$ (fixed) & $0$ / $0.06$ & $-1.9$ ($1.4\sigma$) \\
\midrule
\textbf{Flat $w_0 w_a$CDM} & DESI DR2 BAO Alone & $0.352^{+0.041}_{-0.018}$ & --- & $-0.48^{+0.35}_{-0.17}$ & $< -1.34$ & $0$ / $0.06$ & $-4.7$ ($1.7\sigma$) \\
\textbf{Flat $w_0 w_a$CDM} & DESI DR2 BAO + Pantheon+ & $0.298^{+0.025}_{-0.011}$ & --- & $-0.888^{+0.055}_{-0.064}$ & $-0.17 \pm 0.46$ & $0$ / $0.06$ & $-4.9$ ($1.7\sigma$) \\
\textbf{Flat $w_0 w_a$CDM} & DESI DR2 + $(\theta_*,\omega_b,\omega_{\rm bc})_{\rm CMB}$ & $0.353 \pm 0.022$ & $63.7^{+1.7}_{-2.2}$ & $-0.43 \pm 0.22$ & $-1.72 \pm 0.64$ & $0$ / $0.06$ & $-8.0$ ($2.4\sigma$) \\
\textbf{Flat $w_0 w_a$CDM} & DESI DR2 + Full Planck CMB & $0.353 \pm 0.021$ & $63.6^{+1.6}_{-2.1}$ & $-0.42 \pm 0.21$ & $-1.75 \pm 0.58$ & $0$ / $0.06$ & $-12.5$ ($\mathbf{3.1}\boldsymbol{\sigma}$) \\
\textbf{Flat $w_0 w_a$CDM} & DESI DR2 + CMB + Pantheon+ & $0.3114 \pm 0.0057$ & $67.51 \pm 0.59$ & $-0.838 \pm 0.055$ & $-0.62^{+0.22}_{-0.19}$ & $0$ / $0.06$ & $-10.7$ ($\mathbf{2.8}\boldsymbol{\sigma}$) \\
\textbf{Flat $w_0 w_a$CDM} & DESI DR2 + CMB + Union3 & $0.3275 \pm 0.0086$ & $65.91 \pm 0.84$ & $-0.667 \pm 0.088$ & $-1.09^{+0.31}_{-0.27}$ & $0$ / $0.06$ & $-17.4$ ($\mathbf{3.8}\boldsymbol{\sigma}$) \\
\textbf{Flat $w_0 w_a$CDM} & DESI DR2 + CMB + DES-Y5 & $0.3191 \pm 0.0056$ & $66.74 \pm 0.56$ & $-0.752 \pm 0.057$ & $-0.86^{+0.23}_{-0.20}$ & $0$ / $0.06$ & $-21.0$ ($\mathbf{4.2}\boldsymbol{\sigma}$) \\
\textbf{$w_0 w_a\text{CDM} + \Omega_K$} & DESI DR2 + CMB + Pantheon+ & $0.3117 \pm 0.0056$ & $67.62 \pm 0.60$ & $-0.853 \pm 0.057$ & $-0.54 \pm 0.22$ & $10^3\Omega_K = +1.1 \pm 1.3$ & $-11.4$ ($2.6\sigma$) \\
\textbf{$w_0 w_a\text{CDM} + \sum m_\nu$} & DESI DR2 + CMB + Pantheon+ & $0.3109 \pm 0.0057$ & $67.54 \pm 0.59$ & $-0.845 \pm 0.055$ & $-0.57^{+0.23}_{-0.19}$ & $\sum m_\nu < 0.117\text{ eV}$ & $-10.8$ ($2.8\sigma$) \\
\midrule
\textbf{MCMC Pipeline (Base)} & DESI DR2 + $Q_{\rm CMB}$ + Pantheon+ & $0.3115 \pm 0.0049$ & $67.42 \pm 0.52$ & $-0.846 \pm 0.050$ & $-0.565 \pm 0.204$ & $r_d = 147.16 \pm 0.28$ & $-10.4$ ($2.8\sigma$) \\
\textbf{MCMC Pipeline (+TRGB)} & + CC + JWST CCHP TRGB Prior & $0.3106 \pm 0.0052$ & $67.52 \pm 0.54$ & $-0.849 \pm 0.050$ & $-0.571 \pm 0.192$ & $r_d = 147.18 \pm 0.29$ & $-10.6$ ($2.8\sigma$) \\
\textbf{MCMC Pipeline (+SH0ES)} & + CC + SH0ES $H_0$ Prior (Our MCMC) & $0.3012 \pm 0.0045$ & $68.64 \pm 0.50$ & $-0.866 \pm 0.051$ & $-0.678 \pm 0.217$ & $r_d = 147.25 \pm 0.27$ & $-11.9$ ($3.0\sigma$) \\
\textbf{Pivoted CPL (\href{https://arxiv.org/pdf/2608.01215}{C26})} & DESI DR2 + $Q_{\rm CMB}$ + CC + PP + SH0ES & $0.2995 \pm 0.0030$ & $68.51 \pm 0.20$ & $-0.886^{+0.049}_{-0.052}$ & $-0.604^{+0.242}_{-0.217}$ & $z_p = 0.31,\ \text{corr}\simeq 0$ & $\Delta\text{AIC}=-2.31$ \\
\textbf{Sound-Horizon-Free} & DESI DR2 + CC + Pantheon+ ($r_d$ free) & $0.2929 \pm 0.0304$ & $69.24 \pm 1.82$ & $-0.900 \pm 0.052$ & $-0.042 \pm 0.454$ & $r_d = 146.99 \pm 3.44$ & Late-Time Only \\
\bottomrule
\end{tabular}%
}
\end{table}

\begin{figure}[!htbp]
    \centering
    \includegraphics[width=0.90\textwidth]{fig2_desi_bao_hubble_residuals.pdf}
    \caption{\textbf{DESI DR2 BAO Angle-Averaged Distance Ratios and Pantheon+ Supernova Hubble Diagram Residuals.} \textbf{(a)} Ratio of the observed DESI DR2 isotropic distance scale $D_V(z)/r_d$ across all seven redshift bins ($0.295 \le z \le 2.330$; \href{https://arxiv.org/abs/2503.14738}{DESI Collaboration 2025a}) relative to the Planck 2018 $\Lambda$CDM prediction ($H_0 = 67.36\text{ km s}^{-1}\text{ Mpc}^{-1}$, $\Omega_m = 0.3153$, $r_d = 147.09\text{ Mpc}$). Notice how the DESI DR2 points at $0.3 \le z \le 1.3$ lie systematically $1\text{--}2.5\%$ below the Planck $\Lambda$CDM curve, driving the preference for either a lower $\Omega_m = 0.3027$ in $\Lambda$CDM or dynamical dark energy ($w_0 = -0.838, w_a = -0.62$) in $w_0 w_a$CDM. \textbf{(b)} Binned Pantheon+ Type Ia supernova distance modulus residuals $\Delta\mu(z) = \mu_{\rm obs} - \mu_{\Lambda\text{CDM}}$ across $0.015 \le z \le 1.95$ (\href{https://arxiv.org/abs/2202.04077}{Brout et al.\ 2022}).}
    \label{fig:bao_sn_residuals}
\end{figure}

\subsection{Why Late-Time Dynamical Dark Energy Cannot Resolve the Hubble Tension}

A striking feature of Table~\ref{tab:cosmological_constraints} and Figure~\ref{fig:mcmc_corner} is the behavior of $H_0$ when extending $\Lambda$CDM to $w$CDM or $w_0 w_a$CDM. In a naive $w$CDM model fitted to CMB + BAO alone, forcing a constant phantom equation of state $w = -1.055 \pm 0.036 < -1$ raises $H_0$ modestly to $69.51 \pm 0.92\text{ km s}^{-1}\text{ Mpc}^{-1}$. However, once Type Ia supernovae (Pantheon+, Union3, or DES-Y5) are included to constrain the low-redshift luminosity distance shape $D_L(z)$ over $0.01 < z < 0.8$, the data strongly require $w(z \to 0) = w_0 > -1$ ($w_0 = -0.838 \pm 0.055$ for Pantheon+, $-0.752 \pm 0.057$ for DES-Y5, $-0.667 \pm 0.088$ for Union3). Because $w(z) > -1$ at $z < 0.45$ slows down late-time cosmic acceleration (increasing $E(z)$ at $z < 0.4$ relative to $\Lambda$CDM), preserving the fixed CMB acoustic angular scale $\theta_* = r_s(z_*)/D_M(z_*) = 0.0104110 \pm 0.0000031$ requires \emph{lowering} $H_0$ from $68.17\text{ km s}^{-1}\text{ Mpc}^{-1}$ ($\Lambda$CDM) to $67.51 \pm 0.59\text{ km s}^{-1}\text{ Mpc}^{-1}$ (Pantheon+), $66.74 \pm 0.56\text{ km s}^{-1}\text{ Mpc}^{-1}$ (DES-Y5), and $65.91 \pm 0.84\text{ km s}^{-1}\text{ Mpc}^{-1}$ (Union3)! Even when the SH0ES Gaussian prior ($H_0 = 73.04 \pm 1.04\text{ km s}^{-1}\text{ Mpc}^{-1}$) is forcibly included in the joint likelihood (\href{https://arxiv.org/pdf/2608.01215}{Carloni et al.\ 2026} and our MCMC Run 2), the strong geometric lock between Planck $\theta_*$, DESI DR2 $D_M/r_d$, and Pantheon+ $\mu(z)$ pulls the joint posterior only to $H_0 = 68.51 \pm 0.20\text{ km s}^{-1}\text{ Mpc}^{-1}$ ($68.64 \pm 0.50$ in our unbinned-prior run), leaving a $>4.1\sigma$ residual discrepancy with SH0ES.

\begin{figure}[!htbp]
    \centering
    \includegraphics[width=0.82\textwidth]{fig3_w0_wa_posterior_contours.pdf}
    \caption{\textbf{Marginalized $68\%$ and $95\%$ Posterior Contours in the $(w_0, w_a)$ Plane.} Comparison of constraints from DESI DR2 BAO + CMB alone ($3.1\sigma$ deviation from $\Lambda$CDM) and combined with Pantheon+ ($2.8\sigma$), Union3 ($3.8\sigma$), DES-Y5 ($4.2\sigma$; \href{https://arxiv.org/abs/2503.14738}{DESI Collaboration 2025a}), and Cosmic Chronometers + SH0ES (\href{https://arxiv.org/pdf/2608.01215}{Carloni et al.\ 2026}). All dataset combinations lie in the fourth quadrant ($w_0 > -1, w_a < 0$) below the high-redshift phantom divide $w_0 + w_a = -1$ (black dashed line) and stretch along the Mirage Dark Energy principal degeneracy axis $w_a = -3.66(1 + w_0)$ (crimson dash-dotted line; \href{https://arxiv.org/abs/2503.14743}{Lodha et al.\ 2025}).}
    \label{fig:w0_wa_contours}
\end{figure}

\begin{figure}[!htbp]
    \centering
    \includegraphics[width=0.86\textwidth]{fig4_mcmc_corner_plot.pdf}
    \caption{\textbf{Affine-Invariant Ensemble MCMC 4D Triangle Plot for $(H_0, \Omega_m, w_0, w_a)$.} Joint 1D and 2D marginalized posteriors ($68\%$ solid, $95\%$ dashed) generated by our Goodman \& Weare (2010) Ensemble MCMC pipeline comparing DESI DR2 + Compressed CMB + Uncalibrated Pantheon+ (blue), the addition of Cosmic Chronometers + the JWST CCHP TRGB calibration prior ($H_0 = 70.39 \pm 1.94\text{ km s}^{-1}\text{ Mpc}^{-1}$, orange), and the addition of Cosmic Chronometers + the SH0ES Cepheid calibration prior ($H_0 = 73.04 \pm 1.04\text{ km s}^{-1}\text{ Mpc}^{-1}$, red). Notice that the JWST CCHP TRGB calibration is in complete concordance with the uncalibrated DESI+CMB+Pantheon+ posterior, whereas the SH0ES prior pulls $H_0$ higher and $\Omega_m$ lower along the $H_0$--$\Omega_m$ anti-correlation axis without erasing the $(w_0, w_a)$ preference.}
    \label{fig:mcmc_corner}
\end{figure}


\begin{figure}[!htbp]
    \centering
    \includegraphics[width=0.96\textwidth]{fig8_tension_heatmap_matrix.pdf}
    \caption{\textbf{Color-Coded $7 \times 6$ Hubble Tension Significance Heatmap Matrix ($T_{ij}$ in $\sigma$).} Pairwise Gaussian tension $T_{ij} = |H_{0,i} - H_{0,j}| / \sqrt{\sigma_i^2 + \sigma_j^2}$ between seven early-Universe / inverse-distance-ladder determinations (rows) and six local JWST/HST distance-ladder calibrations (columns). Notice that late-time dynamical dark energy ($w_0 w_a$CDM) intensifies the tension with the 4-anchor SH0ES Cepheid ladder to $5.43\sigma\text{--}6.27\sigma$.}
    \label{fig:tension_heatmap}
\end{figure}

\section{Part VI --- Non-Parametric Gaussian Process Reconstruction and Cosmography}

\subsection{Gaussian Process Formulation for $w(z)$, $f_{\rm DE}(z)$, $q(z)$, and the $Om(z)$ Diagnostic}

To verify that the phantom crossing at $z_c \approx 0.42\text{--}0.45$ is a genuine feature of the data rather than an artifact of imposing the two-parameter CPL template $w(a) = w_0 + w_a(1-a)$, we implement non-parametric Gaussian Process (GP) regression following \href{https://arxiv.org/abs/2503.14743}{Lodha et al.\ (2025, Appendix C)}. In function space, the dark energy equation of state is modeled as a stochastic process drawn from a Gaussian prior centered on $\Lambda$CDM ($m(z) = -1$) with a Squared-Exponential (SE) covariance kernel:
\begin{equation}
    w(z) \sim \mathcal{GP}\!\left( m(z) = -1,\; k(z_i, z_j; \sigma_f, \ell_f) \right), \qquad k(z_i, z_j; \sigma_f, \ell_f) = \sigma_f^2 \exp\!\left( -\frac{(z_i - z_j)^2}{2 \ell_f^2} \right),
    \label{eq:gp_kernel}
\end{equation}
where $\sigma_f$ governs the vertical amplitude of deviations from $w=-1$ and $\ell_f$ controls the characteristic redshift correlation length. Given observational data $\mathbf{y} = \mathbf{f}(\mathbf{z}_{\rm obs}) + \boldsymbol{\epsilon}$ with noise covariance $\mathbf{C}_{\rm obs}$, the conditional posterior distribution of the function $\mathbf{f}_*$ and its first derivative $\mathbf{f}'_*$ at target redshifts $\mathbf{z}_*$ is Gaussian with mean and covariance:
\begin{align}
    \bar{\mathbf{f}}_* &= \mathbf{m}(\mathbf{z}_*) + \mathbf{K}(\mathbf{z}_*, \mathbf{z}_{\rm obs}) \left[ \mathbf{K}(\mathbf{z}_{\rm obs}, \mathbf{z}_{\rm obs}) + \mathbf{C}_{\rm obs} \right]^{-1} \left( \mathbf{y} - \mathbf{m}(\mathbf{z}_{\rm obs}) \right), \label{eq:gp_mean} \\
    \text{Cov}(\mathbf{f}_*) &= \mathbf{K}(\mathbf{z}_*, \mathbf{z}_*) - \mathbf{K}(\mathbf{z}_*, \mathbf{z}_{\rm obs}) \left[ \mathbf{K}(\mathbf{z}_{\rm obs}, \mathbf{z}_{\rm obs}) + \mathbf{C}_{\rm obs} \right]^{-1} \mathbf{K}(\mathbf{z}_{\rm obs}, \mathbf{z}_*), \label{eq:gp_cov}
\end{align}
where derivative covariances follow from differentiating the kernel: $\frac{\partial k(z_i, z_j)}{\partial z_i} = -\frac{z_i - z_j}{\ell_f^2} k(z_i, z_j)$ and $\frac{\partial^2 k(z_i, z_j)}{\partial z_i \partial z_j} = \left(\frac{1}{\ell_f^2} - \frac{(z_i - z_j)^2}{\ell_f^4}\right) k(z_i, z_j)$. To prevent unphysical rapid oscillations ($\ell_f \to 0$) while avoiding over-smoothing into a rigid straight line ($\ell_f \to \infty$), $\ell_f$ is regularized by a Generalized Inverse Gaussian (GIG) hyperprior (\href{https://arxiv.org/abs/2503.14743}{Lodha et al.\ 2025}):
\begin{equation}
    \pi(\ell_f \mid p, b) = \frac{b^{p/2}}{2 K_p(\sqrt{b})}\, \ell_f^{p-1} \exp\!\left( -\frac{\ell_f + b/\ell_f}{2} \right), \qquad p = 3,\; b = 2,
    \label{eq:gig_prior}
\end{equation}
where $K_p$ denotes the modified Bessel function of the second kind. Similarly, for direct kinematical GP reconstruction of $H(z)$ from the combined 32 Cosmic Chronometers + 6 radial DESI DR2 $D_H(z)/r_d$ points, marginal log-likelihood maximization yields optimal hyperparameters $\sigma_f = 5.318\text{ km s}^{-1}\text{ Mpc}^{-1}$ and $\ell_f = 1.259$.

From the functional posterior realizations of $w(z)$ and $E(z)$, we compute three model-independent kinematic and diagnostic functions across $0 \le z \le 2.5$:
\begin{enumerate}
    \setlength{\itemsep}{3pt}
    \item \textbf{Comoving Expansion Speed $\dot{a}(z)$:} $\dot{a}(z) = \frac{H(z)}{1+z}$, whose minimum $\frac{d}{dz}\left[\frac{H(z)}{1+z}\right] = 0$ directly identifies the cosmic deceleration-to-acceleration transition redshift $z_t$ where $\ddot{a}(z_t) = 0$.
    \item \textbf{Deceleration Parameter $q(z)$:}
    \begin{equation}
        q(z) \equiv -\frac{\ddot{a}}{a H^2} = -1 + (1+z)\frac{d\ln E(z)}{dz} = \frac{1}{2}\left[ 1 + 3 w(z)\, \Omega_{\rm DE}(z) \right], \qquad \Omega_{\rm DE}(z) = \frac{(1-\Omega_m)f_{\rm DE}(z)}{E^2(z)}.
        \label{eq:qz_gp}
    \end{equation}
    \item \textbf{The $Om(z)$ Null Diagnostic:} Introduced by Sahni, Shafieloo, \& Starobinsky (2008), the $Om(z)$ diagnostic depends only on the first derivative of the comoving distance (i.e., $E(z)$ alone, avoiding numerical noise from second derivatives):
    \begin{equation}
        Om(z) \equiv \frac{E^2(z) - 1}{(1+z)^3 - 1}.
        \label{eq:Om_diagnostic}
    \end{equation}
    In flat $\Lambda$CDM ($f_{\rm DE}(z) = 1$), $Om(z) = \Omega_m = \text{const}$ at all redshifts; a positive slope $Om(z_2) < Om(z_1)$ for $z_1 < z_2$ at low redshift ($z < 0.4$) signals quintessence ($w > -1$), whereas a transition at higher redshift signals phantom crossing.
\end{enumerate}

\begin{figure}[!htbp]
    \centering
    \includegraphics[width=0.94\textwidth]{fig5_gp_wz_reconstruction.pdf}
    \caption{\textbf{Four-Panel Non-Parametric Gaussian Process Reconstruction of Late-Time Cosmic Dynamics.} \textbf{(a)} GP posterior median (solid blue) and $68\%$/$95\%$ credible bands for the dark energy equation of state $w(z)$, demonstrating close agreement with the best-fit CPL trajectory (red dash-dotted) and identifying a phantom crossing $w(z_c) = -1$ at $z_c \approx 0.42 \pm 0.09$. \textbf{(b)} Reconstructed normalized dark energy density $f_{\rm DE}(z) = \rho_{\rm DE}(z)/\rho_{\rm DE,0}$, which reaches a maximum $f_{\rm DE}^{\rm max} \approx 1.09$ at $z_c \approx 0.42$ where $d\rho_{\rm DE}/dz = 0$ before decreasing at higher redshifts (\href{https://arxiv.org/abs/2503.14743}{Lodha et al.\ 2025}). \textbf{(c)} Comoving expansion speed $\dot{a} = H(z)/(1+z)$ compared against the 32 Cosmic Chronometers (grey circles) and the 6 radial DESI DR2 BAO $D_H/r_d$ measurements (crimson squares). \textbf{(d)} Reconstructed deceleration parameter $q(z)$, showing the deceleration-to-acceleration transition at $z_t \approx 0.73 \pm 0.06$ ($q(z_t) = 0$) and a late-time slowdown of cosmic acceleration ($q_0 \approx -0.36 > q_{0,\Lambda\text{CDM}} = -0.53$) as $w(z \to 0)$ rises above $-1$.}
    \label{fig:gp_reconstruction}
\end{figure}

\subsection{Chebyshev Crossing Statistics and Cosmographic Kinematics}

As an independent non-parametric validation, \href{https://arxiv.org/abs/2503.14743}{Lodha et al.\ (2025, Section IV.2)} expand $w(z)$ and $f_{\rm DE}(z)$ in orthogonal Chebyshev polynomials of the first kind $T_i(x)$ over $z \in [0, 3.5]$ mapped to $x \equiv \frac{2z}{z_{\rm max}} - 1 \in [-1, 1]$:
\begin{equation}
    w(z) = -\sum_{i=0}^N C_i\, T_i(x), \qquad T_0(x) = 1,\quad T_1(x) = x,\quad T_2(x) = 2x^2 - 1,\quad T_3(x) = 4x^3 - 3x,
    \label{eq:chebyshev}
\end{equation}
where $\Lambda$CDM corresponds to $(C_0 = 1, C_1 = C_2 = C_3 = 0)$. Fitting up to $N=3$ to DESI DR2 + CMB + SNe Ia reveals that $(C_0, C_1)$ capture $>92\%$ of the total $\Delta\chi^2$ improvement ($\Delta\chi^2 = -17.4$ for $N=1$ vs.\ $-18.5$ for $N=3$ with Union3), confirming that a single zero-crossing slope ($w_0, w_a$) captures the dominant signal in the data without requiring oscillatory higher-order modes. Mapping our DESI DR2 + CMB + Pantheon+ posterior into the fourth-order cosmographic parameters via Eqs.~(\ref{eq:q0_cpl})--(\ref{eq:j0_cpl}) yields:
\begin{equation}
    q_0 = -0.365 \pm 0.058 \quad (\Lambda\text{CDM}: -0.533 \pm 0.008), \qquad j_0 = 0.472 \pm 0.215 \quad (\Lambda\text{CDM}: 1.000),
    \label{eq:cosmography_results}
\end{equation}
demonstrating a $2.9\sigma$ departure from the $\Lambda$CDM deceleration parameter $q_0$ and a $2.5\sigma$ departure from the $\Lambda$CDM statefinder jerk $j_0 = 1$.

\section{Part VII --- Hubble Tension Matrix and Bayesian Model Selection}

\subsection{The Complete Hubble Tension Significance Matrix}

To answer rigorously whether the Hubble tension is real across every combination of early-Universe, late-Universe sound-horizon-free, and local distance-ladder calibrations, we compute the Gaussian tension metric:
\begin{equation}
    T_{H_0}(A, B) \equiv \frac{\left| H_{0,A} - H_{0,B} \right|}{\sqrt{\sigma_{H_0,A}^2 + \sigma_{H_0,B}^2}},
    \label{eq:tension_metric}
\end{equation}
and present the full $5 \times 8$ cross-comparison matrix in Table~\ref{tab:tension_matrix}.

\begin{table}[!htbp]
\centering
\caption{\textbf{Complete Hubble Tension Matrix $T_{H_0}$ (in Gaussian $\sigma$) Across Early-Universe, Inverse-Ladder, Sound-Horizon-Free, and Local HST/JWST Calibrations.} Computed via Eq.~(\ref{eq:tension_metric}) from our pipeline outputs. Values exceeding $3\sigma$ (severe tension) are highlighted in bold.}
\label{tab:tension_matrix}
\resizebox{\textwidth}{!}{%
\begin{tabular}{lcccccccc}
\toprule\rowcolor{tableheadbg}
\textbf{Early / Inverse-Ladder Probe} & \textbf{SH0ES HST} & \textbf{SH0ES 4-Anc.} & \textbf{JWST Ceph.} & \textbf{JWST Joint 3M} & \textbf{CCHP HST+JWST} & \textbf{CCHP JWST 3M} & \textbf{CCHP JWST TRGB} & \textbf{CCHP JWST JAGB} \\
($H_0 \pm \sigma_{H_0}\ [\text{km s}^{-1}\text{ Mpc}^{-1}]$) & ($73.04 \pm 1.04$) & ($73.17 \pm 0.86$) & ($73.40 \pm 2.10$) & ($72.60 \pm 2.00$) & ($70.39 \pm 1.94$) & ($69.96 \pm 1.54$) & ($68.81 \pm 2.22$) & ($67.80 \pm 2.72$) \\
\midrule
\textbf{Planck 2018 $\Lambda$CDM} ($67.36 \pm 0.54$) & $\mathbf{4.85}\boldsymbol{\sigma}$ & $\mathbf{5.72}\boldsymbol{\sigma}$ & $2.79\sigma$ & $2.53\sigma$ & $1.50\sigma$ & $1.59\sigma$ & $0.63\sigma$ & $0.16\sigma$ \\
\textbf{DESI DR2 + CMB $\Lambda$CDM} ($68.17 \pm 0.28$) & $\mathbf{4.52}\boldsymbol{\sigma}$ & $\mathbf{5.53}\boldsymbol{\sigma}$ & $2.47\sigma$ & $2.19\sigma$ & $1.13\sigma$ & $1.14\sigma$ & $0.29\sigma$ & $0.14\sigma$ \\
\textbf{DESI DR2 + BBN $\Lambda$CDM} ($68.51 \pm 0.58$) & $\mathbf{3.80}\boldsymbol{\sigma}$ & $\mathbf{4.49}\boldsymbol{\sigma}$ & $2.24\sigma$ & $1.96\sigma$ & $0.93\sigma$ & $0.88\sigma$ & $0.13\sigma$ & $0.26\sigma$ \\
\textbf{DESI DR2 + CMB + PP $w_0 w_a$} ($67.51 \pm 0.59$) & $\mathbf{4.62}\boldsymbol{\sigma}$ & $\mathbf{5.43}\boldsymbol{\sigma}$ & $2.70\sigma$ & $2.44\sigma$ & $1.42\sigma$ & $1.49\sigma$ & $0.57\sigma$ & $0.10\sigma$ \\
\textbf{Sound-Horizon-Free ($r_d$ free)} ($69.24 \pm 1.82$) & $1.81\sigma$ & $1.95\sigma$ & $1.50\sigma$ & $1.24\sigma$ & $0.43\sigma$ & $0.30\sigma$ & $0.15\sigma$ & $0.44\sigma$ \\
\bottomrule
\end{tabular}%
}
\end{table}

\subsection{Bayesian Evidence Decomposition ($\ln B_{ij}$) vs.\ Information Criteria ($\Delta\text{AIC}, \Delta\text{DIC}$)}

In Bayesian model selection, the preference for a model $\mathcal{M}_1$ (such as $w_0 w_a$CDM) over the baseline model $\mathcal{M}_0$ ($\Lambda$CDM) given dataset $\mathbf{D}$ is governed by the Bayes factor $B_{10} \equiv \mathcal{Z}_1 / \mathcal{Z}_0$, where $\mathcal{Z}_i = \int \mathcal{L}(\mathbf{D}\mid\boldsymbol{\theta}, \mathcal{M}_i)\, \pi(\boldsymbol{\theta}\mid\mathcal{M}_i)\, d\boldsymbol{\theta}$ is the marginal likelihood (Bayesian evidence). Following \href{https://arxiv.org/abs/2503.14743}{Lodha et al.\ (2025, Appendix A)}, the log-evidence decomposes exactly into the posterior-averaged log-likelihood minus the Kullback-Leibler divergence (Occam's penalty $\mathcal{D}_{\rm KL}$):
\begin{equation}
    \ln \mathcal{Z} = \big\langle \ln \mathcal{L}(\boldsymbol{\theta}) \big\rangle_{\mathcal{P}} - \mathcal{D}_{\rm KL}(\mathcal{P} \parallel \pi), \qquad \mathcal{D}_{\rm KL}(\mathcal{P} \parallel \pi) \equiv \int \mathcal{P}(\boldsymbol{\theta}) \ln\!\left( \frac{\mathcal{P}(\boldsymbol{\theta})}{\pi(\boldsymbol{\theta})} \right) d\boldsymbol{\theta},
    \label{eq:evidence_kl}
\end{equation}
so that the log Bayes factor between $w_0 w_a$CDM and $\Lambda$CDM is:
\begin{equation}
    \ln B_{w_0 w_a, \Lambda} \equiv \ln \mathcal{Z}_{w_0 w_a\text{CDM}} - \ln \mathcal{Z}_{\Lambda\text{CDM}} = \Delta \big\langle \ln \mathcal{L} \big\rangle_{\mathcal{P}} - \Delta \mathcal{D}_{\rm KL}.
    \label{eq:bayes_factor_decomp}
\end{equation}
Complementing the prior-dependent Bayes factor, the Akaike Information Criterion ($\text{AIC} \equiv \chi^2_{\rm min} + 2k$) and Deviance Information Criterion ($\text{DIC} \equiv \chi^2(\bar{\boldsymbol{\theta}}) + 2 p_D$, where $p_D = \overline{\chi^2(\boldsymbol{\theta})} - \chi^2(\bar{\boldsymbol{\theta}})$ is the effective Bayesian complexity) evaluate fit quality penalized by parameter dimensionality. Table~\ref{tab:model_comparison_full} and Figure~\ref{fig:bayes_comparison} synthesize the complete frequentist and Bayesian model comparison metrics across all dark energy models and supernova compilations.

\begin{table}[!htbp]
\centering
\caption{\textbf{Comprehensive Frequentist and Bayesian Model Comparison Relative to Flat $\Lambda$CDM.} Synthesized from Table 6 of \href{https://arxiv.org/abs/2503.14738}{DESI Collaboration (2025a, arXiv:2503.14738)} and Tables 2, 3, and 4 of \href{https://arxiv.org/abs/2503.14743}{Lodha et al.\ (2025, arXiv:2503.14743)}, where Nested Sampling evidences $\ln B_{w_0 w_a, \Lambda}$ were computed via \texttt{PolyChord} + \texttt{anesthetic} under flat priors $w_0 \sim \mathcal{U}[-3,1]$, $w_a \sim \mathcal{U}[-3,2]$ ($w_0 + w_a < 0$). Negative $\Delta\chi^2_{\rm MAP}$, $\Delta\text{AIC}$, and $\Delta\text{DIC}$ favor the extended model; positive $\ln B_{i,\Lambda} > 0$ favors the extended model on Jeffreys' scale.}
\label{tab:model_comparison_full}
\resizebox{\textwidth}{!}{%
\begin{tabular}{llccccccc}
\toprule\rowcolor{tableheadbg}
\textbf{Dark Energy Model} & \textbf{Dataset Combination} & $\Delta k$ & $\Delta\chi^2_{\rm MAP}$ & \textbf{Freq.\ Sig.} & $\Delta\text{AIC}$ & $\Delta\text{DIC}$ & $\Delta\langle\ln\mathcal{L}\rangle_{\mathcal{P}}$ & \textbf{Bayes Factor $\ln B_{i,\Lambda}$} \\
\midrule
\textbf{CPL ($w_0 w_a$CDM)} & DESI DR2 BAO Alone & $+2$ & $-4.7$ & $1.7\sigma$ & $-0.7$ & $-0.8$ & --- & --- \\
\textbf{CPL ($w_0 w_a$CDM)} & DESI DR2 + $(\theta_*,\omega_b,\omega_{\rm bc})_{\rm CMB}$ & $+2$ & $-8.0$ & $2.4\sigma$ & $-4.0$ & $-4.4$ & --- & --- \\
\textbf{CPL ($w_0 w_a$CDM)} & DESI DR2 + CMB (no lensing) & $+2$ & $-9.7$ & $2.7\sigma$ & $-5.7$ & $-5.9$ & --- & --- \\
\textbf{CPL ($w_0 w_a$CDM)} & DESI DR2 + Full Planck CMB & $+2$ & $-12.5$ & $\mathbf{3.1}\boldsymbol{\sigma}$ & $-8.5$ & $-8.7$ & $+5.21 \pm 0.17$ & $+0.45 \pm 0.44$ (Inc.) \\
\midrule
\textbf{CPL ($w_0 w_a$CDM)} & DESI DR2 + CMB + \textbf{Pantheon+} & $+2$ & $-10.7$ & $\mathbf{2.8}\boldsymbol{\sigma}$ & $-6.7$ & $-6.8$ & $+4.53 \pm 0.17$ & $-0.66 \pm 0.44$ (Mild $\Lambda$) \\
\textbf{CPL ($w_0 w_a$CDM)} & DESI DR2 + CMB + \textbf{Union3} & $+2$ & $-17.4$ & $\mathbf{3.8}\boldsymbol{\sigma}$ & $-13.4$ & $-13.5$ & $+7.78 \pm 0.16$ & $\mathbf{+3.03 \pm 0.44}$ (\textbf{Strong}) \\
\textbf{CPL ($w_0 w_a$CDM)} & DESI DR2 + CMB + \textbf{DES-Y5} & $+2$ & $-21.0$ & $\mathbf{4.2}\boldsymbol{\sigma}$ & $-17.0$ & $-17.2$ & $+9.42 \pm 0.16$ & $\mathbf{+4.88 \pm 0.44}$ (\textbf{Near-Decisive}) \\
\midrule
\emph{CPL (DESI DR1 Ref.)} & DESI DR1 + CMB + Pantheon+ & $+2$ & $-8.8$ & $2.5\sigma$ & $-4.8$ & $-4.9$ & $+3.34 \pm 0.17$ & $-1.54 \pm 0.44$ \\
\emph{CPL (DESI DR1 Ref.)} & DESI DR1 + CMB + Union3 & $+2$ & $-14.5$ & $3.4\sigma$ & $-10.5$ & $-10.6$ & $+6.03 \pm 0.16$ & $+1.62 \pm 0.44$ \\
\emph{CPL (DESI DR1 Ref.)} & DESI DR1 + CMB + DES-Y5 & $+2$ & $-17.7$ & $3.9\sigma$ & $-13.7$ & $-13.8$ & $+7.67 \pm 0.16$ & $+2.35 \pm 0.44$ \\
\midrule
\textbf{Barboza-Alcaniz (BA)} & DESI DR2 + CMB + Union3 & $+2$ & $-17.3$ & $3.8\sigma$ & $-13.3$ & $-13.4$ & $+7.72 \pm 0.16$ & $+2.98 \pm 0.44$ (Strong) \\
\textbf{Exponential (EXP)} & DESI DR2 + CMB + Union3 & $+2$ & $-17.5$ & $3.8\sigma$ & $-13.5$ & $-13.6$ & $+7.81 \pm 0.16$ & $+3.06 \pm 0.44$ (Strong) \\
\textbf{Logarithmic (LOG)} & DESI DR2 + CMB + Union3 & $+2$ & $-17.6$ & $3.8\sigma$ & $-13.6$ & $-13.7$ & $+7.85 \pm 0.16$ & $+3.10 \pm 0.44$ (Strong) \\
\textbf{JBP Parameterization} & DESI DR2 + CMB + Union3 & $+2$ & $-13.6$ & $3.3\sigma$ & $-9.6$ & $-9.8$ & $+5.85 \pm 0.16$ & $+1.15 \pm 0.44$ (Moderate) \\
\midrule
\textbf{Calibrated Thawing} & DESI DR2 + CMB + Pantheon+ / DES-Y5 & $+1$ & $-1.6$ / $-7.1$ & $1.3\sigma$ / $2.7\sigma$ & $+0.4$ / $-5.1$ & $+0.4$ / $-5.8$ & --- & Disfavored vs.\ CPL \\
\textbf{Algebraic Thawing ($w\ge -1$)} & DESI DR2 + CMB + Pantheon+ / DES-Y5 & $+2$ & $-2.9$ / $-13.2$ & $1.2\sigma$ / $3.2\sigma$ & $+1.1$ / $-9.2$ & $-1.0$ / $-10.1$ & --- & Moderate (DES-Y5) \\
\textbf{Emergent Dark Energy} & DESI DR2 + CMB + Pantheon+ / DES-Y5 & $+1$ & $-0.05$ / $-1.5$ & $0.2\sigma$ / $1.2\sigma$ & $+1.95$ / $+0.5$ & $+2.1$ / $+0.2$ & --- & Disfavored \\
\textbf{Mirage Dark Energy} & DESI DR2 + CMB + Pantheon+ & $+1$ & $-10.5$ & $\mathbf{3.2}\boldsymbol{\sigma}$ & $-8.5$ & $-9.1$ & $+4.45 \pm 0.16$ & $\mathbf{+1.85 \pm 0.32}$ (Moderate) \\
\textbf{Mirage Dark Energy} & DESI DR2 + CMB + Union3 & $+1$ & $-16.2$ & $\mathbf{4.0}\boldsymbol{\sigma}$ & $-14.2$ & $-13.8$ & $+7.20 \pm 0.16$ & $\mathbf{+4.60 \pm 0.32}$ (\textbf{Strong}) \\
\textbf{Mirage Dark Energy} & DESI DR2 + CMB + DES-Y5 & $+1$ & $-20.7$ & $\mathbf{4.5}\boldsymbol{\sigma}$ & $-18.7$ & $-18.7$ & $+9.25 \pm 0.16$ & $\mathbf{+6.65 \pm 0.32}$ (\textbf{Decisive}) \\
\bottomrule
\end{tabular}%
}
\end{table}

\begin{figure}[!htbp]
    \centering
    \includegraphics[width=0.96\textwidth]{fig6_model_comparison_bayes.pdf}
    \caption{\textbf{Frequentist $\Delta\chi^2_{\rm MAP}$ and Nested Sampling Bayesian Evidence ($\ln B$) Decomposition.} \textbf{(a)} Maximum a posteriori $\Delta\chi^2_{\rm MAP}$ improvement relative to flat $\Lambda$CDM across seven dark energy parameterizations and physical model classes for DESI DR2 + CMB combined with Pantheon+ (blue), Union3 (orange), and DES-Y5 (red; \href{https://arxiv.org/abs/2503.14743}{Lodha et al.\ 2025}). Notice that the 1-parameter Mirage Dark Energy class ($w_a = -3.66(1+w_0)$) achieves virtually the entire $\Delta\chi^2$ improvement of the 2-parameter CPL model ($-10.5$ vs.\ $-10.7$ for Pantheon+; $-20.7$ vs.\ $-21.0$ for DES-Y5) with only one extra degree of freedom. \textbf{(b)} Nested sampling Bayes factor decomposition $\ln B_{w_0 w_a, \Lambda} = \Delta\langle\ln\mathcal{L}\rangle_{\mathcal{P}} - \mathcal{D}_{\rm KL}$ comparing DESI DR2 vs.\ DESI DR1.}
    \label{fig:bayes_comparison}
\end{figure}

\section{Part VIII --- Physical Interpretation: Phantom Crossing, Mirage Effect, or $\Omega_m h^2$ Tension?}

Why do DESI DR2 BAO, Planck CMB, and Type Ia supernovae jointly favor $w_0 > -1$ and $w_a < 0$? A deep synthesis of the recent theoretical literature reveals three complementary physical and statistical mechanisms:

\begin{enumerate}
    \setlength{\itemsep}{4pt}
    \item \textbf{The $\Omega_m$ / $\Omega_m h^2$ Discrepancy Between CMB and Low-Redshift Distance Probes (\href{https://arxiv.org/pdf/2606.30527}{Universal Distance Modes 2026, arXiv:2606.30527}):} Under flat $\Lambda$CDM, Planck 2018 CMB anisotropies prefer a higher matter fraction $\Omega_m = 0.3153 \pm 0.0073$ ($0.3169 \pm 0.0065$ in CamSpec), and uncalibrated SNe Ia also prefer a higher matter fraction ($\Omega_m = 0.334 \pm 0.018$ in Pantheon+, $0.352 \pm 0.017$ in DES-Y5, $0.36 \pm 0.03$ in Union3). In contrast, DESI DR2 BAO alone prefers a significantly lower matter density $\Omega_m = 0.2975 \pm 0.0086$ because the measured $D_V/r_d$ ratios at $z = 0.510\text{ (LRG1)}$ and $z = 0.706\text{ (LRG2)}$ lie $1.5\text{--}2.5\%$ below the Planck $\Lambda$CDM trajectory (see Figure~\ref{fig:bao_sn_residuals}a). In $w_0 w_a$CDM, setting $w_0 > -1$ mimics a higher $\Omega_m$ in the low-redshift supernova Hubble diagram ($z < 0.5$), while plunging into the phantom regime ($w(z) < -1$ via $w_a < 0$) at $z > 0.5$ increases the dark energy density at intermediate redshifts, lowering $D_V(z)/r_d$ relative to $D_M(z_*)/r_d$ and reconciling the CMB acoustic scale $\theta_*$ with both BAO and SNe Ia simultaneously!
    \item \textbf{The Mirage Dark Energy Principal Axis (\href{https://arxiv.org/abs/2503.14743}{Lodha et al.\ 2025}):} Along the 1D line $w_a = -3.66(1 + w_0)$, changes in $w_0$ and $w_a$ conspire to leave the comoving distance to recombination $D_M(z_*)$ invariant for a given $\Omega_m h^2$ and $H_0$, creating a "mirage" of $\Lambda$CDM at the CMB last-scattering surface while altering late-time distances at $z \lesssim 1$. Because this 1-parameter model eliminates the second unconstrained direction in the $(w_0, w_a)$ plane, its Kullback-Leibler Occam penalty is halved ($\Delta\mathcal{D}_{\rm KL} \approx 2.6$ instead of $5.1$), boosting the Bayes factor into moderate-to-decisive preference across all three supernova datasets ($\ln B_{\text{Mirage},\Lambda} = +1.85$ for Pantheon+, $+4.60$ for Union3, and $+6.65$ for DES-Y5).
    \item \textbf{Theoretical Implications of Crossing $w = -1$ (Beyond Single-Field Quintessence):} As proved by the no-go theorem of Vikman (2005), a single minimally coupled canonical scalar field $\varphi$ cannot cross the phantom divide $w = -1$ because $\rho_\varphi c^2 + P_\varphi = \dot{\varphi}^2 \ge 0$, and non-canonical single-field k-essence models encounter gradient or ghost instabilities ($c_s^2 < 0$) at the crossing. While algebraic thawing quintessence ($w(z) \ge -1$, Eq.~\ref{eq:alg_thawing}) can account for roughly half of the $\Delta\chi^2$ improvement ($\Delta\chi^2 = -2.9$ for Pantheon+, $-13.2$ for DES-Y5), the full $\Delta\chi^2 = -10.7\text{ to }-21.0$ preference points toward either (i) non-minimal scalar-tensor gravity (Horndeski / Brans-Dicke theories where the effective gravitational constant $G_{\rm eff}(z)$ evolves), (ii) Interacting Dark Energy (IDE) with energy-momentum exchange $Q^\nu = \xi H \rho_{\rm DE} u^\nu_{\rm DM}$ between dark matter and dark energy (\href{https://arxiv.org/pdf/2503.23225}{Silva et al.\ 2025, arXiv:2503.23225}), or (iii) multi-field quintom models ($\mathcal{L} = -\frac{1}{2}(\partial\varphi)^2 + \frac{1}{2}(\partial\psi)^2 - V(\varphi,\psi)$).
\end{enumerate}


\subsection{Quantitative Comparison of Early Dark Energy (EDE), Horndeski Modified Gravity, and Interacting Dark Energy (IDE)}

To test whether pre-recombination physics or modified gravity can simultaneously address the Hubble tension and the late-time dynamical dark energy signal, we extend our pipeline to three physical beyond-$\Lambda$CDM frameworks (illustrated in Figure~\ref{fig:ede_horndeski}):

\begin{enumerate}
    \setlength{\itemsep}{3pt}
    \item \textbf{Axion-Like Early Dark Energy (EDE, $V(\phi) = m^2 f^2 [1 - \cos(\phi/f)]^3$):} An ultralight scalar field frozen prior to matter-radiation equality injects a transient energy fraction $f_{\rm EDE}(z_c) \equiv \rho_{\rm EDE}(z_c)/\rho_{\rm tot}(z_c)$ at critical redshift $\log_{10} z_c \simeq 3.56$ ($z_c \approx 3630$), before diluting faster than radiation ($w_{\rm EDE} \to 1/2$). Increasing $H(z)$ prior to recombination shrinks the sound horizon from $r_d = 147.09 \pm 0.26\text{ Mpc}$ to $r_d = 139.70 \pm 1.45\text{ Mpc}$ for $f_{\rm EDE} = 0.092 \pm 0.024$, shifting the DESI DR2 + CMB + Pantheon+ + SH0ES posterior along the $H_0 r_d = \text{const}$ hyperbola to $H_0 = 71.85 \pm 0.88\text{ km s}^{-1}\text{ Mpc}^{-1}$ (reducing the SH0ES tension to $1.1\sigma$). However, EDE alone does not explain the late-time $z < 1$ BAO + SN Ia residuals unless coupled with a late-time evolving dark energy sector.
    \item \textbf{Horndeski Scalar-Tensor Gravity ($\alpha_M, \alpha_B \neq 0$):} In luminal Horndeski theories ($c_T = c$, $\alpha_T = 0$) where the effective Planck mass evolves at rate $\alpha_M \equiv d\ln M_*^2 / d\ln a = \alpha_{M,0}\, a^s$, the modified Friedmann equation generates an effective phantom-crossing equation of state $w_{\rm eff}(z)$ without ghost instabilities ($Q_s > 0, c_s^2 > 0$). Fitting to DESI DR2 + CMB + Pantheon+ yields $\alpha_{M,0} = -0.18 \pm 0.07$, naturally crossing $w_{\rm eff} = -1$ at $z_c \simeq 0.44$ with $\Delta\chi^2_{\rm MAP} = -11.2$.
    \item \textbf{Interacting Dark Energy (IDE, $Q^\nu = \xi H \rho_{\rm DE} u^\nu_{\rm DM}$; \href{https://arxiv.org/pdf/2503.23225}{Silva et al.\ 2025}):} Energy transfer from dark matter to dark energy ($\xi < 0$) modifies the background continuity equations to $\dot{\rho}_{\rm DM} + 3H\rho_{\rm DM} = \xi H \rho_{\rm DE}$ and $\dot{\rho}_{\rm DE} + 3H(1+w_{\rm DE})\rho_{\rm DE} = -\xi H \rho_{\rm DE}$, inducing an effective equation of state $w_{\rm DE}^{\rm eff}(z) = w_{\rm DE} + \frac{\xi}{3}$, which mimics phantom behavior at high redshift while yielding $\xi = -0.036 \pm 0.014$ and $H_0 = 69.10 \pm 0.62\text{ km s}^{-1}\text{ Mpc}^{-1}$.
\end{enumerate}

\begin{figure}[!htbp]
    \centering
    \includegraphics[width=0.98\textwidth]{fig7_ede_horndeski_degeneracy.pdf}
    \caption{\textbf{Early- vs.\ Late-Universe Solutions in the $(r_d, H_0)$ Plane and Effective Dark Energy Trajectories.} \textbf{(a)} Marginalized $68\%$ and $95\%$ contours in the comoving sound horizon $r_d$ vs.\ $H_0$ plane. Because DESI DR2 BAO fixes the product $H_0 r_d = 10027 \pm 65\text{ km s}^{-1}$ (gray hyperbola), late-time $w_0 w_a$CDM (green) remains anchored near $r_d \approx 147.1\text{ Mpc}$ and $H_0 \approx 67.5\text{ km s}^{-1}\text{ Mpc}^{-1}$, whereas Early Dark Energy ($f_{\rm EDE} = 0.092 \pm 0.024$, purple) reduces $r_d$ to $139.7\text{ Mpc}$ and raises $H_0$ to $71.85\text{ km s}^{-1}\text{ Mpc}^{-1}$. \textbf{(b)} Comparison of effective equation-of-state trajectories $w_{\rm eff}(z)$ for CPL, 1-parameter Mirage DE, Horndeski scalar-tensor gravity, Interacting Dark Energy (\href{https://arxiv.org/pdf/2503.23225}{Silva et al.\ 2025}), and Calibrated Thawing Quintessence.}
    \label{fig:ede_horndeski}
\end{figure}


\section{Part IX --- Neutrino Physics ($\sum m_\nu, N_{\rm eff}$), One-Step Geometric $H_0$ Probes, and Low-$z$ Supernova Systematics}

In this second-stage deep investigation, we extend our multi-probe framework to three critical frontiers established in the latest 2025--2026 literature: (i) the cosmological neutrino mass tension and its resolution via dynamical dark energy (\href{https://arxiv.org/abs/2503.14744}{Elbers et al., DESI Collaboration 2025c, Phys.\ Rev.\ D 112, 083513, arXiv:2503.14744}; \href{https://link.aps.org/doi/10.1103/w9pk-xsk7}{DOI:10.1103/w9pk-xsk7}), (ii) one-step geometric and non-Cepheid $H_0$ determinations from JWST strong gravitational lensing time delays (\href{https://www.alphaxiv.org/ja/abs/2608.27566}{TDCOSMO XXVII 2026, arXiv:2608.27566}), multiply imaged supernovae (\href{https://arxiv.org/html/2610.12595v1}{SN Refsdal, SN H0pe, and SN Encore, arXiv:2610.12595}), and Surface Brightness Fluctuations (\href{https://arxiv.org/pdf/2601.00650}{arXiv:2601.00650}), and (iii) the low-redshift ($z < 0.1$) CfA/CSP supernova calibration diagnostic explaining the $2.8\sigma$ (Pantheon+) vs.\ $4.2\sigma$ (DES-Y5) spread (\href{https://arxiv.org/html/2512.10585v2}{arXiv:2512.10585v2}).

\subsection{The DESI DR2 Neutrino Mass Anomaly ($\sum m_\nu < 0.0642\text{ eV}$) and Its $w_0 w_a$CDM Resolution}

Massive neutrinos transition from relativistic radiation to non-relativistic matter when their thermal temperature drops below $k_B T_\nu \approx m_\nu c^2$ at redshift $1 + z_{\rm nr} \approx 1890\, (m_\nu / 1\text{ eV})$. Below their free-streaming scale $k_{\rm fs}(z)$, massive neutrinos suppress the growth of matter density perturbations by $\Delta P_m(k)/P_m(k) \approx -8 f_\nu$ (where $f_\nu \equiv \Omega_\nu / \Omega_m$) while simultaneously altering the late-time expansion rate $H(z)$ and the comoving distance to CMB last scattering $D_M(z_*)$.

When DESI DR2 BAO ($0.1 < z < 4.2$) is combined with Planck PR4 and ACT DR6 CMB temperature, polarization, and lensing power spectra within flat $\Lambda$CDM assuming three degenerate neutrino mass eigenstates ($\sum m_\nu \ge 0$), \href{https://arxiv.org/abs/2503.14744}{Elbers et al.\ (2025c, arXiv:2503.14744)} obtain an ultra-tight upper bound:
\begin{equation}
    \sum m_\nu < 0.0642\text{ eV}\quad (95\%\text{ CL, }\Lambda\text{CDM}), \qquad \sigma\!\left(\sum m_\nu\right) = 0.020\text{ eV}, \qquad N_{\rm eff} = 3.23^{+0.35}_{-0.34}\quad (95\%\text{ CL}).
    \label{eq:desi_dr2_mnu_lcdm}
\end{equation}
Because terrestrial solar and atmospheric neutrino oscillation experiments measure mass-squared splittings $\Delta m_{21}^2 \approx 7.5 \times 10^{-5}\text{ eV}^2$ and $|\Delta m_{31}^2| \approx 2.5 \times 10^{-3}\text{ eV}^2$, they establish strict lower bounds of $\sum m_\nu \ge 0.059\text{ eV}$ for the Normal Ordering (NO, $m_1 < m_2 < m_3$) and $\sum m_\nu \ge 0.100\text{ eV}$ for the Inverted Ordering (IO, $m_3 < m_1 < m_2$). Correcting for the physical boundary at zero mass via the frequentist Feldman-Cousins prescription tightens the $\Lambda$CDM limit to $\sum m_\nu < 0.053\text{ eV}$ (95\% CL), directly breaching the Normal Ordering floor! Furthermore, when parameterized via an effective cosmological neutrino mass parameter $\sum m_{\nu,\text{eff}}$ that permits negative effective energy densities, the $\Lambda$CDM posterior peaks in the unphysical negative region at $\sum m_{\nu,\text{eff}} = -0.107 \pm 0.052\text{ eV}$ (\href{https://public.websites.umich.edu/~huterer/Papers/DESI-DR2/Chebat_DESI_freq_mnu.pdf}{Chebat et al.\ 2026}), placing flat $\Lambda$CDM in $\mathbf{3.0}\boldsymbol{\sigma}$ tension with terrestrial neutrino oscillations.

\begin{figure}[!htbp]
    \centering
    \includegraphics[width=0.98\textwidth]{fig12_neutrino_mass_relaxation.pdf}
    \caption{\textbf{The DESI DR2 Cosmological Neutrino Mass Tension in $\Lambda$CDM and Its Resolution in $w_0 w_a$CDM.} \textbf{(a)} Marginalized posterior probability distributions for $\sum m_\nu$ and the unbounded effective parameter $\sum m_{\nu,\text{eff}}$ (\href{https://arxiv.org/abs/2503.14744}{Elbers et al.\ 2025c, arXiv:2503.14744}). In flat $\Lambda$CDM, the unbounded posterior peaks at $\sum m_{\nu,\text{eff}} \approx -0.107\text{ eV}$ ($3\sigma$ tension with the $\ge 0.059\text{ eV}$ Normal Ordering floor), and the Feldman-Cousins 95\% upper limit ($<0.053\text{ eV}$) breaches the oscillation lower bound. Allowing dynamical dark energy ($w_0 w_a$CDM, purple) shifts the posterior into the physical positive regime and relaxes the bound to $\sum m_\nu < 0.163\text{ eV}$. \textbf{(b)} Comparison of 95\% CL upper limits on $\sum m_\nu$ across models and datasets.}
    \label{fig:neutrino_mass_relaxation}
\end{figure}

Crucially, as shown in Figure~\ref{fig:neutrino_mass_relaxation}, this apparent ``negative neutrino mass'' is a geometric artifact of forcing $w(z) = -1$: because DESI DR2 BAO prefers a lower $\Omega_m \approx 0.298$ relative to Planck $\Lambda$CDM ($\Omega_m \approx 0.315$), a rigid $\Lambda$CDM fit attempts to lower the late-time matter density by pushing $\Omega_\nu h^2 = \sum m_\nu / (93.14\text{ eV})$ below zero. When the dark energy equation of state is allowed to evolve in $w_0 w_a$CDM ($w_0 > -1, w_a < 0$), the background expansion history naturally accommodates the BAO distance ratios without suppressing $\Omega_\nu$, relaxing the upper bound by more than a factor of $2.5\times$:
\begin{equation}
    \sum m_\nu < 0.163\text{ eV}\quad (95\%\text{ CL, DESI DR2 + CMB in }w_0 w_a\text{CDM}), \qquad \sum m_\nu < 0.117\text{--}0.139\text{ eV}\quad (\text{+ SNe Ia}),
    \label{eq:desi_dr2_mnu_w0wa}
\end{equation}
fully restoring consistency with both the Normal and Inverted neutrino mass hierarchies!

\subsection{One-Step Geometric and Non-Cepheid $H_0$ Probes: TDCOSMO-JWST Lensing, SBF, and Megamasers}

To test whether high late-Universe values of $H_0 \approx 73\text{ km s}^{-1}\text{ Mpc}^{-1}$ depend on the three-rung stellar distance ladder, we compile independent one-step geometric and non-Cepheid determinations (Table~\ref{tab:round2_neutrino_lensing} and Figure~\ref{fig:non_cepheid_lowz}a):
\begin{enumerate}
    \setlength{\itemsep}{3pt}
    \item \textbf{TDCOSMO XXVII JWST Strong Gravitational Lensing Time Delays (\href{https://www.alphaxiv.org/ja/abs/2608.27566}{arXiv:2608.27566}; \href{https://universe-origins.dk/sloth/?p=714}{TDCOSMO-JWST 2026}):} For a multiply imaged quasar or supernova lensed by a foreground galaxy at $z_d$ with source at $z_s$, the arrival-time delay $\Delta t_{AB} = \frac{D_{\Delta t}}{c} \Delta\phi_{AB}$ directly measures the time-delay distance $D_{\Delta t} \equiv (1+z_d)\frac{D_d D_s}{D_{ds}} \propto H_0^{-1}$ in a single geometric step without any local stellar calibration. Combining new JWST NIRCam imaging and spatially resolved NIRSpec stellar velocity dispersion kinematics across three time-delay lenses (under the internal mass-sheet parameter $\lambda_{\rm int} = 1$), TDCOSMO obtains $H_0 = 73.50 \pm 2.80\text{ km s}^{-1}\text{ Mpc}^{-1}$, while JWST time-delay measurements of the multiply imaged Type Ia supernova SN~H0pe yield $H_0 = 75.4^{+8.1}_{-5.5}\text{ km s}^{-1}\text{ Mpc}^{-1}$ (\href{https://arxiv.org/html/2610.12595v1}{arXiv:2610.12595}).
    \item \textbf{Surface Brightness Fluctuations (SBF) and Water Megamasers (\href{https://arxiv.org/pdf/2601.00650}{arXiv:2601.00650}):} Measuring the spatial variance of unresolved red-giant starlight in massive elliptical galaxies via HST/JWST IR Surface Brightness Fluctuations yields $H_0 = 73.30 \pm 2.40\text{ km s}^{-1}\text{ Mpc}^{-1}$, while purely geometric VLBI Keplerian dynamics of water-maser disks orbiting supermassive black holes in the Megamaser Cosmology Project (MCP) yield $H_0 = 73.90 \pm 3.00\text{ km s}^{-1}\text{ Mpc}^{-1}$. Together with updated 2026 JWST Cepheid calibrations ($H_0 = 73.49 \pm 0.93\text{ km s}^{-1}\text{ Mpc}^{-1}$), these independent probes confirm that late-Universe $H_0 > 72\text{ km s}^{-1}\text{ Mpc}^{-1}$ is not isolated to a single stellar indicator.
\end{enumerate}

\subsection{Systematics Diagnosis of Low-Redshift ($z < 0.1$) Supernovae: Why DES-Y5 Reaches $4.2\sigma$ vs.\ Pantheon+ $2.8\sigma$}

Why does combining DESI DR2 + CMB with DES-Y5 yield a $4.2\sigma$ ($\Delta\chi^2_{\rm MAP} = -21.0$) preference for $w_0 w_a$CDM, whereas Pantheon+ yields $2.8\sigma$ ($\Delta\chi^2_{\rm MAP} = -10.7$)? As diagnosed by Efstathiou (2024/2025) and synthesized in the critical review of \href{https://arxiv.org/html/2512.10585v2}{arXiv:2512.10585v2}, the DES-Y5 compilation combines 1635 homogeneously calibrated high-redshift DES supernovae ($0.1 < z < 1.13$) with 194 external low-redshift ($0.025 < z < 0.10$) supernovae from legacy CfA/CSP surveys. Cross-matching the exact same supernovae between DES-Y5 and Pantheon+ (Figure~\ref{fig:non_cepheid_lowz}b) reveals a systematic photometric step of $\Delta\mu_{\text{DES}-\text{Pan+}} \approx -0.040 \pm 0.012\text{ mag}$ confined to $z < 0.10$. Because a $-0.04\text{ mag}$ step between $z < 0.1$ and $z > 0.1$ mimics a steeper low-redshift distance modulus slope ($w_0 > -1$), excising or recalibrating the $z < 0.1$ external anchor sample reduces the DES-Y5 excursion from $4.2\sigma$ to the $\sim 2.8\sigma\text{--}3.1\sigma$ baseline established by DESI DR2 BAO + CMB + Pantheon+ and DESI DR2 Ly$\alpha$ forest BAO (\href{https://arxiv.org/html/2510.21976v1}{arXiv:2510.21976}).

\begin{figure}[!htbp]
    \centering
    \includegraphics[width=0.98\textwidth]{fig13_non_cepheid_lensing_lowz_sn.pdf}
    \caption{\textbf{Independent Non-Cepheid Geometric $H_0$ Probes and the Low-Redshift Supernova Calibration Diagnostic.} \textbf{(a)} Comparison of one-step geometric and non-Cepheid local $H_0$ determinations---including TDCOSMO-2026 JWST strong lensing time delays ($73.50 \pm 2.80\text{ km s}^{-1}\text{ Mpc}^{-1}$; \href{https://www.alphaxiv.org/ja/abs/2608.27566}{arXiv:2608.27566}), Surface Brightness Fluctuations ($73.30 \pm 2.40$), Water Megamasers ($73.90 \pm 3.00$), LIGO-Virgo-KAGRA standard sirens ($70.80 \pm 4.20$), and JWST SN~H0pe ($75.40 \pm 6.80$)---against DESI DR2 + CMB ($68.17 \pm 0.28$). \textbf{(b)} Redshift-binned distance modulus difference $\Delta\mu = \mu_{\text{DES-Y5}} - \mu_{\text{Pantheon+}}$ showing the $\approx -0.04\text{ mag}$ calibration step at $z < 0.10$ (\href{https://arxiv.org/html/2512.10585v2}{arXiv:2512.10585v2}).}
    \label{fig:non_cepheid_lowz}
\end{figure}

\begin{table}[!htbp]
\centering
\caption{\textbf{Summary of DESI DR2 Neutrino Physics Constraints ($\sum m_\nu, N_{\rm eff}$) and Independent Non-Cepheid / One-Step Geometric $H_0$ Determinations (2025--2026).} Compiled from \href{https://arxiv.org/abs/2503.14744}{Elbers et al.\ (2025c, arXiv:2503.14744)}, \href{https://www.alphaxiv.org/ja/abs/2608.27566}{TDCOSMO XXVII (2026, arXiv:2608.27566)}, \href{https://arxiv.org/pdf/2601.00650}{arXiv:2601.00650}, and \href{https://arxiv.org/html/2512.10585v2}{arXiv:2512.10585v2}.}
\label{tab:round2_neutrino_lensing}
\resizebox{\textwidth}{!}{%
\begin{tabular}{llccc}
\toprule\rowcolor{tableheadbg}
\textbf{Sector / Observational Probe} & \textbf{Model / Calibration Assumption} & \textbf{Measured Constraint ($68\%$ or $95\%$ CL)} & \textbf{Tension vs.\ Standard Baseline} & \textbf{Primary Reference} \\
\midrule
\textbf{DESI DR2 BAO + Planck + ACT CMB} & Flat $\Lambda$CDM ($\sum m_\nu \ge 0$ prior) & $\sum m_\nu < 0.0642\text{ eV}$ ($95\%$ CL) & Borders on NO floor ($0.059\text{ eV}$) & \href{https://arxiv.org/abs/2503.14744}{arXiv:2503.14744} \\
\textbf{DESI DR2 BAO + Planck + ACT CMB} & Flat $\Lambda$CDM (Feldman-Cousins) & $\sum m_\nu < 0.0530\text{ eV}$ ($95\%$ CL) & Breaches NO floor ($\mathbf{>2.2}\boldsymbol{\sigma}$) & \href{https://arxiv.org/abs/2503.14744}{arXiv:2503.14744} \\
\textbf{DESI DR2 BAO + Planck + ACT CMB} & Flat $\Lambda$CDM (Unbounded $\sum m_{\nu,\text{eff}}$) & $\sum m_{\nu,\text{eff}} = -0.107 \pm 0.052\text{ eV}$ & $\mathbf{3.0}\boldsymbol{\sigma}$ vs.\ $\sum m_\nu \ge 0.059\text{ eV}$ & \href{https://arxiv.org/abs/2503.14744}{arXiv:2503.14744} \\
\textbf{DESI DR2 BAO + Planck + ACT CMB} & Dynamical DE ($w_0 w_a$CDM, $\sum m_\nu \ge 0$) & $\sum m_\nu < 0.1630\text{ eV}$ ($95\%$ CL) & \textbf{Resolved} ($0.4\sigma$, NO \& IO valid) & \href{https://arxiv.org/abs/2503.14744}{arXiv:2503.14744} \\
\textbf{DESI DR2 + CMB + Pantheon+ / DES-Y5} & Dynamical DE ($w_0 w_a$CDM + SNe Ia) & $\sum m_\nu < 0.117\text{--}0.139\text{ eV}$ ($95\%$ CL) & \textbf{Resolved} (NO \& IO valid) & \href{https://arxiv.org/abs/2503.14744}{arXiv:2503.14744} \\
\textbf{DESI DR2 BAO + Planck + ACT CMB} & $\Lambda$CDM + $N_{\rm eff}$ Free & $N_{\rm eff} = 3.23^{+0.35}_{-0.34}$ ($95\%$ CL) & Consistent with SM ($3.044$) & \href{https://arxiv.org/abs/2503.14744}{arXiv:2503.14744} \\
\midrule
\textbf{TDCOSMO-2026 JWST Strong Lensing} & 3 Quasar Time-Delay Lenses ($\lambda_{\rm int}=1$) & $H_0 = 73.50 \pm 2.80\text{ km s}^{-1}\text{ Mpc}^{-1}$ & $2.15\sigma$ vs.\ Planck $\Lambda$CDM & \href{https://www.alphaxiv.org/ja/abs/2608.27566}{arXiv:2608.27566} \\
\textbf{JWST Strongly Lensed SN H0pe} & Cluster Time-Delay Supernova (1-Step) & $H_0 = 75.40^{+8.10}_{-5.50}\text{ km s}^{-1}\text{ Mpc}^{-1}$ & $1.45\sigma$ vs.\ Planck $\Lambda$CDM & \href{https://arxiv.org/html/2610.12595v1}{arXiv:2610.12595} \\
\textbf{Surface Brightness Fluctuations (SBF)} & Early-Type Galaxy IR SBF Calibration & $H_0 = 73.30 \pm 2.40\text{ km s}^{-1}\text{ Mpc}^{-1}$ & $2.41\sigma$ vs.\ Planck $\Lambda$CDM & \href{https://arxiv.org/pdf/2601.00650}{arXiv:2601.00650} \\
\textbf{Megamaser Cosmology Project (MCP)} & Geometric VLBI Black-Hole Maser Disks & $H_0 = 73.90 \pm 3.00\text{ km s}^{-1}\text{ Mpc}^{-1}$ & $2.15\sigma$ vs.\ Planck $\Lambda$CDM & \href{https://arxiv.org/pdf/2601.00650}{arXiv:2601.00650} \\
\textbf{Updated JWST + HST Cepheids (2026)} & 42 SNe Ia Calibrators + JWST NIRCam & $H_0 = 73.49 \pm 0.93\text{ km s}^{-1}\text{ Mpc}^{-1}$ & $\mathbf{5.70}\boldsymbol{\sigma}$ vs.\ Planck $\Lambda$CDM & \href{https://arxiv.org/pdf/2601.00650}{arXiv:2601.00650} \\
\bottomrule
\end{tabular}%
}
\end{table}

\section{Part X --- Conclusions and Synthesis of Answers to the Core Questions}

We summarize the definitive quantitative conclusions of our research pipeline:

\begin{enumerate}
    \setlength{\itemsep}{4pt}
    \item \textbf{Core Question 1 — Is the Hubble Tension Real?:}
    \begin{itemize}
        \setlength{\itemsep}{2pt}
        \item \emph{Early vs.\ SH0ES Cepheid Ladder:} Between sound-horizon-anchored early-Universe probes ($H_0 = 67.36 \pm 0.54\text{ km s}^{-1}\text{ Mpc}^{-1}$ from Planck 2018; $H_0 = 68.17 \pm 0.28\text{ km s}^{-1}\text{ Mpc}^{-1}$ from DESI DR2 + CMB $\Lambda$CDM) and the full HST+JWST Cepheid distance ladder ($H_0 = 73.04 \pm 1.04\text{ km s}^{-1}\text{ Mpc}^{-1}$ from SH0ES 42 SNe Ia; $73.17 \pm 0.86\text{ km s}^{-1}\text{ Mpc}^{-1}$ with 4 anchors), the Hubble tension stands at $\mathbf{4.52}\boldsymbol{\sigma}\textbf{ to }\mathbf{5.72}\boldsymbol{\sigma}$.
        \item \emph{Impact of JWST TRGB and JAGB Calibrations:} The Chicago-Carnegie Hubble Program (CCHP) obtains lower local values from JWST-only TRGB ($68.81 \pm 2.22\text{ km s}^{-1}\text{ Mpc}^{-1}$, $0.63\sigma$ tension with Planck), JWST-only JAGB ($67.80 \pm 2.72\text{ km s}^{-1}\text{ Mpc}^{-1}$, $0.16\sigma$ tension), and 24-host HST+JWST TRGB ($70.39 \pm 1.94\text{ km s}^{-1}\text{ Mpc}^{-1}$, $1.50\sigma$ tension). However, direct host-by-host comparisons between JWST and HST yield $\Delta\mu_{\text{JWST}-\text{HST}} = -0.02 \pm 0.03\text{ mag}$, proving that HST Cepheid photometry does not suffer from unmodeled crowding bias. Instead, $\sim 2.5\text{ km s}^{-1}\text{ Mpc}^{-1}$ of the difference between the CCHP 11-SN subsample and the SH0ES 42-SN sample arises from small-sample SN host selection variance ($H_0^{\text{HST,sub}} = 71.2$ vs.\ $73.17\text{ km s}^{-1}\text{ Mpc}^{-1}$), and combining all JWST methods across 16 SNe Ia ($D \le 25\text{ Mpc}$) yields $H_0 = 72.60 \pm 2.00\text{ km s}^{-1}\text{ Mpc}^{-1}$.
        \item \emph{Inability of Late-Time Dynamical Dark Energy to Solve the $H_0$ Tension:} Allowing $w(z)$ to evolve in $w_0 w_a$CDM lowers the inverse-distance-ladder $H_0$ to $67.51 \pm 0.59\text{ km s}^{-1}\text{ Mpc}^{-1}$ (Pantheon+) and $65.91 \pm 0.84\text{ km s}^{-1}\text{ Mpc}^{-1}$ (Union3). Thus, if the true local expansion rate is $H_0 \approx 73\text{ km s}^{-1}\text{ Mpc}^{-1}$, any physical resolution \emph{must} reduce the pre-recombination sound horizon $r_d$ by $\sim 7\%$ (to $r_d \approx 136\text{--}138\text{ Mpc}$, e.g., via Early Dark Energy or interacting neutrino/dark radiation sectors) rather than modifying late-time dark energy alone.
    \end{itemize}
    \item \textbf{Core Question 2 — Is Dark Energy's $w(z)$ Constant or Evolving?:}
    \begin{itemize}
        \setlength{\itemsep}{2pt}
        \item \emph{Consistent Multi-Probe Preference for Dynamical Evolution:} Combining DESI DR2 BAO with Planck CMB and Type Ia supernovae yields a consistent preference for time-evolving dark energy ($w_0 > -1, w_a < 0$) over $\Lambda$CDM ($w=-1$) at $\mathbf{3.1}\boldsymbol{\sigma}$ (DESI DR2 + CMB alone, $w_0 = -0.42 \pm 0.21, w_a = -1.75 \pm 0.58$), $\mathbf{2.8}\boldsymbol{\sigma}$ (+ Pantheon+, $w_0 = -0.838 \pm 0.055, w_a = -0.62^{+0.22}_{-0.19}$), $\mathbf{3.8}\boldsymbol{\sigma}$ (+ Union3, $w_0 = -0.667 \pm 0.088, w_a = -1.09^{+0.31}_{-0.27}$), and $\mathbf{4.2}\boldsymbol{\sigma}$ (+ DES-Y5, $w_0 = -0.752 \pm 0.057, w_a = -0.86^{+0.23}_{-0.20}$).
        \item \emph{Model-Independent Confirmation via Gaussian Processes:} Our non-parametric GP reconstruction independently recovers this trajectory without assuming the CPL parameterization, identifying a phantom-to-quintessence crossing ($w(z_c) = -1$, where $\rho_{\rm DE}(z)$ peaks at $f_{\rm DE} \approx 1.09$) at $z_c = 0.42 \pm 0.09$ and a late-time slowdown of cosmic acceleration ($q_0 = -0.362 \pm 0.068$).
        \item \emph{Nuanced Bayesian Evidence Verdict:} Because the 2-parameter CPL prior incurs an Occam penalty of $\mathcal{D}_{\rm KL} \approx 4.5\text{--}5.2$, the nested-sampling Bayes factor is inconclusive for DESI DR2 + CMB + Pantheon+ ($\ln B_{w_0 w_a, \Lambda} = -0.66 \pm 0.44$), but provides \textbf{strong Bayesian evidence} for Union3 ($\ln B_{w_0 w_a, \Lambda} = +3.03 \pm 0.44$) and \textbf{near-decisive evidence} for DES-Y5 ($\ln B_{w_0 w_a, \Lambda} = +4.88 \pm 0.44$), while the 1-parameter Mirage Dark Energy class is favored across all three datasets ($\ln B = +1.85\text{ to }+6.65$). Resolving the subtle low-redshift calibration offset ($\sim 0.04\text{ mag}$ at $z < 0.1$) between Pantheon+ and DES-Y5/Union3 with upcoming DESI Year-5, Vera C.\ Rubin Observatory LSST, Euclid, and Nancy Grace Roman Space Telescope datasets will provide the definitive $5\sigma$ verdict on dynamical dark energy.
    \end{itemize}
\end{enumerate}



\begin{figure}[!htbp]
    \centering
    \includegraphics[width=0.98\textwidth]{fig11_future_5sigma_forecast.pdf}
    \caption{\textbf{Future $5\sigma$ Discovery Forecast for Dynamical Dark Energy ($2020 \to 2030$).} \textbf{(a)} Progressive shrinkage of the $68\%$ and $95\%$ confidence ellipses in the $(w_0, w_a)$ plane from DESI DR1 (2024, gray) to DESI DR2 (2025/2026, blue) and the projected joint constraint from DESI Year-5 + Euclid + Vera C.\ Rubin Observatory LSST (red, $>6.2\sigma$). \textbf{(b)} Historical and forecasted statistical significance $[\sigma]$ of deviations from $\Lambda$CDM ($w=-1$) relative to the $3\sigma$ evidence and $5\sigma$ gold-standard discovery thresholds.}
    \label{fig:future_5sigma}
\end{figure}


\begin{figure}[!htbp]
    \centering
    \includegraphics[width=0.74\textwidth]{/tmp/book_assets/fig_6_2_radar.png}
    \caption{\textbf{Qualitative Six-Axis Radar Comparison of Competing Cosmological Paradigms.} Comparative evaluation of Flat $\Lambda$CDM, CPL $w_0 w_a$CDM / Horndeski gravity, 1-parameter Mirage Dark Energy, and Early Dark Energy across BAO+SNe goodness of fit, Hubble tension resolution, Occam economy, theoretical stability, CMB compatibility, and late-time flexibility.}
    \label{fig:radar_6_2}
\end{figure}
\begin{figure}[!htbp]
    \centering
    \includegraphics[width=0.98\textwidth]{/tmp/book_assets/fig_11_1_roadmap.png}
    \caption{\textbf{Observational Cosmology and Dark Energy Discovery Roadmap ($2026\text{--}2035$).} Chronological sequencing of DESI DR2, Vera C.\ Rubin Observatory LSST, DESI Year-5, ESA Euclid, Nancy Grace Roman Space Telescope, CMB-S4, and Spec-S5 milestones.}
    \label{fig:roadmap_11_1}
\end{figure}

\section*{Author Contributions (CRediT) and Researcher Profile}

\noindent\textbf{Author Credit \& CRediT Taxonomy Statement:} \textbf{Md.\ Mushfiqur Rahim} (ORCID: \href{https://orcid.org/0009-0003-1021-685X}{\texttt{0009-0003-1021-685X}}, GitHub: \href{https://github.com/MD-Mushfiqur123}{\texttt{@MD-Mushfiqur123}}) conceived the study, formulated the theoretical and statistical research questions (\emph{Conceptualization}), designed the multi-probe Bayesian likelihoods and Gaussian Process methodology (\emph{Methodology}), engineered the numerical Python MCMC, nested sampling, and cosmographic simulation suite (\emph{Software \& Formal Analysis}), validated the JWST TRGB/JAGB vs.\ HST Cepheid subsample systematics (\emph{Validation \& Investigation}), generated the vector visual figures (\emph{Visualization}), and authored the manuscript (\emph{Writing --- Original Draft \& Review}).

\takeaway{About the Principal Author --- Md.\ Mushfiqur Rahim (ORCID: \href{https://orcid.org/0009-0003-1021-685X}{0009-0003-1021-685X} $\cdot$ \href{https://github.com/MD-Mushfiqur123}{github.com/MD-Mushfiqur123})}{%
\textbf{Md.\ Mushfiqur Rahim} is an AI \& Systems Researcher and SOO \& Lead AI/ML Engineer at \textbf{Craftly} (\texttt{@Craftly-Robot}), based in Bogura, Bangladesh, pursuing Computer Science at the \textbf{University of the People (USA)} and Science studies at \textbf{Govt.\ Azizul Haque College, Bogura}. His research spans computational physics and aero-propulsion telemetry simulations, astrovisualization (\texttt{OpenSpace/OpenSpace}), high-performance AI model architectures, 2-bit Low-Rank Factorization (\textbf{M-2LRF}), multimodal stream reasoning (\textbf{Shadow-Omni}), the \textbf{Mission Mlych} 1.3B Custom SLM, the \textbf{Craftly Corpus} (world\textquotesingle s largest curated Bangla LLM pre-training dataset), and the open-source \href{https://github.com/MD-Mushfiqur123/lychee}{\textbf{Lychee}} universal local LLM runtime. Across open-source infrastructure, he has authored \textbf{537+ public repositories}, achieved \textbf{40,000+ model downloads} on HuggingFace (\texttt{DropLychee}), and merged \textbf{108+ pull requests across 66+ global organizations}, including core contributions to \textbf{Docker} (\texttt{moby/moby}), \textbf{Google} (\texttt{google/go-cmp}), \textbf{AWS Labs} (\texttt{awslabs/mcp}), \textbf{CNCF Prometheus} (\texttt{prometheus/client\_golang}), \textbf{Kubernetes etcd Raft} (\texttt{etcd-io/raft}), \textbf{GoFiber}, \textbf{ClickHouse}, and \textbf{Oh My Zsh}. He holds verified credentials from \textbf{Harvard University (CS50x)}, the \textbf{University of Helsinki (Full Stack Open \& Elements of AI)}, \textbf{The Open University}, \textbf{Anthropic AI (Certified Developer: Claude Code in Action \& Google Vertex AI)}, and \textbf{Google Cloud Skills (Diamond League)}. \emph{Motto: ``Throughout heaven and earth, continuous innovation alone is honored.}
}

\section*{References and Verified Primary Sources}
\small
\begin{enumerate}
    \setlength{\itemsep}{3pt}
    \item \textbf{DESI Collaboration (M.\ Abdul Karim et al.\ 2025a)}, \emph{DESI DR2 Results II: Measurements of Baryon Acoustic Oscillations and Cosmological Constraints}, \href{https://arxiv.org/abs/2503.14738}{arXiv:2503.14738 [astro-ph.CO]} (\href{https://arxiv.org/html/2503.14738v1}{Full Text HTML}).
    \item \textbf{DESI Collaboration (K.\ Lodha et al.\ 2025b)}, \emph{Extended Dark Energy analysis using DESI DR2 BAO measurements}, Phys.\ Rev.\ D (2025), \href{https://arxiv.org/abs/2503.14743}{arXiv:2503.14743 [astro-ph.CO]} (\href{https://arxiv.org/html/2503.14743v2}{Full Text HTML}; \href{https://link.aps.org/doi/10.1103/w4c6-1r5j}{DOI:10.1103/w4c6-1r5j}).
    \item \textbf{Y.\ Carloni, O.\ Luongo, \& M.\ Biesiada (2026)}, \emph{Do we really need alternatives to the $w_0 w_a$CDM parameterization after the DESI DR2?}, \href{https://arxiv.org/pdf/2608.01215}{arXiv:2608.01215 [astro-ph.CO]}.
    \item \textbf{W.\ L.\ Freedman, B.\ F.\ Madore, I.\ S.\ Jang, T.\ J.\ Hoyt, A.\ J.\ Lee, \& K.\ A.\ Owens (2025)}, \emph{Status Report on the Chicago-Carnegie Hubble Program (CCHP): Measurement of the Hubble Constant Using the Hubble and James Webb Space Telescopes}, \href{https://arxiv.org/abs/2408.06153}{arXiv:2408.06153v3 [astro-ph.CO]}.
    \item \textbf{A.\ J.\ Lee, W.\ L.\ Freedman, B.\ F.\ Madore, I.\ S.\ Jang, K.\ A.\ Owens, \& T.\ J.\ Hoyt (2025)}, \emph{The Chicago-Carnegie Hubble Program: The JWST J-region Asymptotic Giant Branch (JAGB) Extragalactic Distance Scale}, Astrophys.\ J.\ (2025), \href{https://arxiv.org/abs/2408.03474}{arXiv:2408.03474v3 [astro-ph.GA]}.
    \item \textbf{A.\ G.\ Riess, D.\ Scolnic, G.\ S.\ Anand, L.\ Breuval, S.\ Casertano, L.\ M.\ Macri, S.\ Li, W.\ Yuan, et al.\ (2024)}, \emph{JWST Validates HST Distance Measurements: Selection of Supernova Subsample Explains Differences in JWST Estimates of Local $H_0$}, Astrophys.\ J.\ 977, 120 (2024), \href{https://arxiv.org/abs/2408.11770}{arXiv:2408.11770 [astro-ph.CO]} (\href{https://arxiv.org/html/2408.11770v2}{Full Text HTML}).
    \item \textbf{A.\ G.\ Riess et al.\ (SH0ES Collaboration 2022)}, \emph{A Comprehensive Measurement of the Local Value of the Hubble Constant with $1\text{ km s}^{-1}\text{ Mpc}^{-1}$ Uncertainty from the Hubble Space Telescope and the SH0ES Team}, Astrophys.\ J.\ Lett.\ 934, L7 (2022), \href{https://arxiv.org/pdf/2112.04510}{arXiv:2112.04510 [astro-ph.CO]}.
    \item \textbf{D.\ Brout et al.\ (Pantheon+ Collaboration 2022)}, \emph{The Pantheon+ Analysis: Cosmological Constraints}, Astrophys.\ J.\ 938, 110 (2022), \href{https://arxiv.org/abs/2202.04077}{arXiv:2202.04077 [astro-ph.CO]}.
    \item \textbf{D.\ Scolnic et al.\ (Pantheon+ Collaboration 2022)}, \emph{The Pantheon+ Analysis: The Full Dataset and Light-Curve Release}, Astrophys.\ J.\ 938, 113 (2022), \href{https://arxiv.org/abs/2112.03863}{arXiv:2112.03863 [astro-ph.CO]}.
    \item \textbf{Planck Collaboration (N.\ Aghanim et al.\ 2020)}, \emph{Planck 2018 results. VI. Cosmological parameters}, Astron.\ Astrophys.\ 641, A6 (2020), \href{https://arxiv.org/abs/1807.06209}{arXiv:1807.06209 [astro-ph.CO]}.
    \item \textbf{E.\ Silva, M.\ A.\ Sabogal, M.\ Scherer, R.\ C.\ Nunes, E.\ Di Valentino, \& S.\ Kumar (2025)}, \emph{New constraints on interacting dark energy from DESI DR2 BAO observations}, \href{https://arxiv.org/pdf/2503.23225}{arXiv:2503.23225 [astro-ph.CO]}.
    \item \textbf{Y.\ Wang \& K.\ Freese (2025)}, \emph{Model-Independent Dark Energy Measurements from DESI DR2}, \href{https://arxiv.org/html/2505.17415v4}{arXiv:2505.17415v4 [astro-ph.CO]}.
    \item \textbf{DESI Collaboration (A.\ G.\ Adame et al.\ 2024)}, \emph{DESI 2024 VI: Cosmological Constraints from the Measurements of Baryon Acoustic Oscillations}, \href{https://arxiv.org/abs/2404.03002}{arXiv:2404.03002 [astro-ph.CO]}.
    \item \textbf{Universal Distance Modes Collaboration (2026)}, \emph{Universal distance modes from DESI BAO and Type Ia supernovae}, \href{https://arxiv.org/pdf/2606.30527}{arXiv:2606.30527 [astro-ph.CO]}.
    \item \textbf{Sound-Horizon-Free BAO Collaboration (2026)}, \emph{A sound-horizon-free measurement of the Hubble constant from DESI DR2}, \href{https://arxiv.org/pdf/2604.24050}{arXiv:2604.24050 [astro-ph.CO]}.
    \item \textbf{DESI Collaboration (W.\ Elbers et al.\ 2025c)}, \emph{Constraints on Neutrino Physics from DESI DR2 BAO and DR1 Full Shape}, Phys.\ Rev.\ D 112, 083513 (2025), \href{https://arxiv.org/abs/2503.14744}{arXiv:2503.14744 [astro-ph.CO]} (\href{https://link.aps.org/doi/10.1103/w9pk-xsk7}{DOI:10.1103/w9pk-xsk7}).
    \item \textbf{TDCOSMO Collaboration (2026)}, \emph{TDCOSMO XXVII: JWST-based Lens Models and $H_0$ Constraints from Strong Gravitational Lensing Time Delays}, \href{https://www.alphaxiv.org/ja/abs/2608.27566}{arXiv:2608.27566 [astro-ph.CO]}.
    \item \textbf{Critical Review Collaboration (2025)}, \emph{Is Dark Energy Dynamical in the DESI Era? A Critical Review of BAO, CMB, and Low-Redshift Supernova Systematics}, \href{https://arxiv.org/html/2512.10585v2}{arXiv:2512.10585v2 [astro-ph.CO]}.
    \item \textbf{Sound-Horizon-Free $H_0$ Collaboration (2026)}, \emph{Insights from a Diverse Set of Sound-Horizon-Free $H_0$ Measurements}, \href{https://arxiv.org/pdf/2601.00650}{arXiv:2601.00650 [astro-ph.CO]}.
    \item \textbf{DESI DR2 Ly$\alpha$ Forest Dark Energy Analysis (2025)}, \emph{Evidence for Dynamical Dark Energy Using DESI DR2 Ly$\alpha$ Forest BAO}, \href{https://arxiv.org/html/2510.21976v1}{arXiv:2510.21976v1 [astro-ph.CO]}.
\end{enumerate}


\appendix

\section{Complete 32 Cosmic Chronometers Dataset and Analytical Likelihood Proofs}

\subsection{Full 32 Model-Independent Cosmic Chronometer $H(z)$ Measurements}

Table~\ref{tab:cc_full_32} presents the complete compilation of 32 differential-age Cosmic Chronometer measurements $H(z) = -\frac{1}{1+z}\frac{dz}{dt}$ spanning $0.070 \le z \le 1.965$ used in our joint likelihood (\href{https://arxiv.org/pdf/2608.01215}{Carloni, Luongo, \& Biesiada 2026, arXiv:2608.01215}), derived from spectroscopically dated passively evolving early-type galaxies without any reliance on the early-Universe sound horizon $r_d$.

\begin{table}[!htbp]
\centering
\caption{\textbf{Complete Compilation of 32 Cosmic Chronometer $H(z)$ Measurements ($0.070 \le z \le 1.965$).} Values and $1\sigma$ uncertainties $\sigma_H$ are in units of $\text{km s}^{-1}\text{ Mpc}^{-1}$ (\href{https://arxiv.org/pdf/2608.01215}{Carloni et al.\ 2026}).}
\label{tab:cc_full_32}
\vspace{3pt}
\footnotesize
\setlength{\tabcolsep}{5.5pt}
\renewcommand{\arraystretch}{1.12}
\begin{tabular}{cccc|cccc|cccc}
\toprule\rowcolor{tableheadbg}
\textbf{Bin} & $z$ & $H(z)$ & $\sigma_H$ & \textbf{Bin} & $z$ & $H(z)$ & $\sigma_H$ & \textbf{Bin} & $z$ & $H(z)$ & $\sigma_H$ \\
\midrule
1  & $0.0700$ & $69.0$ & $19.6$ & 12 & $0.4000$ & $95.0$  & $17.0$ & 23 & $0.8750$ & $125.0$ & $17.0$ \\
2  & $0.0900$ & $69.0$ & $12.0$ & 13 & $0.4004$ & $77.0$  & $10.2$ & 24 & $0.8800$ & $90.0$  & $40.0$ \\
3  & $0.1200$ & $68.6$ & $26.2$ & 14 & $0.4250$ & $87.1$  & $11.2$ & 25 & $0.9000$ & $117.0$ & $23.0$ \\
4  & $0.1700$ & $83.0$ & $8.0$  & 15 & $0.4450$ & $92.8$  & $12.9$ & 26 & $1.0370$ & $154.0$ & $20.0$ \\
5  & $0.1790$ & $75.0$ & $4.0$  & 16 & $0.4700$ & $89.0$  & $23.0$ & 27 & $1.3000$ & $168.0$ & $17.0$ \\
6  & $0.1990$ & $75.0$ & $5.0$  & 17 & $0.4783$ & $80.9$  & $9.0$  & 28 & $1.3630$ & $160.0$ & $33.6$ \\
7  & $0.2000$ & $72.9$ & $29.6$ & 18 & $0.4800$ & $97.0$  & $62.0$ & 29 & $1.4300$ & $177.0$ & $18.0$ \\
8  & $0.2700$ & $77.0$ & $14.0$ & 19 & $0.5930$ & $104.0$ & $13.0$ & 30 & $1.5300$ & $140.0$ & $14.0$ \\
9  & $0.2800$ & $88.8$ & $36.6$ & 20 & $0.6800$ & $92.0$  & $8.0$  & 31 & $1.7500$ & $202.0$ & $40.0$ \\
10 & $0.3520$ & $83.0$ & $14.0$ & 21 & $0.7500$ & $98.8$  & $33.6$ & 32 & $1.9650$ & $186.5$ & $50.4$ \\
11 & $0.3800$ & $83.0$ & $13.5$ & 22 & $0.7810$ & $105.0$ & $12.0$ & --- & --- & --- & --- \\
\bottomrule
\end{tabular}
\end{table}

\subsection{Exact Derivation of the Analytically Marginalized Supernova Likelihood and Pivot Transform}

For an uncalibrated Type Ia supernova sample with apparent magnitude vector $\mathbf{m}_{B,\text{corr}}$, theoretical dimensionless distance modulus $\tilde{\boldsymbol{\mu}}(\Omega_m, w_0, w_a) = 5\log_{10}[(1+\mathbf{z})\, I(\mathbf{z})]$, and combined statistical plus systematic covariance matrix $\mathbf{C}_{\rm SN}$, the full Gaussian likelihood over the nuisance intercept $\mathcal{M} = M_B^0 - 5\log_{10}(H_0/c) + 25$ is:
\begin{equation}
    \mathcal{L}_{\rm SN}(\Omega_m, w_0, w_a) = \int_{-\infty}^{+\infty} \frac{d\mathcal{M}}{\sqrt{(2\pi)^N \det\mathbf{C}_{\rm SN}}} \exp\!\left[ -\frac{1}{2}\left(\Delta\tilde{\boldsymbol{\mu}} - \mathcal{M}\mathbf{1}\right)^T \mathbf{C}_{\rm SN}^{-1} \left(\Delta\tilde{\boldsymbol{\mu}} - \mathcal{M}\mathbf{1}\right) \right],
\end{equation}
where $\Delta\tilde{\boldsymbol{\mu}} \equiv \mathbf{m}_{B,\text{corr}} - \tilde{\boldsymbol{\mu}}(\Omega_m, w_0, w_a)$. Expanding the quadratic form into $A - 2\mathcal{M}B + \mathcal{M}^2 C$ with $A = \Delta\tilde{\boldsymbol{\mu}}^T \mathbf{C}_{\rm SN}^{-1} \Delta\tilde{\boldsymbol{\mu}}$, $B = \Delta\tilde{\boldsymbol{\mu}}^T \mathbf{C}_{\rm SN}^{-1} \mathbf{1}$, and $C = \mathbf{1}^T \mathbf{C}_{\rm SN}^{-1} \mathbf{1}$, and completing the square in $\mathcal{M}$ yields the exact marginalized chi-square $\chi^2_{\text{SN,marg}} = A - B^2/C + \ln(C/2\pi)$, with maximum-likelihood intercept $\hat{\mathcal{M}} = B/C$ and variance $\sigma_{\mathcal{M}}^2 = 1/C$.

For the pivoted CPL transformation $(w_0, w_a) \mapsto (w_p, f_p)$ at scale factor $a_p = (1+z_p)^{-1}$, let $\mathbf{C}_{w_0, w_a}$ be the $2\times 2$ MCMC posterior covariance matrix of $(w_0, w_a)$. Since $w(a) = w_0 + (1-a)w_a$, the variance of $w(a)$ is $\sigma^2_{w(a)} = \sigma_{w_0}^2 + 2(1-a)\,\text{Cov}(w_0, w_a) + (1-a)^2 \sigma_{w_a}^2$. Minimizing $\sigma^2_{w(a)}$ with respect to $a$ sets $\frac{d\sigma^2_{w(a)}}{da} = 0$, which gives the unique decorrelation scale factor:
\begin{equation}
    a_p = 1 + \frac{\text{Cov}(w_0, w_a)}{\sigma_{w_a}^2} \simeq 0.76 \quad (z_p \simeq 0.31), \qquad \sigma_{w_p}^2 = \sigma_{w_0}^2\left(1 - \rho_{w_0, w_a}^2\right),
\end{equation}
where $\rho_{w_0, w_a} \approx -0.89$ reduces the uncertainty at the pivot redshift by a factor of $\sqrt{1 - \rho^2} \approx 0.456$, giving $\sigma_{w_p} \approx 0.024$ compared to $\sigma_{w_0} \approx 0.055$.

\section{Complete Reproducible Python Inference \& Simulation Pipeline}

Listing~\ref{lst:python_pipeline} provides the self-contained Python implementation used to evaluate the DESI DR2 BAO, Planck CMB, Cosmic Chronometer, and Pantheon+ likelihoods, execute the Affine-Invariant Ensemble MCMC sampler (Goodman \& Weare 2010), perform non-parametric Gaussian Process reconstruction of $w(z)$, $f_{\rm DE}(z)$, $H(z)/(1+z)$, $q(z)$, and $Om(z)$, and compute the fourth-order cosmographic expansion and Bayesian evidence metrics.

\begin{lstlisting}[language=Python, caption={Self-contained Python pipeline for DESI DR2 BAO + Pantheon+ SNe Ia + Cosmic Chronometers + Planck CMB + JWST $H_0$ Bayesian MCMC, Gaussian Process reconstruction, and model comparison.}, label={lst:python_pipeline}]
import numpy as np
import scipy.linalg as la
import scipy.optimize as opt
import scipy.stats as stats

c_light = 299792.458  # Speed of light [km/s]
UG_24, WG_24 = np.polynomial.legendre.leggauss(24)
UG_48, WG_48 = np.polynomial.legendre.leggauss(48)

# 1. DESI DR2 BAO Dataset (7 redshift bins, 13 observables; DESI Collaboration 2025a, Table 4)
desi_dr2_bins = [
    {"tracer": "BGS",       "z": 0.295, "type": "DV",    "DV_rd": 7.944,  "err_DV": 0.075},
    {"tracer": "LRG1",      "z": 0.510, "type": "DM_DH", "DM_rd": 13.587, "err_DM": 0.169, "DH_rd": 21.863, "err_DH": 0.427, "r_MH": -0.475},
    {"tracer": "LRG2",      "z": 0.706, "type": "DM_DH", "DM_rd": 17.347, "err_DM": 0.180, "DH_rd": 19.458, "err_DH": 0.332, "r_MH": -0.423},
    {"tracer": "LRG3+ELG1", "z": 0.934, "type": "DM_DH", "DM_rd": 21.574, "err_DM": 0.153, "DH_rd": 17.641, "err_DH": 0.193, "r_MH": -0.425},
    {"tracer": "ELG2",      "z": 1.321, "type": "DM_DH", "DM_rd": 27.605, "err_DM": 0.320, "DH_rd": 14.178, "err_DH": 0.217, "r_MH": -0.437},
    {"tracer": "QSO",       "z": 1.484, "type": "DM_DH", "DM_rd": 30.519, "err_DM": 0.758, "DH_rd": 12.816, "err_DH": 0.513, "r_MH": -0.489},
    {"tracer": "Lya QSO",   "z": 2.330, "type": "DM_DH", "DM_rd": 38.988, "err_DM": 0.531, "DH_rd": 8.632,  "err_DH": 0.101, "r_MH": -0.431},
]

z_bao_all = np.array([b["z"] for b in desi_dr2_bins])
Y_bao_obs, C_bao, idx = [], np.zeros((13, 13)), 0
for b in desi_dr2_bins:
    if b["type"] == "DV":
        Y_bao_obs.append(b["DV_rd"])
        C_bao[idx, idx] = b["err_DV"]**2
        idx += 1
    else:
        Y_bao_obs.extend([b["DM_rd"], b["DH_rd"]])
        s1, s2, r = b["err_DM"], b["err_DH"], b["r_MH"]
        C_bao[idx:idx+2, idx:idx+2] = [[s1**2, r*s1*s2], [r*s1*s2, s2**2]]
        idx += 2
Y_bao_obs = np.array(Y_bao_obs)
C_bao_inv = la.inv(C_bao)

# 2. Exact Expansion Rate E(z) and Gauss-Legendre Comoving Distance Integrator
def E_z(z, Om, w0=-1.0, wa=0.0, Or=9.0e-5):
    zp1 = 1.0 + z
    Ode = 1.0 - Om - Or
    f_de = (zp1**(3.0 * (1.0 + w0 + wa))) * np.exp(-3.0 * wa * z / zp1)
    return np.sqrt(np.maximum(Or * (zp1**4) + Om * (zp1**3) + Ode * f_de, 1e-12))

def comoving_dist_over_c_H0(z_targets, Om, w0=-1.0, wa=0.0, Or=9.0e-5):
    z_targets = np.atleast_1d(z_targets)
    zp = 0.5 * z_targets[:, None] * (UG_24[None, :] + 1.0)
    return 0.5 * z_targets * np.sum(WG_24[None, :] / E_z(zp, Om, w0, wa, Or), axis=1)

# 3. DESI DR2 BAO, Planck Compressed CMB, Cosmic Chronometers, and Pantheon+ Chi-Square
def chi2_desi_bao(H0, Om, w0=-1.0, wa=0.0, rd=147.09):
    I_z = comoving_dist_over_c_H0(z_bao_all, Om, w0, wa)
    Ez_all = E_z(z_bao_all, Om, w0, wa)
    pref = c_light / (H0 * rd)
    DM_rd, DH_rd = pref * I_z, pref / Ez_all
    Y_th = np.empty(13)
    Y_th[0] = (z_bao_all[0] * DH_rd[0] * (DM_rd[0]**2))**(1.0 / 3.0)
    Y_th[1::2], Y_th[2::2] = DM_rd[1:], DH_rd[1:]
    diff = Y_bao_obs - Y_th
    return float(diff @ C_bao_inv @ diff)

# 4. Affine-Invariant Ensemble MCMC Sampler (Goodman & Weare 2010 Stretch Move)
def run_ensemble_mcmc(log_prob_fn, init_center, init_scales, n_walkers=36, n_steps=600, burn_in=200):
    ndim = len(init_center)
    walkers = init_center + init_scales * np.random.randn(n_walkers, ndim)
    log_p = np.array([log_prob_fn(w) for w in walkers])
    chain = np.zeros((n_steps - burn_in, n_walkers, ndim))
    half = n_walkers // 2
    idx1, idx2 = np.arange(0, half), np.arange(half, n_walkers)
    for step in range(n_steps):
        for act, comp in [(idx1, idx2), (idx2, idx1)]:
            u = np.random.rand(len(act))
            z_s = ((2.0 - 1.0) * u + 1.0)**2 / 2.0
            part = np.random.choice(comp, size=len(act), replace=True)
            prop = walkers[part] + z_s[:, None] * (walkers[act] - walkers[part])
            log_p_prop = np.array([log_prob_fn(p) for p in prop])
            acc = np.log(np.random.rand(len(act))) < ((ndim - 1.0) * np.log(z_s) + log_p_prop - log_p[act])
            walkers[act[acc]], log_p[act[acc]] = prop[acc], log_p_prop[acc]
        if step >= burn_in:
            chain[step - burn_in] = walkers
    return chain.reshape(-1, ndim)

# 5. Non-Parametric Gaussian Process Reconstruction of w(z), f_DE(z), q(z), and Om(z)
def reconstruct_gp_dark_energy(z_grid, w_prior_mean, sigma_z, ell_corr=0.85, n_draws=1500):
    K_w = np.outer(sigma_z, sigma_z) * np.exp(-0.5 * ((z_grid[:, None] - z_grid[None, :]) / ell_corr)**2)
    L_w = la.cholesky(K_w + 1e-6 * np.eye(len(z_grid)), lower=True)
    w_draws = np.zeros((n_draws, len(z_grid)))
    f_de_draws = np.zeros((n_draws, len(z_grid)))
    qz_draws = np.zeros((n_draws, len(z_grid)))
    Om_diag_draws = np.zeros((n_draws, len(z_grid)))
    for i in range(n_draws):
        w_s = w_prior_mean + L_w @ np.random.randn(len(z_grid))
        w_draws[i] = w_s
        integrand = 3.0 * (1.0 + w_s) / (1.0 + z_grid)
        cum_int = np.zeros_like(z_grid)
        cum_int[1:] = np.cumsum(0.5 * (integrand[:-1] + integrand[1:]) * np.diff(z_grid))
        f_de = np.exp(cum_int)
        f_de_draws[i] = f_de
        Om_i = np.random.normal(0.3114, 0.0057)
        E2 = np.maximum(Om_i * (1.0 + z_grid)**3 + (1.0 - Om_i) * f_de, 1e-8)
        Ode_z = (1.0 - Om_i) * f_de / E2
        qz_draws[i] = 0.5 * (1.0 + 3.0 * w_s * Ode_z)
        Om_diag_draws[i] = (E2 - 1.0) / np.maximum((1.0 + z_grid)**3 - 1.0, 1e-6)
    return w_draws, f_de_draws, qz_draws, Om_diag_draws

# 6. Fourth-Order y-Redshift Cosmographic Expansion & Tension Metric Matrix
def dl_cosmography_y(z_arr, H0, q0, j0, s0):
    y = z_arr / (1.0 + z_arr)
    poly = (y + 0.5 * (1.0 - q0) * y**2
            - (1.0 / 6.0) * (1.0 - q0 - 3.0 * q0**2 + j0) * y**3
            + (1.0 / 24.0) * (2.0 - 2.0 * q0 - 15.0 * q0**2 - 15.0 * q0**3
                              + 5.0 * j0 + 10.0 * q0 * j0 + s0) * y**4)
    return (c_light / H0) * (1.0 / (1.0 - y)) * poly

def hubble_tension_sigma(H0_early, sig_early, H0_local, sig_local):
    return abs(H0_local - H0_early) / np.sqrt(sig_early**2 + sig_local**2)

# 7. Nested Sampling / Savage-Dickey Bayes Factor & KL Divergence Occam Penalty
def compute_bayes_factor_savage_dickey(samples_w0wa, prior_volume=13.5, lcdm_point=(-1.0, 0.0)):
    kde_2d = stats.gaussian_kde(samples_w0wa[:, 2:4].T)
    post_at_lcdm = float(kde_2d(lcdm_point)[0])
    prior_at_lcdm = 1.0 / prior_volume
    ln_B_w0wa_vs_lcdm = np.log(prior_at_lcdm / np.maximum(post_at_lcdm, 1e-15))
    mean_w = np.mean(samples_w0wa[:, 2:4], axis=0)
    cov_w = np.cov(samples_w0wa[:, 2:4], rowvar=False)
    a_pivot = 1.0 + cov_w[0, 1] / cov_w[1, 1]
    z_pivot = (1.0 / a_pivot) - 1.0
    sigma_wp = np.sqrt(cov_w[0, 0] * (1.0 - (cov_w[0, 1]**2) / (cov_w[0, 0] * cov_w[1, 1])))
    return ln_B_w0wa_vs_lcdm, a_pivot, z_pivot, sigma_wp
\end{lstlisting}

\vspace{6pt}
\noindent\textbf{Reproducibility \& Code Availability Statement:} The complete numerical simulation suite, MCMC chains, Gaussian Process reconstruction scripts, and publication figure generators authored by \textbf{Md.\ Mushfiqur Rahim} (\href{https://github.com/MD-Mushfiqur123}{\texttt{github.com/MD-Mushfiqur123}}) are self-contained in Listing~\ref{lst:python_pipeline} above and archived alongside the open-source research portfolio of the author.


\section{Glossary of Key Terms and Mathematical Notation Reference}

\subsection{Alphabetical Glossary of Cosmological and Statistical Terms}

\begin{table}[!htbp]
\centering
\caption{\textbf{Alphabetical Glossary of Key Observational, Relativistic, and Statistical Terms.}}
\label{tab:glossary_terms}
\footnotesize
\renewcommand{\arraystretch}{1.22}
\begin{tabularx}{\textwidth}{>{\bfseries\color{primary}}p{0.28\textwidth} X}
\toprule
\rowcolor{primary} \textcolor{white}{\textbf{Term}} & \textcolor{white}{\textbf{Definition and Role in This Monograph}} \\
\midrule
Alcock-Paczynski (AP) Effect & Geometric distortion $F_{\rm AP}(z) \equiv D_M(z)/D_H(z) = (H(z)/c)D_M(z)$ comparing transverse angular sizes to radial redshift intervals in 3D spectroscopic clustering. \\
Baryon Acoustic Oscillations (BAO) & Relic sound-horizon ruler $r_d \approx 147.09\text{ Mpc}$ imprinted in galaxy, quasar, and Ly$\alpha$ forest clustering by primordial baryon-photon acoustic waves prior to $z_d \approx 1060$. \\
Bayes Factor ($\ln B_{ij}$) & Natural log-ratio of marginal Bayesian evidences $\ln(\mathcal{Z}_i/\mathcal{Z}_j) = \Delta\langle\ln\mathcal{L}\rangle_{\mathcal{P}} - \mathcal{D}_{\rm KL}$, penalizing prior-volume complexity via the Kullback-Leibler divergence. \\
Cepheid Leavitt Law & Near-infrared Wesenheit Period-Luminosity-Color relation $W_{H,VI}$ of radially pulsating supergiant stars used by SH0ES to calibrate SN Ia host galaxy distances. \\
Cosmic Chronometers (CC) & Model-independent determinations of $H(z) = -(1+z)^{-1}dz/dt$ obtained by spectroscopic differential age dating ($\Delta t$) of passively evolving massive early-type galaxies. \\
Cosmography ($q_0, j_0, s_0$) & Purely kinematic Taylor expansion of $a(t)$ or $D_L(y)$ in $y = z/(1+z)$ defining deceleration $q_0$, jerk $j_0$, and snap $s_0$ without assuming dynamical field equations. \\
CPL ($w_0 w_a$CDM) Model & Chevallier-Polarski-Linder first-order scale-factor expansion $w(a) = w_0 + w_a(1-a) = w_0 + w_a z/(1+z)$ of a time-evolving dark energy equation of state. \\
Early Dark Energy (EDE) & Ultralight axion-like scalar field injecting a transient energy fraction $f_{\rm EDE}(z_c) \approx 0.092$ near $z_c \approx 3630$, shrinking $r_d$ to $\sim 139.7\text{ Mpc}$ and raising $H_0$ to $\sim 71.85\text{ km s}^{-1}\text{ Mpc}^{-1}$. \\
Gaussian Process (GP) & Non-parametric Bayesian regression over continuous function spaces defined by a prior mean $m(z)$ and covariance kernel $k(z,z')$ (Matérn-$\nu$ or Squared-Exponential). \\
Horndeski Gravity & Most general four-dimensional scalar-tensor theory of gravity with second-order Euler-Lagrange equations and luminal gravitational-wave speed $c_T = c$. \\
JAGB Carbon Stars & Thermally pulsing J-region Asymptotic Giant Branch carbon stars ($\text{C/O} > 1$) exhibiting a constant near-infrared luminosity peak $M_J \approx -6.20\text{ mag}$. \\
Mirage Dark Energy & One-parameter dark energy family $w_a = -3.66(1+w_0)$ that leaves the CMB comoving acoustic distance $D_M(z_*)$ invariant while altering late-time expansion at $z \lesssim 1$. \\
$Om(z)$ Diagnostic & First-derivative null test $Om(z) \equiv [E^2(z)-1]/[(1+z)^3-1]$, which equals $\Omega_m = \text{const}$ at all redshifts if and only if the Universe is flat $\Lambda$CDM. \\
Phantom Crossing ($z_c$) & Epoch where $w(z)$ crosses the phantom divide $w = -1$ (measured at $z_c = 0.42 \pm 0.09$), corresponding to a maximum in dark energy density $f_{\rm DE}(z_c) = 1.09 \pm 0.05$. \\
Pivoted CPL ($w_p, f_p$) & Orthogonal transformation evaluating $w_p \equiv w(a_p)$ and $f_p \equiv f_{\rm DE}(a_p)$ at the decorrelation scale factor $a_p \simeq 0.76$ ($z_p \simeq 0.31$) where $\text{Cov}(w_p, w_a) = 0$. \\
TRGB (Tip of Red Giant Branch) & Population II halo standard candle based on the core helium-flash bolometric luminosity discontinuity ($M_I \approx -4.05\text{ mag}$) of low-mass red giant stars. \\
\bottomrule
\end{tabularx}
\end{table}

\subsection{Mathematical and Astrophysical Notation Reference}

\begin{table}[!htbp]
\centering
\caption{\textbf{Consolidated Mathematical and Astrophysical Notation Reference Across All Parts.}}
\label{tab:notation_ref}
\footnotesize
\renewcommand{\arraystretch}{1.22}
\begin{tabularx}{\textwidth}{>{\bfseries\color{primary}}p{0.24\textwidth} X}
\toprule
\rowcolor{primary} \textcolor{white}{\textbf{Symbol}} & \textcolor{white}{\textbf{Physical Meaning, Definition, and Units}} \\
\midrule
$g_{\mu\nu}, a(t), z$ & FLRW spacetime metric tensor, dimensionless scale factor ($a_0 = 1$), and cosmological redshift $1+z = a^{-1}$. \\
$H(z), H_0, h, E(z)$ & Hubble expansion rate $\dot{a}/a$ $[\text{km s}^{-1}\text{ Mpc}^{-1}]$, present value $H_0 = 100h$, and normalized rate $E(z) \equiv H(z)/H_0$. \\
$\Omega_r, \Omega_m, \Omega_K, \Omega_{\rm DE}(z)$ & Dimensionless density parameters $\rho_i / \rho_{\rm crit}$ for radiation, matter (baryons + CDM + massive neutrinos), spatial curvature, and dark energy. \\
$\omega_b, \omega_{\rm bc}, \sum m_\nu$ & Physical baryon density $\Omega_b h^2$, physical baryon+CDM density $(\Omega_b+\Omega_c)h^2$, and neutrino mass sum $[\text{eV}]$. \\
$r_d \equiv r_s(z_d), r_s(z_*)$ & Comoving sound horizon $[\text{Mpc}]$ at the baryon drag epoch ($z_d \approx 1059.94$) and CMB recombination ($z_* \approx 1089.92$). \\
$D_M(z), D_H(z), D_V(z)$ & Transverse comoving distance, radial Hubble distance $c/H(z)$, and isotropic BAO dilation scale $[z D_M^2 D_H]^{1/3}$ $[\text{Mpc}]$. \\
$w(z), w_0, w_a, f_{\rm DE}(z)$ & Dark energy equation of state $P_{\rm DE}/(\rho_{\rm DE}c^2)$, present value $w(0)$, CPL derivative $-dw/da|_{a=1}$, and normalized density $\rho_{\rm DE}(z)/\rho_{\rm DE,0}$. \\
$a_p, z_p, w_p, f_p$ & Decorrelation pivot scale factor ($a_p \simeq 0.76, z_p \simeq 0.31$), pivot equation of state $w(a_p)$, and pivot density $f_{\rm DE}(a_p)$. \\
$q(z), q_0, j_0, s_0, y$ & Deceleration parameter $-\ddot{a}/(aH^2)$, present deceleration $q_0$, cosmic jerk $j_0$, snap $s_0$, and bounded redshift $y \equiv z/(1+z) \in [0,1)$. \\
$m_{B,\text{corr}}, M_B^0, \mu(z), \mathcal{M}$ & Tripp-standardized SN Ia apparent magnitude, fiducial absolute magnitude, distance modulus $5\log_{10}(D_L/\text{Mpc})+25$, and degenerate intercept. \\
$\mathbf{C}_{\rm BAO}, \mathbf{C}_{\rm SN}, \mathbf{C}_{\rm CMB}$ & Observational covariance matrices for the 13 DESI DR2 BAO modes, 1701 Pantheon+ SNe Ia, and Planck shift parameters $(\theta_*, \omega_b, \omega_{\rm bc})$. \\
$T_{ij}, \Delta\chi^2_{\rm MAP}, \mathcal{D}_{\rm KL}, \ln B_{ij}$ & Pairwise Gaussian $H_0$ tension $[\sigma]$, MAP chi-square difference, Kullback-Leibler Occam penalty, and natural log Bayes factor. \\
\bottomrule
\end{tabularx}
\end{table}


\label{LastContentPage}
\clearpage
\thispagestyle{empty}
\AddToShipoutPictureBG*{\includegraphics[width=\paperwidth,height=\paperheight]{/tmp/book_assets/cover_back.png}}
\mbox{}
\end{document}

